\documentclass[11pt,a4paper]{article}
\usepackage[utf8]{inputenc}
\usepackage[T1]{fontenc}
\usepackage[czech]{babel}
\usepackage[margin=2.2cm]{geometry}
\usepackage{amsmath,amssymb,amsthm}
\usepackage{parskip}
\usepackage{enumitem}
\setlist{nosep,leftmargin=1.4em}
\usepackage{booktabs}
\usepackage{array}
\usepackage{xcolor}
\definecolor{linkblue}{RGB}{20,60,150}
\definecolor{defbg}{RGB}{240,244,252}
\definecolor{defframe}{RGB}{100,130,190}
\definecolor{markcol}{RGB}{176,84,0}
\definecolor{markbg}{RGB}{253,246,236}
\usepackage[most]{tcolorbox}
\newtcolorbox{defbox}[1][]{breakable,colback=defbg,colframe=defframe,boxrule=0.6pt,arc=1.5mm,left=2mm,right=2mm,top=1mm,bottom=1mm,title={#1},fonttitle=\bfseries}
\newtcolorbox{poznbox}[1][]{breakable,colback=markbg,colframe=markcol,boxrule=0.6pt,arc=1.5mm,left=2mm,right=2mm,top=1mm,bottom=1mm,title={#1},fonttitle=\bfseries}
\newtheorem{veta}{Věta}[section]
\theoremstyle{definition}
\newtheorem{definice}{Definice}[section]
\usepackage[colorlinks=true,linkcolor=linkblue,urlcolor=linkblue]{hyperref}
\usepackage{algorithm}
\usepackage{algpseudocode}

\floatname{algorithm}{Algoritmus}
\renewcommand{\algorithmicrequire}{\textbf{Vstup:}}
\renewcommand{\algorithmicensure}{\textbf{Výstup:}}

% Značka pro látku, která není ve slidech NDBI023 / NAIL116.
\newcommand{\mimo}{\textcolor{markcol}{\textsf{\footnotesize[$\blacklozenge$~nad rámec slidů]}}}
\newcommand{\mimoq}[1]{\textcolor{markcol}{\textsf{\footnotesize[$\blacklozenge$~nad rámec slidů: #1]}}}

\title{\textbf{Dobývání znalostí}\\[2mm]
\large Učební text k~okruhu státní závěrečné zkoušky}
\author{SZZ Umělá inteligence, MFF UK}
\date{srpen 2026}

\begin{document}
\maketitle

\begin{center}
\begin{minipage}{0.92\textwidth}
\small
Text pokrývá okruh \emph{Dobývání znalostí} magisterské SZZ oboru Umělá inteligence.
Kapitoly odpovídají \textbf{jedna ku jedné} větám oficiálního znění okruhu (viz mapovací
tabulka níže). Vychází ze slidů přednášek \textbf{NDBI023 Dobývání znalostí}
a~\textbf{NAIL116 Sociální sítě a~jejich analýza} (I.~Mrázová, MFF UK) a~je doplněn
o~standardní látku z~učebnic Han, Kamber, Pei: \emph{Data Mining --- Concepts and
Techniques}, Tan, Steinbach, Kumar: \emph{Introduction to Data Mining}, Aggarwal:
\emph{Data Mining --- The Textbook} a~pro sítě Easley, Kleinberg: \emph{Networks, Crowds
and Markets}, Newman: \emph{Networks} a~Zafarani, Abbasi, Liu: \emph{Social Media Mining}.
Obsahuje definice, pseudokódy, odvození, číselné příklady a~srovnávací tabulky; každá
kapitola končí shrnutím pro ústní zkoušku.
\end{minipage}
\end{center}

\begin{center}
\begin{minipage}{0.92\textwidth}
\begin{poznbox}[Rozsah textu a~značení]
\small
Rozsah kapitol odpovídá váze látky: nejvíce prostoru dostávají metody dobývání znalostí
(kap.~\ref{sec:metody}) a~analýza sociálních sítí (kap.~\ref{sec:sna}) --- tedy to, čemu
se obě přednášky věnují nejpodrobněji a~co okruh jmenuje nejkonkrétněji.

Značkou \mimo{} je označena látka, která \textbf{není ve slidech} ani jednoho z~obou kurzů
(slidy jsou k~dispozici kompletní), ale je v~textu ponechána --- buď proto, že ji
\emph{jmenuje oficiální znění okruhu}, nebo proto, že bez ní nedává okolní výklad smysl.
Slouží jako nabídka ke~škrtnutí: co označeno není, je ve slidech doloženo. Následuje-li za
značkou dvojtečka, upřesňuje, čeho přesně se výhrada týká --- typicky že souhrnná tabulka
nebo rámec jsou moje, ale jednotlivé položky v~nich ze slidů pocházejí.

Slidy pokrývají: analýzu dat a~shlukování, asociační pravidla, rozhodovací stromy
a~ensembly, SVM s~vyhodnocováním klasifikátorů, bayesovské modely a~ELM (NDBI023);
základní pojmy, modelování vlivu, homofilii a~afiliaci, analýzu vazeb a~komplexní sítě
s~detekcí komunit a~predikcí vazeb (NAIL116). Nejnápadnější mezera je
\textbf{proces KDD a~metodologie} (kap.~\ref{sec:paradigmata}). Kapitola
o~\textbf{charakteristických a~diskriminačních pravidlech} (kap.~\ref{sec:pravidla}) je
postavena především na slidech kurzu \emph{Internet a~klasifikační metody} (Holeňa,
přednášky \texttt{ikm1}--\texttt{ikm6}, zejména \texttt{ikm5}); podle Han--Kamber v~ní
zůstává jen vymezení obou druhů pravidel s~mírami t-weight/d-weight. Slidy
\texttt{nmr1}--\texttt{nmr5} z~téže složky patří jinému kurzu a~látku kapitoly
nepokrývají (viz poznámka na jejím začátku).
\end{poznbox}
\end{minipage}
\end{center}

\vspace{2mm}

\section*{Mapování okruhu na kapitoly}
\addcontentsline{toc}{section}{Mapování okruhu na kapitoly}

\begin{center}
\small
\begin{tabular}{@{}>{\raggedright\arraybackslash}p{8.6cm}>{\raggedright\arraybackslash}p{7.4cm}@{}}
\toprule
\textbf{Věta oficiálního znění okruhu} & \textbf{Kapitola v~tomto textu} \\
\midrule
Základní paradigmata dobývání znalostí.
  & \ref{sec:paradigmata}~Základní paradigmata dobývání znalostí \\
Příprava dat; výběr atributů a~metody pro analýzu jejich relevance.
  & \ref{sec:priprava}~Příprava dat, výběr atributů a~analýza jejich relevance \\
Metody pro dobývání znalostí; asociační pravidla, přístupy založené na principu učení
s~učitelem a~klastrová analýza.
  & \ref{sec:metody}~Metody pro dobývání znalostí (\ref{sec:metody}.1~asociační pravidla,
    \ref{sec:metody}.2--\ref{sec:metody}.3~učení s~učitelem, \ref{sec:metody}.4~klastrová
    analýza) \\
Metody pro extrakci charakteristických diskriminačních pravidel a~měření jejich
zajímavosti.
  & \ref{sec:pravidla}~Charakteristická a~diskriminační pravidla a~jejich zajímavost \\
Reprezentace, vyhodnocování a~vizualizace získaných znalostí.
  & \ref{sec:reprez}~Reprezentace, vyhodnocování a~vizualizace znalostí \\
Modely pro analýzu sociálních sítí; míry centrality, detekce komunit.
  & \ref{sec:sna}~Modely pro analýzu sociálních sítí, centralita a~detekce komunit \\
Praktické využití technik pro dobývání znalostí a~analýzu sociálních sítí.
  & \ref{sec:praxe}~Praktické využití technik \\
\bottomrule
\end{tabular}
\end{center}

\vspace{1mm}
\noindent\small\textit{Poznámka k~terminologii.} Text používá české termíny; anglický
ekvivalent je uveden kurzívou v~závorce při prvním výskytu, protože obě přednášky i~celá
literatura jsou anglicky.\normalsize

\vspace{2mm}
\tableofcontents
\newpage

%=====================================================================
\section{Základní paradigmata dobývání znalostí}
\label{sec:paradigmata}
%=====================================================================

Kapitola vymezuje obor, jeho proces a~základní dělení úloh a~dat --- rámec, do něhož se zasazují
všechny ostatní kapitoly.

\subsection{Motivace a~vymezení oboru}

Organizace shromažďují data v~objemech, které daleko přesahují kapacitu člověka je
prohlédnout. Klasické dotazování (SQL, OLAP) umí odpovědět na otázku, kterou už umíme
položit; \emph{dobývání znalostí} odpovídá na otázky, které jsme položit neuměli --- hledá
v~datech struktury, o~jejichž existenci jsme předem nevěděli.

Ilustrativní příklad: systém prošel 20 milionů výkonů účtovaných zdravotní pojišťovně,
214 tisíc z~nich označil za podivné a~14 tisíc nejpodezřelejších předal k~ruční kontrole.
Revizní lékaři, kteří jsou schopni prohlédnout řádově stovky případů, z~předložených
případů okamžitě označili 8 tisíc za podvodné. Stroj tu nenahradil experta ---
\emph{zaostřil} jeho pozornost.

\begin{defbox}[Dobývání znalostí z~databází]
\emph{Dobývání znalostí z~databází} (\emph{Knowledge Discovery in Databases}, KDD) je
netriviální proces získávání implicitních, dosud neznámých a~potenciálně užitečných
informací z~dat.

Klíčová jsou všechna čtyři slova:
\begin{itemize}
\item \textbf{netriviální} --- nejde o~prostý dotaz ani agregaci;
\item \textbf{implicitní} --- znalost v~datech není explicitně zapsána;
\item \textbf{dosud neznámá} --- výsledek, který každý zná, je bezcenný;
\item \textbf{potenciálně užitečná} --- musí vést k~akci nebo rozhodnutí.
\end{itemize}
\end{defbox}

Termín \emph{dobývání dat} (\emph{data mining}, DM) se v~úzkém pojetí vztahuje pouze
k~jednomu kroku procesu KDD (vlastní aplikaci analytických algoritmů), v~běžné praxi se
však oba pojmy používají zaměnitelně. Obor se konstituoval v~90.~letech 20.~století na
průsečíku tří disciplín:

\begin{itemize}
\item \textbf{umělá inteligence} --- zejména metody strojového učení (indukce pravidel,
      stromy, neuronové sítě);
\item \textbf{databázové systémy} --- ukládání a~efektivní přístup k~velkým datovým
      souborům, datové sklady, vyhledávání informací;
\item \textbf{statistika} --- modelování a~analýza závislostí nalezených v~datech,
      testování hypotéz.
\end{itemize}

Čtvrtou --- a~v~praxi rozhodující --- složkou je nutnost využít zpracovaná data k~podpoře
(strategického) rozhodování v~podniku. Bez ní se dobývání znalostí redukuje na akademické
cvičení.

\textbf{Vztah k~datovým skladům a~OLAP.} \mimo{} Typickým zdrojem dat pro KDD je
\emph{datový sklad} (\emph{data warehouse}) --- integrovaná, subjektově orientovaná, časově
variantní a~stálá databáze konsolidovaná z~provozních systémů. Data se v~něm organizují do
\emph{datové kostky}: \emph{dimenze} (čas, produkt, region, \dots) $\times$ \emph{míry}
(tržby, počty kusů); v~relační implementaci jde o~\emph{hvězdicové schéma} (tabulka faktů
obklopená tabulkami dimenzí), fyzicky MOLAP (multidimenzionální uložení) nebo ROLAP
(relační). \textbf{OLAP} (\emph{On-Line Analytical Processing}) nad kostkou umožňuje
agregační operace \emph{roll-up} / \emph{drill-down} a~řezy \emph{slice \& dice} ---
odpovídá ale jen na otázky, které si analytik umí \emph{předem položit}. Dobývání znalostí
jde o~krok dál: hledá vzory, o~nichž předem nevíme, že v~datech jsou. Souvislost
generalizace v~OLAP s~indukcí orientovanou na atributy viz kap.~\ref{sec:pravidla}.

\subsection{Proces KDD}

\mimoq{pětifázový rámec jako celek; náplň jednotlivých fází ve slidech je}
Dobývání znalostí je \textbf{interaktivní a~iterativní} proces: výsledky jednoho
kroku běžně vedou k~návratu ke krokům předchozím. Klasické schéma (Fayyad et al.) má pět
fází:

\begin{center}
\small
\begin{tabular}{@{}p{3.3cm}p{4.2cm}p{7.6cm}@{}}
\toprule
\textbf{Fáze} & \textbf{Vstup $\to$ výstup} & \textbf{Náplň} \\
\midrule
Výběr (\emph{selection}) & původní data $\to$ vybraná data & výběr relevantních záznamů
  a~atributů z~datového skladu \\
Předzpracování (\emph{preprocessing}) & vybraná data $\to$ předzpracovaná data & čištění,
  chybějící hodnoty, šum, odlehlé hodnoty, integrace zdrojů \\
Transformace (\emph{transformation}) & předzpracovaná data $\to$ transformovaná data &
  normalizace, diskretizace, odvozené atributy, redukce dimenze \\
Vlastní dobývání (\emph{data mining}) & transformovaná data $\to$ vzory & aplikace
  algoritmů (stromy, asociace, klastry, \dots) \\
Interpretace (\emph{interpretation}) & vzory $\to$ znalost & vyhodnocení z~pohledu
  koncového uživatele, vizualizace, nasazení \\
\bottomrule
\end{tabular}
\end{center}

Cílem prvních tří fází (souhrnně \emph{příprava dat}, kap.~\ref{sec:priprava}) je z~dat
uložených ve složité struktuře vybudovat \textbf{jednu tabulku} obsahující relevantní údaje
o~zkoumaných objektech (klientech banky, zákaznících, \dots). Tento krok typicky spotřebuje
60--80~\% celkového času projektu.

\textbf{Sedm kroků očima manažera.} \mimo{} Organizační pohled na tentýž proces (impuls $=$
problém, výstup $=$ znalost k~jeho řešení): \emph{(1)}~\emph{tým} --- experti na oblast, na data
a~na metody KDD, žádná role není zastupitelná; \emph{(2)}~\emph{specifikace problému} --- převod
obchodního zadání (\uv{zvýšit zisk} $\to$ \uv{predikovat odchod zákazníka});
\emph{(3)}~\emph{sběr dat} včetně externích o~prostředí (sezóna, kampaně, počasí) --- může vést
k~přeformulování problému; \emph{(4)}~\emph{volba metod} včetně vizualizačních;
\emph{(5)}~\emph{předzpracování}; \emph{(6)}~\emph{dolování} --- opakovaně, výsledky ovlivňují
parametry dalších běhů; \emph{(7)}~\emph{interpretace} --- analytická zpráva, případně akce.

\subsection{Typy úloh dobývání znalostí}

\mimo{} Podle vztahu nalezené znalosti k~cílovému konceptu:

\begin{center}
\small
\begin{tabular}{@{}p{2.9cm}p{6.0cm}p{6.4cm}@{}}
\toprule
\textbf{Typ úlohy} & \textbf{Cíl} & \textbf{Charakter výsledku} \\
\midrule
Klasifikace / predikce & zařadit nové případy, resp.\ odhadnout budoucí vývoj entity &
  preferujeme \emph{přesnost pokrytí} před jednoduchostí; hodně znalosti, obtížně
  interpretovatelné \\
Popis (charakterizace, profil) & najít dominantní strukturu či vztahy v~datech &
  požadujeme \emph{srozumitelnost} při plném pokrytí konceptu; méně znalosti a~méně přesné \\
Hledání \uv{nugetů} & najít zajímavou (novou, překvapivou) znalost & znalost \emph{nemusí}
  pokrývat celý koncept; hodnotí se překvapivost \\
\bottomrule
\end{tabular}
\end{center}

\emph{Predikce:} počasí, cena akcií, spotřeba elektřiny. \emph{Popis:} profil bonitního klienta.
\emph{Nuget:} neobvyklá kombinace zboží v~košíku, podezřelý vzorec transakcí.

\begin{defbox}[Deskriptivní vs.\ prediktivní úlohy]
\textbf{Deskriptivní} (\emph{descriptive}) úlohy charakterizují obecné vlastnosti dat:
charakterizace a~diskriminace, asociační a~korelační analýza, shlukování, detekce odlehlých
hodnot, analýza vývoje.

\textbf{Prediktivní} (\emph{predictive}) úlohy provádějí inferenci nad daty za účelem
predikce: klasifikace (diskrétní cílový atribut), regrese (spojitý cílový atribut),
predikce časových řad.

Ortogonální dělení podle dostupnosti cílové proměnné: \emph{učení s~učitelem}
(\emph{supervised}) --- klasifikace, regrese; \emph{učení bez učitele}
(\emph{unsupervised}) --- shlukování, asociační pravidla, detekce odlehlých hodnot;
\emph{poloučené} (\emph{semi-supervised}) --- malá označkovaná a~velká neoznačkovaná část.
\end{defbox}

Toto dělení je klíčem ke kap.~\ref{sec:metody}: okruh v~ní jmenuje po jednom representantu
učení bez učitele nad transakčními daty (asociační pravidla), učení s~učitelem a~učení bez
učitele nad metrickými daty (klastrová analýza).

\subsection{Typy dat}

\textbf{Podle typu atributu} --- určuje, jaké operace jsou smysluplné a~jaké metody lze
použít:

\begin{center}
\small
\begin{tabular}{@{}p{3.3cm}p{2.2cm}p{5.2cm}p{4.0cm}@{}}
\toprule
\textbf{Škála} & \textbf{Operace} & \textbf{Příklad} & \textbf{Statistika} \\
\midrule
Nominální (\emph{nominal}) & $=$, $\neq$ & barva, druh ovoce & modus, četnosti \\
Ordinální (\emph{ordinal}) & $=$, $<$ & známka, \{nízký, střední, vysoký\} & medián,
  kvantily \\
Intervalová (\emph{interval-scaled}) & $+$, $-$ & teplota ve $^{\circ}$C, datum & průměr,
  směr.\ odchylka \\
Poměrová (\emph{ratio-scaled}) & $\times$, $/$ & příjem, hmotnost, četnost & geometrický
  průměr \\
\bottomrule
\end{tabular}
\end{center}

U~intervalové škály \emph{mají intervaly v~celém rozsahu stejnou důležitost} (rozdíl věku 10
a~20 let je týž jako mezi 40 a~50). Poměrová má navíc absolutní nulu, takže má smysl poměr;
poměrové atributy s~exponenciálním chováním (počet mikroorganismů $Ae^{Bt}$) se logaritmují
a~dál se s~nimi zachází jako s~intervalovými. \emph{Nominální} atribut o~$v$ hodnotách se pro
numerické metody převede na $v$ binárních; \emph{ordinální} se zpracovává jako intervalový.

\textbf{Podle struktury datového souboru:}
\begin{itemize}
\item \textbf{Tabulková / relační data} --- objekty $\times$ atributy; standardní vstup
      většiny metod.
\item \textbf{Transakční data} --- proměnná délka záznamu (nákupní košík); vstup asociačních
      pravidel.
\item \textbf{Temporální data} (\emph{temporal}) --- např.\ časové řady kurzů akcií.
      Typickou úlohou je predikce budoucích hodnot; řada se transformuje do tabulky
      posuvným oknem: vstupy $y(t_0),y(t_1),y(t_2),y(t_3)$ $\to$ výstup $y(t_4)$;
      $y(t_1),\dots,y(t_4)\to y(t_5)$ atd.
\item \textbf{Prostorová data} (\emph{spatial}) --- geografické informační systémy;
      implicitní relace sousedství (hodnoty atributů sousedních objektů se od sebe příliš
      neliší), užitečné zejména pro interpretaci.
\item \textbf{Strukturní data} (\emph{structural}) --- např.\ chemické sloučeniny zapsané
      řetězci SMILES nebo soubory faktů v~Prologu; úlohou je třeba klasifikace látek podle
      struktury.
\item \textbf{Textová a~webová data, multimédia, grafy a~sítě} --- vyžadují specifickou
      reprezentaci (vektorový model, grafové modely, viz kap.~\ref{sec:sna}).
\end{itemize}

\subsection{Metodologie dobývání znalostí}
\label{ssec:metodologie}

\mimo{} Metodologie poskytuje jednotný rámec nezávislý na oboru podnikání a~umožňuje přenos
zkušeností z~úspěšných projektů. Vznikají u~výrobců BI nástrojů (SEMMA od SAS, 5A od SPSS),
nebo jako softwarově neutrální standardy (CRISP-DM).

\textbf{SEMMA} (SAS Enterprise Miner), pět fází: \emph{Sample} --- výběr dat (vzorkování,
imputace, rozdělení na trénovací/testovací/validační množinu); \emph{Explore} --- vizuální
průzkum a~redukce (odlehlé hodnoty, vícerozměrné histogramy, shlukování, Kohonenovy mapy);
\emph{Modify} --- příprava objektů, hodnot a~proměnných; \emph{Model} --- aplikace metod
(stromy, regrese, neuronové sítě); \emph{Assess} --- vyhodnocení spolehlivosti a~užitečnosti.

\textbf{CRISP-DM} (\emph{CRoss-Industry Standard Process for Data Mining}) --- otevřený,
oborově, nástrojově i~aplikačně neutrální model; financovala iniciativa ESPRIT, konsorcium ISL
(později SPSS, dnes IBM), Teradata, NCR, Daimler AG a~OHRA. Šest fází, jejichž \textbf{pořadí
není striktní} --- výsledky jedné fáze určují další kroky a~některé je nutné opakovat (proces se
kreslí jako cyklus s~vnějším cyklem nasazení a~monitoringu):

\begin{enumerate}
\item \textbf{Porozumění problému} (\emph{Business Understanding}) --- obchodní cíle a~kritéria
      úspěchu; inventura zdrojů (lidé, informace, software, kapacita), předpoklady a~omezení,
      rizika, náklady a~přínosy; cíle dobývání znalostí; plán projektu.
\item \textbf{Porozumění datům} (\emph{Data Understanding}) --- sběr, popis (množství,
      vlastnosti, formát), průzkum a~vizualizace (rozdělení atributů, vztahy dvojic,
      subpopulace, statistiky), ověření kvality.
\item \textbf{Příprava dat} (\emph{Data Preparation}) --- výběr, čištění, odvozené atributy,
      integrace, agregace, formátování; iterativně a~v~různém pořadí.
\item \textbf{Modelování} (\emph{Modeling}) --- volba techniky (její předpoklady mohou vynutit
      další úpravy dat), návrh testu, sestavení modelu s~různými parametry, posouzení.
\item \textbf{Vyhodnocení} (\emph{Evaluation}) --- výsledky vůči \emph{obchodním} kritériím
      úspěchu, schválení modelů, revize procesu, další kroky.
\item \textbf{Nasazení} (\emph{Deployment}) --- prezentace ve formě nasaditelné uživatelem;
      plán nasazení, monitoring a~údržba; závěrečná zpráva a~revize projektu (včetně slepých
      uliček).
\end{enumerate}

\textbf{ASUM-DM} (IBM, 2015) reaguje na hlavní výtku vůči CRISP-DM --- malou podporu
projektového řízení. Fáze \emph{Analyze}, \emph{Design}, \emph{Configure \& Build},
\emph{Deploy}, \emph{Operate \& Optimize}, průřezově \emph{Project Management}.

\begin{center}
\small
\begin{tabular}{@{}llll@{}}
\toprule
& \textbf{SEMMA} & \textbf{CRISP-DM} & \textbf{ASUM-DM} \\
\midrule
Původ & SAS Institute & ESPRIT / konsorcium & IBM (2015) \\
Neutralita & vázáno na SAS EM & nástrojově neutrální & vázáno na IBM \\
Byznys fáze & chybí & ano (fáze 1 a~5) & ano \\
Projektové řízení & ne & slabé & explicitní, průřezové \\
Počet fází & 5 & 6 & 5 + PM \\
\bottomrule
\end{tabular}
\end{center}

\subsection{Shrnutí ke~zkoušce}

\begin{itemize}
\item \textbf{KDD} $=$ netriviální proces získávání implicitních, dosud neznámých
      a~potenciálně užitečných informací z~dat; interaktivní a~iterativní.
\item \textbf{Pět fází KDD}: výběr $\to$ předzpracování $\to$ transformace $\to$ dobývání
      $\to$ interpretace. Příprava dat spotřebuje 60--80~\% času.
\item \textbf{Kořeny oboru}: umělá inteligence (strojové učení), databáze, statistika
      $+$ potřeba podpory rozhodování. \textbf{OLAP} odpovídá na předem položené otázky,
      KDD hledá vzory, o~nichž nevíme.
\item \textbf{Tři typy úloh}: klasifikace/predikce (přesnost $>$ jednoduchost), popis
      (srozumitelnost, plné pokrytí), hledání nugetů (překvapivost, nemusí pokrývat).
\item \textbf{Deskriptivní vs.\ prediktivní} úlohy; \textbf{s~učitelem vs.\ bez učitele}
      vs.\ \textbf{poloučené}.
\item \textbf{Škály atributů}: nominální, ordinální, intervalová, poměrová --- určují
      přípustné operace i~metody. Nominální $\to$ $v$ binárních atributů; poměrový
      s~exponenciálním chováním $\to$ logaritmovat.
\item \textbf{Metodologie}: SEMMA (Sample, Explore, Modify, Model, Assess), CRISP-DM
      (6 fází, pořadí není striktní, cyklické), ASUM-DM (CRISP-DM $+$ projektové řízení).
\item Nejčastější doplňující otázka: \emph{proč pořadí fází CRISP-DM není striktní?} ---
      protože výsledky fáze určují další kroky; např.\ zjištěná kvalita dat vede zpět
      k~přeformulování cíle.
\end{itemize}

%=====================================================================
\newpage
\section{Příprava dat, výběr atributů a~analýza jejich relevance}
\label{sec:priprava}
%=====================================================================

Okruh spojuje \emph{přípravu dat} (aby byla data vůbec zpracovatelná) a~\emph{výběr atributů
s~analýzou relevance} (vybrat ty, které nesou informaci). Druhé je speciálním případem
prvního (redukce dimenze), okruh je ale jmenuje zvlášť --- čeká se znalost konkrétních
kritérií a~umění je spočítat.

\subsection{Proč je příprava dat klíčová}

Předzpracování je pro úspěch aplikace \emph{klíčové}, je obtížné a~časově náročné a~těží ze
spolupráce s~doménovým expertem. Cíl: vybrat (nebo vytvořit) z~dostupných dat ta relevantní
pro danou úlohu a~reprezentovat je ve formě vhodné pro zvolený algoritmus; výsledkem je
\emph{jedna} datová tabulka objekty~$\times$ atributy.

\mimo{} Rozměry kvality dat: \emph{přesnost} (\emph{accuracy}), \emph{úplnost}
(\emph{completeness}), \emph{konzistence}, \emph{aktuálnost} (\emph{timeliness}),
\emph{důvěryhodnost}, \emph{interpretovatelnost} --- \uv{garbage in, garbage out}.

\subsection{Čištění dat: chybějící hodnoty}

Základní strategie (v~pořadí od nejjednodušší k~nejpřesnější):

\begin{enumerate}
\item \textbf{Ignorovat objekty} s~chybějící hodnotou --- jednoduché, ale při mnoha atributech
      ztratíme většinu dat a~zavádí zkreslení, není-li chybění náhodné.
\item \textbf{Nová hodnota \uv{nevím}} (\emph{I don't know}) --- absence se stane
      plnohodnotnou kategorií; vhodné, nese-li sama informaci (nevyplněný příjem u~žádosti
      o~úvěr).
\item \textbf{Existující hodnotou atributu}: nejčastější (modus, u~numerických průměr či
      medián), úměrně poměru všech hodnot (náhodně podle empirického rozdělení), nebo
      libovolnou.
\item \textbf{Modelem} (regrese, rozhodovací strom, $k$-NN) z~ostatních atributů ---
      nejpřesnější, výpočetně nejdražší, riskuje umělou závislost mezi atributy.
\end{enumerate}

\mimo{} Terminologicky se rozlišuje MCAR (\emph{missing completely at random}), MAR
(\emph{missing at random} --- chybění závisí na pozorovaných atributech) a~MNAR
(\emph{missing not at random} --- závisí na nepozorované hodnotě samotné). Jen u~MCAR je
prosté vypuštění nezkreslené.

Některé metody umí chybějící hodnoty zpracovat samy: C4.5 je při konstrukci stromu ignoruje
(kritérium dělení počítá jen ze známých hodnot) a~při klasifikaci je odhaduje z~ostatních
vzorů; CART chybějící hodnoty do rozhodovacího kritéria nezapočítává; naivní Bayes umí
klasifikovat i~nekompletně popsané vzory (kap.~\ref{sec:metody}).

\subsection{Čištění dat: šum a~odlehlé hodnoty}

\emph{Šum} (\emph{noise}) je náhodná chyba nebo rozptyl v~měřené veličině. Techniky
vyhlazování: \textbf{binning} (hodnoty se rozdělí do košů a~nahradí průměrem, mediánem nebo
hranicí koše), \textbf{regrese} (hodnoty se nahradí predikcí regresní funkce),
\textbf{shlukování} (hodnoty mimo shluky se označí za odlehlé) a~\textbf{kombinace stroje
a~člověka} (podezřelé hodnoty se předloží ke kontrole).

\begin{defbox}[Rozpětí, kvartily, mezikvartilové rozpětí, odlehlá hodnota]
\emph{Rozpětí} (\emph{range}) množiny dat je rozdíl nejvyšší a~nejnižší hodnoty.

\emph{Kvartily} $Q_1,Q_2,Q_3$ dělí data na čtyři části: 25~\% pozorování má hodnotu menší
než dolní kvartil $Q_1$, 50~\% pod mediánem $Q_2$ a~75~\% pod horním kvartilem $Q_3$.

\emph{Mezikvartilové rozpětí} (\emph{interquartile range}): $\mathrm{IQR}=Q_3-Q_1$.

\emph{Odlehlá hodnota} (\emph{outlier}) je extrémní hodnota ležící mimo celkový vzorec dat.
Běžné pravidlo: odlehlá je každá hodnota
\[ x > Q_3 + 1{,}5\cdot \mathrm{IQR} \qquad\text{nebo}\qquad x < Q_1 - 1{,}5\cdot \mathrm{IQR}. \]
\end{defbox}

\textbf{Výpočet kvartilů} pro $n$ vzestupně seřazených hodnot (diskrétní data): pro $Q_1$
dělíme $n$ čtyřmi --- je-li $n/4$ celé číslo, bereme střed mezi odpovídajícím a~následujícím
členem; není-li celé, zaokrouhlíme nahoru a~vezmeme odpovídající člen. Pro $Q_3$ analogicky
s~$3n/4$. U~spojitých dat se místo zaokrouhlení interpoluje.

\begin{defbox}[Příklad --- krabicový graf a~odlehlé hodnoty]
Hladina glukózy v~krvi u~30 žen (mmol/l), stonkový diagram (klíč: $2\mid 1$ znamená 2,1):
\begin{center}\ttfamily
2 | 2 2 3 3 5 7 \hfill (6)\\
3 | 1 2 6 7 7 7 8 8 8 8 9 9 9 \hfill (13)\\
4 | 0 0 0 0 4 5 6 7 8 \hfill (9)\\
5 | 1 5 \hfill (2)
\end{center}
$n=30$, $30/4=7{,}5$ $\to$ zaokrouhlíme na 8, osmé číslo je 3,2, tedy $Q_1=3{,}2$. Medián
$Q_2=3{,}8$. Pro $Q_3$: $3\cdot 30/4=22{,}5\to 23$, 23.\ číslo je 4,0, tedy $Q_3=4{,}0$.

$\mathrm{IQR}=4{,}0-3{,}2=0{,}8$; dolní mez $=3{,}2-1{,}5\cdot 0{,}8=2{,}0$; horní mez
$=4{,}0+1{,}5\cdot 0{,}8=5{,}2$.

\textbf{Odlehlá hodnota je 5,5.} Nejvyšší hodnota, která odlehlá není, je 5,1 --- v~krabicovém
grafu (\emph{box plot}) k~ní vede vous, odlehlá hodnota se vyznačí křížkem.
\end{defbox}

Násobek 1,5 je konvence (s~násobkem 1 je kritérium přísnější). \emph{Příklad:} želví samci
$Q_1=400$, $Q_3=580$~g $\Rightarrow$ $\mathrm{IQR}=180$, meze 220 a~760 --- z~hodnot 400,
260, 550, 640~g není odlehlá žádná. Samice $Q_1=260$, $Q_3=340$ $\Rightarrow$ meze 180
a~420 --- odlehlé jsou 170~g a~440~g.

\subsection{Integrace dat}

\mimo{} Při \emph{integraci} (\emph{data integration}) spojujeme více zdrojů do jedné
tabulky. Hlavní problémy: \textbf{problém identity entity} (atribut \texttt{cust\_id}
v~jedné databázi odpovídá \texttt{customer\_number} v~druhé; nutná metadata a~párování
záznamů); \textbf{redundance} (atribut odvoditelný z~jiných --- roční příjem $=$
12$\times$měsíční; detekuje se korelační analýzou u~numerických a~$\chi^2$ testem
u~kategoriálních atributů, viz odd.~\ref{ssec:filtr}); \textbf{konflikty hodnot} (různé
jednotky, škály, kódování); \textbf{duplicity} (tentýž objekt vícekrát s~drobně odlišnými
hodnotami).

\subsection{Transformace a~standardizace atributů}

Standardizace dává všem atributům \emph{stejný vliv} na výpočet vzdáleností (nutí je do
společného rozsahu). Pro $x_i=(0{,}1;\,20)$ a~$x_j=(0{,}9;\,720)$ je
\[ \mathrm{dist}(x_i,x_j)=\sqrt{(0{,}9-0{,}1)^2+(720-20)^2}=700{,}000457, \]
tedy vzdálenost je zcela ovládnuta druhým atributem; první nemá prakticky žádný vliv.

\begin{defbox}[Základní standardizační transformace]
Nechť $f$ je atribut a~$x_{if}$ jeho hodnota u~$i$-tého objektu.

\textbf{Desítkové škálování} (\emph{decimal scaling}) do $[-1,1]$: dělíme hodnoty nejmenší
mocninou 10, která udrží všechny transformované hodnoty v~$[-1,1]$.

\textbf{Min-max normalizace} (\emph{range standardization}) do $[0,1]$:
\[ \mathrm{range}(x_{if})=\frac{x_{if}-\min(f)}{\max(f)-\min(f)}, \]
do $[-1,1]$ pak $\mathrm{range}_{[-1,1]}(x_{if})=2\cdot \mathrm{range}(x_{if})-1$.

\textbf{Standardizace podle průměrné absolutní odchylky} --- nulový průměr a~jednotková
střední absolutní odchylka:
\[ m_f=\tfrac1n\big(x_{1f}+\cdots+x_{nf}\big),\qquad
s_f=\tfrac1n\big(|x_{1f}-m_f|+\cdots+|x_{nf}-m_f|\big),\qquad
z(x_{if})=\frac{x_{if}-m_f}{s_f}. \]

\textbf{Standardizace podle směrodatné odchylky} (\emph{std-score}) --- nulový průměr
a~jednotková korigovaná výběrová směrodatná odchylka:
\[ \mathrm{cor}s_f=\sqrt{\frac{1}{n-1}\sum_{i=1}^n\big(x_{if}-m_f\big)^2},\qquad
\mathrm{std}(x_{if})=\frac{x_{if}-m_f}{\mathrm{cor}s_f}. \]
\end{defbox}

Volba mezi z-skóry: střední absolutní odchylka je \emph{robustnější} k~odlehlým hodnotám než
směrodatná (odchylky se neumocňují), takže outliery stlačí ostatní data méně.

\subsection{Diskretizace numerických atributů}

Řada metod (asociační pravidla, ID3, CHAID, bayesovská klasifikace s~diskrétními atributy)
vyžaduje kategoriální vstup. \emph{Diskretizace} rozdělí rozsah numerického atributu na
intervaly. Metody se liší ve čtyřech ohledech: \emph{strategií} (dělení intervalů shora,
nebo slučování zdola), \emph{počtem} intervalů, \emph{typem} intervalů a~\emph{kritériem}
kvality (minimální klasifikační chyba, entropie, $\chi^2$ test).

\begin{itemize}
\item \textbf{Ekvidistantní} (\emph{equi-width}) --- rozsah se rozdělí na intervaly stejné
      délky. Nejjednodušší, ale citlivé na odlehlé hodnoty a~nerovnoměrné rozdělení.
\item \textbf{Ekvifrekvenční} (\emph{equi-depth}) \mimo{} --- intervaly se stejným počtem
      objektů.
\item \textbf{S~využitím informace o~třídě} --- hranice se volí tak, aby intervaly byly co
      nejčistší; kritériem je entropie (rekurzivní dělení, jako v~C4.5) nebo $\chi^2$
      (slučování zdola, jako v~CHAID).
\end{itemize}

\subsubsection{Fuzzy diskretizace}

Ostré hranice pošlou dva blízké objekty do různých kategorií. \emph{Fuzzy diskretizace} je
\uv{rozostří}: jedné numerické hodnotě mohou odpovídat \emph{dvě} diskretizované hodnoty se
součtem charakteristických funkcí 1. Vznikají \uv{fuzzy}-objekty s~vahou $<1$ danou
\emph{součinem} charakteristických funkcí všech atributů:
\[ w = \prod_j \mu_j. \]
Rizika: ztráta informace a~\emph{kontradikce} --- objekty se stejnými diskretizovanými
hodnotami, ale různými třídami.

\begin{defbox}[Procedura \emph{Interval} (fuzzy diskretizace)]
\begin{enumerate}
\item Má-li posloupnost hodnot stejnou třídu, vytvoř z~ní interval s~charakteristickou
      funkcí rovnou 1 na celém intervalu.
\item Pro každý interval $\mathit{Int}_i$ patřící do třídy \uv{?}:
      \begin{itemize}
      \item patří-li sousední intervaly $\mathit{Int}_{i-1}$ a~$\mathit{Int}_{i+1}$ do
            \emph{téže} třídy, spoj je do jednoho intervalu
            $\mathit{Int}_{i-1}\cup \mathit{Int}_i\cup \mathit{Int}_{i+1}$
            (charakteristická funkce $=1$ na celém intervalu);
      \item jinak vytvoř \emph{dva fuzzy-intervaly}: spojením $\mathit{Int}_{i-1}\cup
            \mathit{Int}_i$ s~charakteristickou funkcí rovnou 1 na
            $[L_{i-1},U_{i-1}]$ a~lineárně klesající
            $\frac{U_{i+1}-x}{U_{i+1}-L_{i-1}}$ na přechodové oblasti, a~analogicky
            spojením $\mathit{Int}_i\cup \mathit{Int}_{i+1}$ s~funkcí lineárně rostoucí
            $\frac{x-L_{i-1}}{U_{i+1}-L_{i-1}}$.
      \end{itemize}
\item Zajisti spojité pokrytí celého rozsahu (mezery mezi intervaly se chápou jako
      intervaly třídy \uv{?}).
\end{enumerate}
\end{defbox}

\subsubsection{Seskupování hodnot kategoriálních atributů (KEX)}

Má-li kategoriální atribut příliš mnoho hodnot, je užitečné je seskupit. \textbf{Cíl:}
vytvořit takové grupy $A(\mathit{Grp})$, aby $P(\mathrm{Class}\mid A(\mathit{Grp}))$
\emph{významně} lišilo od apriorního $P(\mathrm{Class})$ --- a~vytvořit jich tolik, kolik je
tříd, plus jednu navíc (\uv{?}).

\begin{defbox}[Algoritmus seskupování hodnot (KEX)]
\textbf{Hlavní cyklus:}
\begin{enumerate}
\item vytvoř uspořádaný seznam hodnot uvažovaného atributu;
\item pro každou hodnotu: \emph{(a)}~spočítej četnosti výskytu objektů s~touto hodnotou
      v~jednotlivých třídách; \emph{(b)}~přiřaď hodnotě kód třídy procedurou
      \emph{Assign};
\item vytvoř grupy hodnot procedurou \emph{Group}.
\end{enumerate}

\textbf{Procedura \emph{Assign}:} patří-li pro danou hodnotu všechny objekty do téže třídy,
přiřaď hodnotě kód této třídy; jinak --- liší-li se rozdělení tříd pro danou hodnotu
\emph{významně} (podle $\chi^2$ testu) od apriorního rozdělení tříd --- přiřaď kód
nejčastější třídy, a~pokud ne, přiřaď kód \uv{?}.

\textbf{Procedura \emph{Group}:} vytvoř grupu z~hodnot, kterým byl přiřazen stejný kód
třídy.
\end{defbox}

\subsection{Redukce dat}

\subsubsection{Příliš mnoho objektů}

Dávkové zpracování se stává problematickým. Řešení: \emph{(a)}~\textbf{vzorek} objektů;
\emph{(b)}~uložení dovolující přístup bez nutnosti mít data celá v~paměti; \emph{(c)}~\textbf{více
modelů nad podmnožinami} a~jejich kombinace (ensembly, kap.~\ref{sec:metody}). Klíčová je
\textbf{reprezentativnost} vzorku --- rozdělení hodnot atributů na vzorku má souhlasit
s~rozdělením na celých datech.

\textbf{Nevyvážené třídy} vedou k~preferenci majoritní třídy. Řešení: různé \emph{váhy} typů
chyb (cenově citlivé učení) nebo výběr z~tříd s~\emph{různou pravděpodobností} (převzorkování
minority, podvzorkování majority).

\subsubsection{Příliš mnoho atributů}

Dvě principiálně odlišné cesty:

\begin{itemize}
\item \textbf{Redukce transformací} --- z~existujících atributů vytvoříme \emph{méně nových},
      typicky jako lineární kombinace původních (analýza hlavních komponent, PCA). Nevýhody:
      je nutné znát hodnoty \emph{všech} původních atributů a~nové atributy nemusí mít
      jasnou interpretaci.
\item \textbf{Redukce selekcí} (\emph{feature selection}) --- z~existujících atributů
      vybereme jen ty nejdůležitější. Interpretace se zachovává a~nepotřebné atributy se
      vůbec nemusí měřit.
\end{itemize}

\mimo{} PCA hledá ortogonální směry maximálního rozptylu (vlastní vektory kovarianční
matice s~největšími vlastními čísly); je to \emph{nesupervizovaná} transformace, takže může
zahodit směr, který je pro predikci klíčový, ale má malý rozptyl.

\subsection{Analýza relevance atributů: filtry a~wrappery}
\label{ssec:filtr}

Hledáme atributy, které nejlépe přispívají ke správné klasifikaci objektů.

\begin{defbox}[Filtry vs.\ wrappery]
\textbf{Filtrační metody} (\emph{filter}) --- pro každý atribut spočítáme charakteristiku
odrážející jeho přínos ke klasifikaci, atributy podle ní seřadíme a~vybereme určitý počet
nejlepších. Charakteristika je \emph{nezávislá} na následně použitém modelu; metoda je
rychlá.

\textbf{Obalové metody} (\emph{wrapper}) --- pomocí algoritmu strojového učení se postaví
model nad \emph{podmnožinou} atributů a~model se vyhodnotí; použije se nejlepší
z~postavených modelů (s~odpovídající podmnožinou atributů). Metoda je výpočetně drahá, ale
zohledňuje interakci atributů \emph{s~konkrétním} modelem.

Prohledávání prostoru podmnožin:
\begin{itemize}
\item \uv{zdola nahoru} (\emph{bottom-up}, \emph{dopředná selekce}) --- začneme od modelů
      nad jednotlivými atributy a~postupně přidáváme další;
\item \uv{shora dolů} (\emph{top-down}, \emph{zpětná eliminace}) --- začneme od modelu nad
      úplnou množinou atributů a~postupně odstraňujeme nepotřebné.
\end{itemize}
\end{defbox}

\textbf{Zásadní omezení filtrů:} každý atribut se hodnotí \emph{samostatně}, takže vzájemný
vliv atributů se zachycuje obtížně. Kritérium lze počítat pro \emph{množinu} atributů, ale
kombinatorická explozivnost to omezuje na hladové dopředné či zpětné prohledávání.
\mimo{} Důsledek: filtr propustí dvojici silně korelovaných atributů (oba mají vysoké skóre,
druhý ale nepřináší nic nového) a~zahodí atribut užitečný jen \emph{v~kombinaci} s~jiným
(XOR-vztah).

\subsubsection{Kritéria na kontingenčních tabulkách}
\label{ssec:kontingencni}

Vztah vstupního atributu $A$ (hodnoty $v_1,\dots,v_R$) a~cílového atributu $C$ (třídy
$v_1,\dots,v_S$) se popíše kontingenční tabulkou o~četnostech $a_{kl}$ kombinace
$A(v_k)\wedge C(v_l)$, s~marginály
\[ r_k=\sum_{l=1}^{S}a_{kl},\qquad s_l=\sum_{k=1}^{R}a_{kl},\qquad
n=\sum_{k=1}^{R}\sum_{l=1}^{S}a_{kl}. \]
Odhady pravděpodobností: $P(A(v_k)\wedge C(v_l))=a_{kl}/n$, $P(A(v_k))=r_k/n$,
$P(C(v_l))=s_l/n$. Očekávaná četnost při nezávislosti je $e_{kl}=r_k s_l/n$.

\begin{defbox}[$\chi^2$ jako kritérium relevance]
\[ \chi^2=\sum_{k=1}^{R}\sum_{l=1}^{S}\frac{(a_{kl}-e_{kl})^2}{e_{kl}}
   = n\sum_{k=1}^{R}\sum_{l=1}^{S}\frac{\big(a_{kl}-\tfrac{r_ks_l}{n}\big)^2}{r_ks_l}
   \qquad\text{(čím větší, tím lepší).}\]
Alternativní tvar pro ruční výpočet: $\chi^2=\sum_k\sum_l \frac{a_{kl}^2}{e_{kl}}-n$.

\textbf{Jako statistický test:} při platnosti nulové hypotézy o~nezávislosti
\[ H_0:\ P(A=v_k\wedge C=v_l)=P(A=v_k)\,P(C=v_l)\quad \forall k,l \]
má statistika $\chi^2$ rozdělení s~$(R-1)(S-1)$ stupni volnosti. Je-li
$\chi^2\ge\chi^2_{(R-1)(S-1)}(\alpha)$, nulovou hypotézu na hladině významnosti $\alpha$
zamítáme a~přijímáme alternativní hypotézu o~\emph{závislosti}.
\end{defbox}

\begin{defbox}[Příklad --- $\chi^2$ test na čtyřpolní tabulce]
\begin{center}
\begin{tabular}{@{}llrrr@{}}
\toprule
& & \multicolumn{2}{c}{Poskytnutí úvěru} & \\
& & ANO & NE & $r_k$ \\
\midrule
Výše příjmu & vysoký & 50 & 0 & 50 \\
            & nízký & 30 & 40 & 70 \\
\midrule
$s_l$ & & 80 & 40 & 120 \\
\bottomrule
\end{tabular}
\end{center}

Očekávané četnosti: $e_{11}=50\cdot 80/120=33{,}3$; $e_{12}=50\cdot 40/120=16{,}7$;
$e_{21}=70\cdot 80/120=46{,}7$; $e_{22}=70\cdot 40/120=23{,}3$. Rozdíly jsou tedy
$+16{,}7$, $-16{,}7$, $-16{,}7$, $+16{,}7$ a~hodnota statistiky
$\chi^2=42{,}857$.

Na hladině významnosti $\alpha=0{,}05$ je kritická hodnota $\chi^2_{(1)}(0{,}05)=3{,}84$.
Protože $42{,}857>3{,}84$, \textbf{mezi výší příjmu a~poskytnutím úvěru je závislost}.
\end{defbox}

\begin{defbox}[Fisherův exaktní test]
$\chi^2$ test lze použít jen pro dostatečně velké četnosti --- musí platit
$(r_k s_l)/n\ge 5$ pro všechna $k,l$. Pro čtyřpolní tabulky s~malými četnostmi se používá
\emph{Fisherův test}. Pravděpodobnost, že čtyřpolní tabulka má při daných marginálech
$r_1,r_2,s_1,s_2$ právě četnosti $a_{kl}$:
\[ p=\frac{r_1!\,r_2!\,s_1!\,s_2!}{n!\,a_{11}!\,a_{12}!\,a_{21}!\,a_{22}!}
     =\frac{\binom{s_1}{a_{11}}\binom{s_2}{a_{12}}}{\binom{n}{r_1}}. \]
Pravděpodobnosti se sečtou přes všechny uvažované hodnoty skutečných četností při daných
marginálech (nechť $m=\min(r_1,s_1)$):
\[ P=\sum_{i=0}^{m}\frac{r_1!\,r_2!\,s_1!\,s_2!}{n!\,(a_{11}-i)!\,(a_{12}+i)!\,(a_{21}+i)!\,(a_{22}-i)!}. \]
Je-li $P\le\alpha$, nulová hypotéza se na hladině $\alpha$ zamítá.

\textbf{Příklad.} Pro tabulku s~marginály $r=(12,18)$, $s=(6,24)$, $n=30$ vychází
$P_0=0{,}031$ a~$P_1=0{,}173$, tedy $P=P_0+P_1=0{,}204>0{,}05$ ---
\textbf{nulovou hypotézu o~nezávislosti přijímáme}. (Kontrola: součet $P_0,\dots,P_6$
$=0{,}031+0{,}173+0{,}340+0{,}302+0{,}128+0{,}024+0{,}002=1{,}000$.)
\end{defbox}

\subsubsection{Kritéria založená na entropii}

\begin{defbox}[Entropie, vzájemná informace, informační míra závislosti]
\textbf{Entropie} atributu $A$ jako vážený průměr entropií jeho hodnot
(čím \emph{menší}, tím lepší):
\[ H(A)=\sum_{k=1}^{R}\frac{r_k}{n}\,H\big(A(v_k)\big),\qquad
   H\big(A(v_k)\big)=-\sum_{l=1}^{S}\frac{a_{kl}}{r_k}\log_2\frac{a_{kl}}{r_k}. \]

\textbf{Vzájemná informace} (\emph{mutual information}):
\[ \mathrm{MI}(A,C)=\sum_{k=1}^{R}\sum_{l=1}^{S}\frac{a_{kl}}{n}
   \log_2\frac{n\,a_{kl}}{r_k s_l}. \]

\textbf{Informační míra závislosti} (\emph{information measure of dependence}) ---
normalizovaná vzájemná informace (čím \emph{větší}, tím lepší):
\[ \mathrm{ID}(A,C)=\frac{\mathrm{MI}(A,C)}{H(C)}
   =\frac{\mathrm{MI}(A,C)}{-\sum_{l=1}^{S}\frac{s_l}{n}\log_2\frac{s_l}{n}}. \]
\end{defbox}

Normalizace v~$\mathrm{ID}$ je podstatná: samotná $\mathrm{MI}$ systematicky zvýhodňuje
atributy s~mnoha hodnotami (v~limitě je identifikátor objektu \uv{nejinformativnější}
atribut). Totéž zvýhodnění řeší poměrný informační zisk v~C4.5 (kap.~\ref{sec:metody}).

\subsection{Vztahy mezi atributy: regrese a~korelace}
\label{ssec:regrese}

Zatímco kontingenční tabulky popisují vztah \emph{kategoriálních} atributů, u~numerických se
používá regresní a~korelační analýza. \emph{Korelační analýza} odpovídá na otázku, zda mezi
dvěma numerickými veličinami \emph{existuje} lineární vztah; \emph{lineární regrese} určuje
jeho \emph{parametry}.

\textbf{Důvody pro regresní analýzu:} nákladné měření výstupu (predikovat jej z~dostupných
vstupů); vstupy známé dříve než výstup (odhad); řízené vstupy pomohou najít chování výstupu;
hledaný kauzální vztah mezi vstupem a~výstupem.

\begin{defbox}[Lineární regrese jedné proměnné metodou nejmenších kvadrátů]
Vzory $[x_1,y_1],\dots,[x_p,y_p]$, regresní rovnice $y=\beta x+\alpha$, kvadratická chyba
\[ E=\sum_{i=1}^{p}\varepsilon_i^2=\sum_{i=1}^{p}\big(y_i-\hat y_i\big)^2
   =\sum_{i=1}^{p}\big(y_i-\beta x_i-\alpha\big)^2. \]
Položíme parciální derivace nule:
\[ \frac{\partial E}{\partial\alpha}=-2\sum_i\big(y_i-\beta x_i-\alpha\big)=0,\qquad
   \frac{\partial E}{\partial\beta}=-2\sum_i\big(y_i-\beta x_i-\alpha\big)x_i=0, \]
odkud po úpravě $\alpha=\bar y-\beta\bar x$ a
\[ \beta=\frac{\sum_{i=1}^{p}(x_i-\bar x)(y_i-\bar y)}{\sum_{i=1}^{p}(x_i-\bar x)^2}
   \qquad(\text{tj.\ výběrová kovariance dělená výběrovým rozptylem }x). \]
Predikce: $\hat y=\beta x+\alpha$.
\end{defbox}

\begin{defbox}[Korelační koeficient]
Posouzení \uv{množství} lineární závislosti:
\[ \rho(x,y)=\frac{\sigma_{xy}}{\sqrt{\sigma_x^2\sigma_y^2}}=\frac{\sigma_{xy}}{\sigma_x\sigma_y},
\qquad
\sigma_{xy}=\frac{1}{n-1}\sum_i (x_i-\bar x)(y_i-\bar y), \]
\[ \sigma_x^2=\frac{1}{n-1}\sum_i (x_i-\bar x)^2,\qquad
   \sigma_y^2=\frac{1}{n-1}\sum_i (y_i-\bar y)^2. \]
Platí $\rho\in[-1,1]$; $|\rho|=1$ znamená přesný lineární vztah, $\rho=0$ žádný
\emph{lineární} vztah (nikoli žádnou závislost --- $\rho$ nezachytí nelineární vztahy).
\end{defbox}

\textbf{Vícerozměrná lineární regrese.} Předpokládáme lineární závislost vysvětlované
veličiny $Y$ na několika vysvětlujících veličinách $X_1,\dots,X_m$:
$\hat y=\beta_0+\beta_1x_1+\cdots+\beta_mx_m$, maticově $\hat Y=X\beta$, kde $X$ je
rozšířená matice vstupních vzorů (první sloupec jedničky) a~$\beta=(\beta_0,\dots,\beta_m)$.
Kvadratická odchylka $E=(Y-X\beta)^{\!\top}(Y-X\beta)$; z~$\partial E/\partial\beta=0$
dostáváme normální rovnice $X^{\!\top}\!X\beta=X^{\!\top}Y$, tedy
\[ \beta=\big(X^{\!\top}\!X\big)^{-1}X^{\!\top}Y. \]
Pro složité praktické úlohy je tento výpočet nákladný $\Rightarrow$ používají se aproximativní
(iterativní) řešení.

\textbf{Nelineární regrese} předpokládá složitější funkční závislost (kvadratickou,
exponenciální, \dots). Nejdůležitějším případem je \emph{logistická regrese}, která
modeluje kategoriální (dvouhodnotový) výstup --- protože jde současně o~klasifikátor, je
zpracována v~kap.~\ref{sec:metody}, odd.~\ref{ssec:logreg}.

\subsection{Shrnutí ke~zkoušce}

\begin{itemize}
\item \textbf{Cíl přípravy dat}: vybrat relevantní data, reprezentovat je pro zvolený
      algoritmus, výsledkem \emph{jedna} tabulka objekty $\times$ atributy.
\item \textbf{Chybějící hodnoty}: ignorovat objekt / nová hodnota \uv{nevím} / existující
      hodnota (modus, poměrně, libovolně) / doplnit modelem. Umět říct, kdy je která volba
      vhodná (absence jako informace $\to$ \uv{nevím}).
\item \textbf{Odlehlé hodnoty}: $\mathrm{IQR}=Q_3-Q_1$, outlier mimo
      $[Q_1-1{,}5\,\mathrm{IQR},\,Q_3+1{,}5\,\mathrm{IQR}]$. Umět spočítat kvartily
      z~vzestupně seřazených dat (pravidlo $n/4$ a~$3n/4$, zaokrouhlení nahoru) a~nakreslit
      box plot.
\item \textbf{Standardizace}: desítkové škálování, min-max do $[0,1]$ a~$[-1,1]$, z-skóre
      podle střední absolutní odchylky (robustnější), std-skóre podle směrodatné odchylky.
      Motivace: bez ní vzdálenost ovládne atribut s~největším rozsahem.
\item \textbf{Diskretizace}: ekvidistantní, ekvifrekvenční, s~využitím třídy (entropie,
      $\chi^2$); \textbf{fuzzy diskretizace} --- rozostřené hranice, jedné hodnotě odpovídají
      dvě kategorie se součtem charakteristických funkcí 1, váha fuzzy-objektu je součin
      $\prod_j\mu_j$; rizika: ztráta informace a~kontradikce. \textbf{KEX} seskupuje hodnoty
      kategoriálního atributu podle $\chi^2$ testu proti apriornímu rozdělení tříd
      (procedury \emph{Assign} a~\emph{Group}).
\item \textbf{Redukce dat}: příliš mnoho objektů $\to$ vzorkování (reprezentativnost!),
      efektivní uložení, více modelů nad podmnožinami; nevyvážené třídy $\to$ váhy chyb nebo
      různé pravděpodobnosti výběru. Příliš mnoho atributů $\to$ \emph{transformace} (PCA:
      méně nových atributů jako lineární kombinace, ale nutné znát všechny původní a~chybí
      interpretace) vs.\ \emph{selekce} (podmnožina původních, interpretace zůstává).
\item \textbf{Filtry vs.\ wrappery}: filtr počítá charakteristiku pro každý atribut
      samostatně a~je nezávislý na modelu (rychlý, ale nezachytí interakce); wrapper staví
      model nad podmnožinou a~vyhodnocuje jej (drahý, zohledňuje interakce). Prohledávání
      \emph{bottom-up} (přidávání) nebo \emph{top-down} (odebírání).
\item \textbf{Filtrační kritéria}: $\chi^2$ (větší $=$ lepší), entropie $H(A)$ (menší $=$
      lepší), $\mathrm{MI}(A,C)$, $\mathrm{ID}(A,C)=\mathrm{MI}/H(C)$ (větší $=$ lepší).
      Umět zdůvodnit normalizaci $\mathrm{ID}$: bez ní jsou zvýhodněny atributy s~mnoha
      hodnotami.
\item \textbf{$\chi^2$ test}: $(R-1)(S-1)$ stupňů volnosti, podmínka $r_ks_l/n\ge 5$; pro
      malé četnosti \textbf{Fisherův exaktní test}. Umět projít výpočet na čtyřpolní tabulce
      (příklad s~úvěrem: $\chi^2=42{,}857>3{,}84$ $\Rightarrow$ závislost).
\item \textbf{Regrese a~korelace}: metoda nejmenších kvadrátů, $\beta$ jako kovariance dělená
      rozptylem, $\alpha=\bar y-\beta\bar x$; vícerozměrně
      $\beta=(X^{\!\top}\!X)^{-1}X^{\!\top}Y$; korelační koeficient
      $\rho=\sigma_{xy}/(\sigma_x\sigma_y)$ měří jen \emph{lineární} závislost.
\end{itemize}

%=====================================================================
\newpage
\section{Metody pro dobývání znalostí}
\label{sec:metody}
%=====================================================================

Nejobsáhlejší věta okruhu jmenuje tři skupiny metod pro tři situace: \emph{asociační pravidla}
(transakční data bez cílového atributu), \emph{učení s~učitelem} (data s~cílovým atributem)
a~\emph{klastrová analýza} (metrická data bez cílového atributu). Podle nich je kapitola
rozdělena.

\subsection{Asociační pravidla a~frekventované vzory}
\label{ssec:ar}

\subsubsection{Analýza nákupního košíku}

\emph{Analýza nákupního košíku} (\emph{market basket analysis}, MBA) je archetypální úloha:
\uv{které položky se v~košíku vyskytují společně?} Výsledky se vyjadřují ve formě pravidel
a~lze je okamžitě použít --- při plánování a~rozvržení prodejny, nabídce kuponů a~časově
omezených slev, nebo \uv{balíčkování} produktů.

\textbf{Položka} (\emph{item}) je nabízený produkt nebo služba, \textbf{transakce} obsahuje
jednu nebo více položek. \textbf{Frekvenční tabulka} udává počet společných výskytů každých
dvou položek v~transakcích; na její diagonále je počet transakcí obsahujících danou položku.

\begin{defbox}[Příklad --- frekvenční tabulka]
Transakce v~obchodě s~potravinami:
\begin{center}
\small
\begin{tabular}{@{}rl@{\hspace{1.2cm}}lrrrr@{}}
\toprule
Zák. & Položky & & chléb & máslo & mléko & káva \\
\midrule
1 & chléb, máslo          & chléb & 4 & 3 & 1 & 2 \\
2 & mléko, chléb, máslo   & máslo & 3 & 4 & 1 & 2 \\
3 & chléb, káva           & mléko & 1 & 1 & 1 & 0 \\
4 & chléb, máslo, káva    & káva  & 2 & 2 & 0 & 3 \\
5 & káva, máslo           &       &   &   &   &   \\
\bottomrule
\end{tabular}
\end{center}
Z~tabulky je povaha prodeje vidět okamžitě: \emph{chléb a~máslo se kupují spolu}, \emph{mléko
se nikdy nekupuje s~kávou}.
\end{defbox}

Pravidla mají formu \texttt{IF} \emph{podmínka} \texttt{THEN} \emph{závěr}. Mají být
\textbf{snadno pochopitelná} (vztah lze ověřit) a~\textbf{použitelná} (vedou k~zásahu),
naopak nemají být \textbf{triviální} (každý to zná) ani \textbf{nevysvětlitelná} (nemají
vysvětlení a~nevedou k~akci).

\subsubsection{Support, confidence a~lift}

\begin{defbox}[Support, confidence, lift]
\textbf{Support} (podpora) --- jak často lze pravidlo použít:
\[ \mathrm{sup}(R)=\frac{\text{počet transakcí obsahujících } i \text{ a } j}
   {\text{počet všech transakcí}}\cdot 100\,\%. \]

\textbf{Confidence} (spolehlivost) --- jak spolehlivé jsou výsledky pravidla:
\[ \mathrm{conf}(R)=\frac{\text{počet transakcí obsahujících } i \text{ a } j}
   {\text{počet transakcí obsahujících } i}\cdot 100\,\%. \]

\textbf{Lift} (zdvih) --- kolikrát je lepší použít pravidlo pro predikci než jen předpokládat
jeho výsledek:
\[ \mathrm{lift}(R)=\frac{p(i\wedge j)}{p(i)\,p(j)}. \]
Je-li $\mathrm{lift}<1$, je pravidlo pro predikci \emph{horší než náhodná volba}
a~\textbf{negace závěru} může dát lepší pravidlo (\texttt{IF} \emph{podmínka} \texttt{THEN
NOT} \emph{závěr}).
\end{defbox}

\begin{defbox}[Příklad --- support, confidence, lift]
Nad daty z~předchozího příkladu (5 transakcí):

\emph{Pravidlo 1:} \uv{kdo kupuje chléb, kupuje i~máslo.}
$\mathrm{sup}=3/5=60\,\%$, $\mathrm{conf}=3/4=75\,\%$.

\emph{Pravidlo 2:} \uv{kdo kupuje kávu, kupuje i~máslo.}
$\mathrm{sup}=2/5=40\,\%$, $\mathrm{conf}=2/3=66\,\%$.

\emph{Pravidlo 3:} \uv{kdo kupuje mléko, kupuje i~máslo.}
$\mathrm{sup}=1/5=20\,\%$, $\mathrm{conf}=1/1=100\,\%$,
$\mathrm{lift}=\frac{1/5}{(1/5)\cdot(4/5)}=\frac54=1{,}25$.

Pravidlo 3 má \emph{stoprocentní} spolehlivost, ale support jen 20~\% --- opírá se o~jedinou
transakci. Lift 1,25 říká, že zdvih proti náhodě je jen malý (máslo kupuje 80~\% zákazníků
tak jako tak). To je základní past: \emph{vysoká confidence sama o~sobě nestačí}.
\end{defbox}

\subsubsection{Volba položek, taxonomie a~virtuální položky}

Transakční data bývají hůř kvalitní a~vyžadují rozsáhlejší předzpracování; relevance položek
se mění v~čase. Volba úrovně je kompromis mezi rostoucím počtem kombinací a~přímou
použitelností výsledků (konkrétní položky), při zachování dostatečného supportu.

\textbf{Taxonomie} je hierarchie kategorií: zpočátku obecnější položky, později pravidla pro
konkrétní položky, a~to jen z~transakcí, které je obsahují. Uvažované položky mají mít
přibližně stejnou frekvenci --- \emph{vzácné posunout výše} v~taxonomii, \emph{častější nechat
níže} (aby nejfrekventovanější nedominovaly pravidlům).

\textbf{Virtuální položky} přesahují taxonomii: rozostřují hranici mezi typy produktů (značky
--- Calvin Klein) a~nesou informaci o~transakci samotné --- anonymní (den, čas) nebo
podepsanou (zákazník a~jeho chování). Riziko \emph{redundance} v~pravidlech: virtuální položka
odpovídá právě jedné položce z~taxonomie (\uv{IF produkt\_Škoda THEN Škoda}), nebo se
v~pravidle sejde virtuální i~obecná položka (\uv{IF produkt\_Škoda AND malé auto THEN stan}
místo \uv{IF Fabia THEN stan}). \emph{Srovnání prodejen:} virtuální položka označí typ
prodejny (neodpovídá žádnému produktu) --- vezmi data z~nových prodejen a~přibližně stejně
velká data ze stávajících, aplikuj MBA na oba soubory a~zaměř se na pravidla s~virtuálními
položkami.

\textbf{Prořezávání podle minimálního supportu} eliminuje málo frekventované položky (akce se
musí týkat dost transakcí): buď eliminovat vzácné položky a~jejich pravidla, nebo taxonomií
vytvořit obecné položky splňující prah. Proměnný minimální support vede ke \emph{kaskádovému
efektu}.

\subsubsection{Algoritmus Apriori}

\begin{defbox}[Frekventovaná množina položek a~Apriori princip]
\emph{Frekventovaná množina položek} (\emph{frequent itemset}) je kombinace (konjunkce)
kategorií, která na datech dosáhne předem stanovené frekvence --- supportu \emph{minsup}.

\textbf{Vlastnost uzavřenosti dolů} (\emph{downward closure property}) --- klíčová idea
Apriori: \emph{každá podmnožina frekventované množiny je také frekventovaná}. Ekvivalentně:
je-li množina nefrekventovaná, jsou nefrekventované i~všechny její nadmnožiny.

Při hledání kombinací délky $k$ se proto využijí známé kombinace délky $k-1$ ---
\emph{prohledávání do šířky}. Aby byla kombinace délky $k$ vygenerována, požaduje se, aby
\emph{všechny} její podkombinace délky $k-1$ splňovaly požadavek na frekvenci.
\end{defbox}

\begin{defbox}[Algoritmus APRIORI (R.~Agrawal)]
\begin{algorithmic}[1]
\State $L_1\gets$ všechny kategorie, které dosahují alespoň požadované frekvence
\State $k\gets 2$
\While{$L_{k-1}\neq\emptyset$}
  \State $C_k\gets \textsc{Apriori-Gen}(L_{k-1})$ \Comment{generování kandidátů}
  \State $L_k\gets$ ty kombinace z~$C_k$, které dosáhly alespoň požadované frekvence
  \State $k\gets k+1$
\EndWhile
\end{algorithmic}

\textbf{Funkce \textsc{Apriori-Gen}$(L_{k-1})$:}
\begin{algorithmic}[1]
\For{každou dvojici kombinací $\mathit{Comb}_p,\mathit{Comb}_q\in L_{k-1}$}
  \If{$\mathit{Comb}_p$ a~$\mathit{Comb}_q$ se shodují v~$k-2$ kategoriích}
     \State přidej $\mathit{Comb}_p\cup \mathit{Comb}_q$ do $C_k$ \Comment{krok spojení}
  \EndIf
\EndFor
\For{každou kombinaci $\mathit{Comb}\in C_k$} \Comment{krok prořezání}
  \If{některá její podkombinace délky $k-1$ není v~$L_{k-1}$}
     \State odstraň $\mathit{Comb}$ z~$C_k$
  \EndIf
\EndFor
\end{algorithmic}
\end{defbox}

\textbf{Generování pravidel.} Po nalezení kombinací s~vyhovující frekvencí se vytvoří
asociační pravidla (platí \emph{frekvence nadkombinace} $\le$ \emph{frekvence kombinace}).
Každá kombinace $\mathit{Comb}$ se rozdělí na všechny možné dvojice podkombinací
$\mathit{Ant}$ a~$\mathit{Suc}$ tak, že
\[ \mathit{Suc}=\mathit{Comb}\setminus \mathit{Ant},\qquad
   \mathit{Ant}\cap \mathit{Suc}=\emptyset,\quad \mathit{Ant}\cup \mathit{Suc}=\mathit{Comb}. \]
Pro pravidlo $\mathit{Ant}\Rightarrow \mathit{Suc}$ pak \emph{support} odpovídá frekvenci
kombinace $\mathit{Comb}$ a~\emph{confidence} poměru frekvencí kombinací $\mathit{Comb}$
a~$\mathit{Ant}$.

\begin{center}
\small
\begin{tabular}{@{}p{7.6cm}p{8.0cm}@{}}
\toprule
\textbf{Výhody MBA} & \textbf{Nevýhody MBA} \\
\midrule
jasné a~pochopitelné výsledky --- IF-THEN pravidla okamžitě použitelná & výpočetní náklady:
  prostor asociačních pravidel je exponenciální, $O(2^m)$ pro $m$ položek \\
dobývání bez cílových výstupních hodnot (bez apriorní znalosti) & nutnost taxonomií
  a~virtuálních položek \\
umí zpracovat data proměnné délky & obtížné určení vhodného počtu položek; položky by měly
  mít přibližně stejnou frekvenci \\
jednoduchý a~snadno pochopitelný výpočet & \emph{produkuje nadměrné množství pravidel} ---
  tisíce, desetitisíce, miliony \\
\bottomrule
\end{tabular}
\end{center}

\subsubsection{Vícenásobný minimální support: MS-Apriori}

\textbf{Problém vzácných položek.} Jediný \emph{minsup} předpokládá, že všechny položky mají
podobnou povahu a~frekvenci. V~praxi se ale některé položky vyskytují velmi často a~jiné
vzácně (kuchyňský robot vs.\ chléb a~mléko). Důsledek:
\begin{itemize}
\item je-li \emph{minsup} příliš \textbf{vysoký}, pravidla se vzácnými položkami se nenajdou;
\item je-li příliš \textbf{nízký} (aby se našla pravidla s~vzácnými i~častými položkami),
      hrozí \emph{kombinatorická explozivnost} --- časté položky se pospojují všemi možnými
      způsoby.
\end{itemize}

\begin{defbox}[Model vícenásobných minimálních supportů (MIS)]
Minimální support pravidla se vyjádří pomocí \emph{minimálních supportů položek}
(\emph{minimum item support}, MIS) těch položek, které v~pravidle vystupují. Každá položka
může mít vlastní $\mathrm{MIS}(i)$; zadáním různých hodnot vyjádří uživatel různé požadavky
na support pro různá pravidla.

\textbf{Minsup pravidla} $R: a_1,\dots,a_k\Rightarrow a_{k+1},\dots,a_r$ je
\emph{nejnižší} hodnota MIS mezi jeho položkami: pravidlo splňuje minimální support, je-li
jeho skutečný support $\ge \min\big(\mathrm{MIS}(a_1),\dots,\mathrm{MIS}(a_r)\big)$.

\textbf{Omezení rozdílu supportů:} aby se v~téže množině neobjevily velmi časté a~velmi
vzácné položky, požaduje se
\[ \max_{i\in s}\{\mathrm{sup}(i)\}-\min_{i\in s}\{\mathrm{sup}(i)\}\le\varphi, \]
kde $0\le\varphi\le 1$ je uživatelem zadaný maximální rozdíl supportů (stejný pro všechny
množiny).
\end{defbox}

\begin{defbox}[Příklad --- minsup pravidla podle MIS]
Položky \emph{chléb}, \emph{obuv}, \emph{oděvy} s~$\mathrm{MIS}(\text{chléb})=2\,\%$,
$\mathrm{MIS}(\text{obuv})=0{,}1\,\%$, $\mathrm{MIS}(\text{oděvy})=0{,}2\,\%$.

$R_1$: oděvy $\Rightarrow$ chléb $[\mathrm{sup}=0{,}15\,\%]$ minsup \textbf{nesplňuje},
protože $0{,}15\,\%<\mathrm{MIS}(\text{oděvy})=0{,}2\,\%$.

$R_2$: oděvy $\Rightarrow$ obuv $[\mathrm{sup}=0{,}15\,\%]$ minsup \textbf{splňuje},
protože $0{,}15\,\%\ge\mathrm{MIS}(\text{obuv})=0{,}1\,\%$.
\end{defbox}

\textbf{Uzavřenost dolů v~novém modelu neplatí!} MIS totiž dovoluje vyšší minimální supporty
pro pravidla s~jen frekventovanými položkami a~nižší pro pravidla s~méně frekventovanými.
Příklad: $\mathrm{MIS}(1)=10\,\%$, $\mathrm{MIS}(2)=20\,\%$, $\mathrm{MIS}(3)=5\,\%$,
$\mathrm{MIS}(4)=6\,\%$. Množina $\{1,2\}$ se supportem 9~\% je nefrekventovaná, ale
$\{1,2,3\}$ a~$\{1,2,4\}$ frekventované být mohou --- $\{1,2\}$ tedy \emph{nesmíme}
zahodit.

\textbf{Řešení:} položky se seřadí \emph{vzestupně podle hodnot MIS} a~toto pořadí se pak
používá ve všech dalších operacích algoritmu. Každá množina $w$ má formu
$\{w[1],\dots,w[k]\}$, kde $\mathrm{MIS}(w[1])\le\cdots\le\mathrm{MIS}(w[k])$.

\begin{defbox}[Algoritmus MS-Apriori]
\begin{algorithmic}[1]
\State $M\gets \textsc{sort}(I,\mathit{MS})$ \Comment{setřídění položek podle $\mathrm{MIS}(i)$}
\State $L\gets \textsc{init-pass}(M,T)$ \Comment{první průchod daty; $L$ je množina semen}
\State $F_1\gets \{\{i\}\mid i\in L,\ i.\mathrm{count}/n\ge \mathrm{MIS}(i)\}$
\For{$k=2$; $F_{k-1}\neq\emptyset$; $k{+}{+}$}
  \If{$k=2$}
     \State $C_k\gets \textsc{level2-candidate-gen}(L,\varphi)$
  \Else
     \State $C_k\gets \textsc{MScandidate-gen}(F_{k-1},\varphi)$
  \EndIf
  \For{každou transakci $t\in T$}
     \For{každého kandidáta $c\in C_k$}
        \If{$c\subseteq t$} \State $c.\mathrm{count}{+}{+}$ \EndIf
        \If{$c\setminus\{c[1]\}\subseteq t$} \State $c.\mathrm{tailCount}{+}{+}$
           \Comment{pro generování pravidel} \EndIf
     \EndFor
  \EndFor
  \State $F_k\gets \{c\in C_k\mid c.\mathrm{count}/n\ge \mathrm{MIS}(c[1])\}$
\EndFor
\State \Return $F=\bigcup_k F_k$
\end{algorithmic}
\end{defbox}

\textbf{První průchod daty} (\textsc{init-pass}): algoritmus zjistí support každé položky,
pak v~setříděném pořadí najde \emph{první} položku $i$, která splní $\mathrm{MIS}(i)$,
a~vloží ji do množiny semen $L$. Pro každou další položku $j$ za $i$ platí: je-li
$j.\mathrm{count}/n\ge \mathrm{MIS}(i)$, patří $j$ také do $L$.

\emph{Příklad:} $\mathrm{MIS}(1)=10\,\%$, $\mathrm{MIS}(2)=20\,\%$, $\mathrm{MIS}(3)=5\,\%$,
$\mathrm{MIS}(4)=6\,\%$; 100 transakcí, četnosti $\{3\}=6$, $\{4\}=3$, $\{1\}=9$,
$\{2\}=25$. Pak $L=\{3,1,2\}$ a~$F_1=\{\{3\},\{2\}\}$. Položka 4 není v~$L$, protože
$4.\mathrm{count}/n<\mathrm{MIS}(3)=5\,\%$; $\{1\}$ není v~$F_1$, protože
$1.\mathrm{count}/n<\mathrm{MIS}(1)=10\,\%$.

\textbf{Generování kandidátů úrovně 2} (\textsc{level2-candidate-gen}): pro každou položku
$l\in L$ v~daném pořadí, splňuje-li $l.\mathrm{count}/n\ge\mathrm{MIS}(l)$, se pro každou
položku $h$ za $l$ přidá kandidát $\{l,h\}$, pokud
$h.\mathrm{count}/n\ge\mathrm{MIS}(l)$ \emph{a} $|\mathrm{sup}(h)-\mathrm{sup}(l)|\le\varphi$.
Podstatné je, že se používá $L$, ne $F_1$ --- $F_1$ totiž neobsahuje položky, které splňují
MIS nějaké dřívější položky, ale ne své vlastní (v~příkladu položka 1).

\textbf{Generování kandidátů vyšších úrovní} (\textsc{MScandidate-gen}) má dvě části.
\emph{Krok spojení} je jako v~Apriori: spojí se $f_1=\{i_1,\dots,i_{k-2},i_{k-1}\}$
a~$f_2=\{i_1,\dots,i_{k-2},i'_{k-1}\}$, které se liší jen v~poslední položce, pokud
$i_{k-1}<i'_{k-1}$ a~$|\mathrm{sup}(i_{k-1})-\mathrm{sup}(i'_{k-1})|\le\varphi$.
\emph{Krok prořezání}: pro každou $(k-1)$-podmnožinu $s$ kandidáta $c$ --- není-li $s\in
F_{k-1}$, lze $c$ smazat, \textbf{s~jedinou výjimkou}: neobsahuje-li $s$ položku $c[1]$
(taková $s$ je právě jedna), tedy položku s~nejnižší hodnotou MIS, nesmíme $c$ smazat ---
nevíme totiž, zda $s$ nesplňuje $\mathrm{MIS}(c[1])$, víme jen, že nesplňuje
$\mathrm{MIS}(c[2])$; to platí, jen pokud $\mathrm{MIS}(c[2])=\mathrm{MIS}(c[1])$.

\textbf{Proč \texttt{tailCount}.} U~Apriori platí: je-li $f$ frekventovaná a~$f_{sub}\subseteq
f$, je i~$f_{sub}$ frekventovaná --- k~určení confidence pravidel tedy stačí supporty
frekventovaných množin. U~MS-Apriori to nestačí; problém vzniká, když by v~závěru pravidla
byla položka s~\emph{nejnižší} hodnotou MIS. Proto se zvlášť počítá
$c.\mathrm{tailCount}$ --- support množiny $c$ bez první položky.

\textbf{Jak zadat hodnoty MIS:} \emph{(a)}~$\lambda\cdot\mathrm{sup}(i)$ podle skutečného
supportu položky ($0\le\lambda\le 1$ stejné pro všechny); \emph{(b)}~seskupit položky do bloků
s~podobnými frekvencemi a~dát bloku jednu hodnotu MIS. Trik: $\mathrm{MIS}=100\,\%$ generování
pravidel s~těmito položkami zakáže. Model vícenásobných minsupů \emph{zahrnuje} model jediného
supportu jako speciální případ.

\subsubsection{Třídní asociační pravidla (CAR)}

Běžné dolování necílí na žádný atribut --- najde všechna pravidla a~libovolná položka může být
podmínkou i~závěrem. Někdy ale uživatele zajímají právě \emph{cílové} atributy či třídy (které
slovo je asociováno s~kterým tématem dokumentu).

\begin{defbox}[Třídní asociační pravidlo]
Nechť $T$ je množina $n$ transakcí, každá navíc označkovaná třídou $y$; $I$ je množina všech
položek, $Y$ množina všech tříd a~$I\cap Y=\emptyset$.

\emph{Třídní asociační pravidlo} (\emph{class association rule}, CAR) je implikace
$X\to y$, kde $X\subseteq I$ a~$y\in Y$. Support a~confidence jsou definovány stejně jako
u~běžných asociačních pravidel.
\end{defbox}

\begin{defbox}[Příklad --- CAR na textové kolekci]
\begin{center}
\small
\begin{tabular}{@{}ll@{}}
\toprule
Dokument & Třída \\
\midrule
1: Student, Učit, Škola & Vzdělávání \\
2: Student, Škola & Vzdělávání \\
3: Učit, Škola, Město, Hra & Vzdělávání \\
4: Baseball, Basketbal & Sport \\
5: Basketbal, Hráč, Divák & Sport \\
6: Baseball, Trenér, Hra, Tým & Sport \\
7: Basketbal, Tým, Město, Hra & Sport \\
\bottomrule
\end{tabular}
\end{center}
Pro $\mathit{minsup}=20\,\%$ a~$\mathit{minconf}=60\,\%$ jsou příklady CAR:
\[ \text{Student, Škola}\to\text{Vzdělávání}\ [\mathrm{sup}=2/7,\ \mathrm{conf}=2/2],\qquad
   \text{Hra}\to\text{Sport}\ [\mathrm{sup}=2/7,\ \mathrm{conf}=2/3]. \]
\end{defbox}

CAR lze dolovat \textbf{přímo v~jednom kroku} (na rozdíl od běžných asociačních pravidel).
Klíčovou operací je nalezení všech \emph{pravidlových položek} (\emph{rule-items}) tvaru
$(\mathit{condset},y)$, kde $\mathit{condset}\subseteq I$ a~$y\in Y$, se supportem nad
\emph{minsup}; každá pravidlová položka reprezentuje pravidlo
$\mathit{condset}\to y$. Apriori lze na generování CAR snadno upravit.

Také zde lze použít ideu vícenásobných minimálních supportů: uživatel zadá
\emph{různé minimální supporty různým třídám} (tím efektivně různý minsup pravidlům každé
třídy) --- např.\ pravidla třídy \emph{Ano} s~minsup 5~\% a~třídy \emph{Ne} s~minsup 10~\%.
Nastavením minimálního třídního supportu na 101~\% se generování pravidel pro danou třídu
zakáže.

\subsubsection{Dolování sekvenčních vzorů}

Asociační pravidla neuvažují \emph{pořadí} transakcí. V~mnoha aplikacích je ale pořadí
podstatné: v~MBA je zajímavé vědět, zda lidé kupují položky v~posloupnosti (nejprve postel,
o~nějakou dobu později povlečení); v~dolování z~webových logů (\emph{web usage mining}) se
hledají navigační vzory podle posloupností návštěv stránek.

\begin{defbox}[Sekvence, velikost, délka, podsekvence]
Nechť $I=\{i_1,\dots,i_m\}$ je množina položek. \emph{Sekvence} je uspořádaný seznam
množin položek. \emph{Množina položek} (\emph{prvek} sekvence) je neprázdná
$X\subseteq I$; sekvenci $s$ značíme $\langle a_1a_2\dots a_r\rangle$, kde $a_i$ je prvek.
Bez újmy na obecnosti předpokládáme položky v~prvku v~lexikografickém pořadí.

\emph{Velikost} sekvence je počet jejích prvků, \emph{délka} počet položek; sekvence délky
$k$ se nazývá $k$-sekvence.

$s_1=\langle a_1\dots a_r\rangle$ je \emph{podsekvencí} $s_2=\langle b_1\dots b_v\rangle$
(a~$s_2$ \emph{nadsekvencí} $s_1$, též $s_2$ \emph{obsahuje} $s_1$), existují-li celá čísla
$1\le j_1<j_2<\cdots<j_r\le v$ taková, že $a_1\subseteq b_{j_1},\dots,a_r\subseteq b_{j_r}$.
\end{defbox}

\emph{Příklad.} Pro $I=\{1,\dots,9\}$ je $\langle\{3\}\{4,5\}\{8\}\rangle$ podsekvencí
$\langle\{6\}\{3,7\}\{9\}\{4,5,8\}\{3,8\}\rangle$, protože $\{3\}\subseteq\{3,7\}$,
$\{4,5\}\subseteq\{4,5,8\}$ a~$\{8\}\subseteq\{3,8\}$. Naopak
$\langle\{3\}\{8\}\rangle$ \emph{není} obsažena v~$\langle\{3,8\}\rangle$ ani naopak.
Velikost $\langle\{3\}\{4,5\}\{8\}\rangle$ je 3, délka 4.

\textbf{Úloha:} pro danou množinu $S$ vstupních datových sekvencí (sekvenční databázi) najít
všechny sekvence, které mají uživatelem zadaný minimální support; taková sekvence se nazývá
\emph{frekventovaná sekvence} nebo \emph{sekvenční vzor}. Support sekvence je podíl
datových sekvencí v~$S$, které ji obsahují.

Algoritmus \textbf{GSP} (\emph{Generalized Sequential Pattern}) generuje sekvenční vzory
způsobem obdobným Apriori --- prohledávání do šířky s~generováním a~testováním kandidátů,
kde kandidáti délky $k$ vznikají spojováním frekventovaných sekvencí délky $k-1$.

\textbf{Analýza časových řad pomocí MBA} --- \uv{příčina--následek}: zavedou se nové položky
pro transakce \emph{před} uvažovanou událostí (příčiny) a~\emph{po} ní (následky), duplicity
se odstraní. \emph{Okno} je snímek dat za období (transakce z~posledního měsíce) --- zachycuje
trendy u~vzácných položek. Může být nutná identifikace zákazníka.

\textbf{Disociační pravidla} \uv{IF $A$ AND NOT $B$ THEN $C$}: zavedou se položky
komplementární k~původním. Nevýhody: dvojnásobný počet položek, větší transakce a~negované
položky jsou častější než původní (pravidla s~vše\-mi negovanými se obtížně využívají).

\subsubsection{FP-Growth a~FP-strom}

Apriori používá přístup \emph{generuj a~testuj}: generování kandidátů je nákladné (prostorově
i~časově) a~nákladné je i~počítání supportu --- kontrola podmnožin je výpočetně drahá
a~databáze se prochází \emph{opakovaně} (I/O). \textbf{FP-Growth} umožňuje hledat
frekventované množiny \emph{bez generování kandidátů} ve dvou krocích: \emph{(1)}~postavit
kompaktní datovou strukturu \emph{FP-strom} (dvěma průchody daty); \emph{(2)}~extrahovat
frekventované množiny přímo z~FP-stromu jeho procházením.

\begin{defbox}[FP-strom]
Uzly FP-stromu odpovídají položkám a~mají počítadlo. FP-Growth čte transakce po jedné
a~mapuje každou na \emph{cestu} ve stromu. Používá se \emph{fixní pořadí} položek, takže
cesty se mohou překrývat, mají-li transakce stejný \emph{prefix}; v~takovém případě se
počítadla inkrementují. Mezi uzly obsahujícími tutéž položku se udržují ukazatele, které
tvoří jednosměrné spojové seznamy. Čím více cest se překrývá, tím vyšší je komprese
a~tím spíše se FP-strom vejde do paměti.
\end{defbox}

\textbf{Konstrukce (2 průchody).} \emph{Průchod 1:} zjisti support každé položky, zahoď
nefrekventované, seřaď frekventované \emph{klesajícím} supportem (např.\ $a,b,c,d,e$) --- toto
pořadí maximalizuje sdílení prefixů. \emph{Průchod 2:} postav strom. Např.\ $\{a,b\}$ dá cestu
\texttt{null}$\to a\to b$ s~počítadly 1; $\{b,c,d\}$ dá \emph{disjunktní} cestu
\texttt{null}$\to b\to c\to d$ (bez společného prefixu) a~ukazatel mezi oba uzly $b$;
$\{a,c,d,e\}$ má společný prefix $a$, takže se cesty překryjí a~počítadlo $a$ se inkrementuje.

\textbf{Velikost FP-stromu.} Obvykle menší než původní data (transakce sdílejí prefixy).
\emph{Nejlepší případ:} všechny transakce stejné $\Rightarrow$ jediná cesta. \emph{Nejhorší:}
každá transakce unikátní $\Rightarrow$ strom je alespoň tak velký jako data a~nároky jsou
\emph{vyšší} (ukazatele a~počítadla). Velikost závisí na pořadí položek --- klesající support
je jen \emph{heuristika}, nemusí dát nejmenší strom.

\begin{defbox}[Extrakce frekventovaných množin z~FP-stromu]
Postup \emph{zdola nahoru}, od listů ke kořeni, metodou \emph{rozděl a~panuj}: nejprve množiny
končící na $e$, pak na $de$, \dots, pak na $d$, $cd$ atd. Pro každou položku:
\begin{enumerate}
\item získej \textbf{prefixový podstrom} (\emph{prefix path sub-tree}) pomocí spojových seznamů;
\item ověř frekventovanost sečtením počítadel po spojovém seznamu (v~příkladu pro
      $\mathit{minSup}=2$ je $\mathrm{count}(e)=3$, takže $\{e\}$ je frekventovaná);
\item je-li frekventovaná, hledej rekurzivně množiny končící na $de$, $ce$, $be$, $ae$ pomocí
      \emph{podmíněného FP-stromu}.
\end{enumerate}

\textbf{Podmíněný FP-strom} (\emph{conditional FP-tree}) pro $e$ je strom, který by vznikl, kdyby
se uvažovaly \emph{jen transakce obsahující} $e$ a~$e$ se z~nich odstranilo. Z~prefixového
podstromu se získá ve třech krocích: aktualizuj počítadla tak, aby odrážela počet transakcí
s~$e$; odstraň uzly $e$; odstraň nefrekventované položky (má-li $b$ support 1, má jej i~$be$,
takže $b$ lze odstranit). Pak se z~něj hledají množiny končící na $de$, $ce$, $ae$ ($be$ ne ---
$b$ ve stromu není) a~pro každou se rekurzivně opakuje týž postup (např.\ $e\to de\to ade$).
\end{defbox}

\begin{center}
\small
\begin{tabular}{@{}p{7.6cm}p{8.0cm}@{}}
\toprule
\textbf{Výhody FP-Growth} & \textbf{Nevýhody FP-Growth} \\
\midrule
jen 2 průchody daty & FP-strom se nemusí vejít do paměti \\
\uv{komprimuje} data & stavba FP-stromu je nákladná (ale poté se množiny čtou snadno) \\
žádné generování kandidátů, výrazně rychlejší než Apriori & prořezávat lze jen na úrovni
  jednotlivých položek; support lze spočítat teprve po vložení \emph{všech} dat do stromu \\
\bottomrule
\end{tabular}
\end{center}

\mimo{} Doplňkově: \emph{uzavřená} frekventovaná množina (\emph{closed}) je taková, jejíž
žádná vlastní nadmnožina nemá stejný support; \emph{maximální} (\emph{maximal}) je taková,
jejíž žádná vlastní nadmnožina není frekventovaná. Maximální množiny jsou nejkompaktnější
reprezentace (ztrácejí ale informaci o~supportech podmnožin), uzavřené jsou kompromis ---
bezeztrátová reprezentace všech supportů.

\subsection{Přístupy založené na principu učení s~učitelem}
\label{ssec:supervised}

\subsubsection{Formulace úlohy}

\begin{defbox}[Klasifikace a~predikce]
Je dána trénovací množina $\{(x_1,y_1),\dots,(x_p,y_p)\}$, kde $x_i$ je vektor hodnot
vstupních atributů a~$y_i$ hodnota \emph{cílového} atributu. Úkolem je najít model
$\hat y=f(x)$, který dobře předpovídá cílový atribut i~pro dosud neviděné vstupy.

Je-li cílový atribut \emph{diskrétní}, jde o~\textbf{klasifikaci}, je-li \emph{spojitý},
o~\textbf{regresi} (predikci). Řešení má vždy dvě fáze: \emph{(1)}~\textbf{indukce}
(učení, trénování) modelu z~trénovacích dat; \emph{(2)}~\textbf{aplikace} modelu na nová
data.
\end{defbox}

Zásadní je, že model se hodnotí na \emph{nových} datech --- schopnost model naučit trénovací
data nazpaměť je bezcenná (\emph{přeučení}, \emph{overfitting}). Metody odhadu chyby
a~míry kvality jsou v~kap.~\ref{sec:reprez}, odd.~\ref{ssec:eval}.

\subsubsection{Rozhodovací stromy}

\begin{defbox}[Rozhodovací strom]
Nechť $D$ je databáze $\{d_1,\dots,d_p\}$, $\{A_1,\dots,A_n\}$ množina atributů
a~$C=\{C_1,\dots,C_m\}$ množina tříd. \emph{Rozhodovací strom} pro $D$ je strom, kde:
\begin{itemize}
\item každý vnitřní uzel je označen atributem $A_a$;
\item každá hrana je označena predikátem (aplikovatelným na atribut rodičovského uzlu);
\item každý list má označení třídy $C_j$.
\end{itemize}
Strom dělí prostor příznaků na \emph{obdélníkové oblasti}; vzory se klasifikují podle
oblasti, do níž patří.
\end{defbox}

\begin{center}
\small
\begin{tabular}{@{}p{7.4cm}p{8.2cm}@{}}
\toprule
\textbf{Výhody} & \textbf{Nevýhody} \\
\midrule
snadná implementace a~interpretace & obtížné zpracování spojitých dat (nutná kategorizace
  $\to$ obdélníkové oblasti) \\
efektivnost & nelze vždy použít --- hranice tříd může být diagonální a~strom s~osově
  zarovnanými řezy ji umí jen schodovitě aproximovat (v~příkladu bankovního úvěru vede
  hranice napříč rovinou \emph{zůstatek na kontě} $\times$ \emph{příjem}) \\
extrakce jednoduchých pravidel & obtížné zpracování chybějících dat \\
použitelnost i~na velké datové soubory & přeučení $\Rightarrow$ nutné \emph{prořezávání} \\
  & neuvažují vzájemné korelace mezi atributy \\
\bottomrule
\end{tabular}
\end{center}

\textbf{Terminologie:} \emph{rozhodovací atributy} pro dělení datové množiny označují
rozhodovací uzly stromu; \emph{rozhodovací predikáty} označují hrany; \emph{zastavovací
kritérium} je např.\ \uv{všechny vzory z~redukované množiny patří do téže třídy}.

\begin{defbox}[Algoritmus TDIDT (\emph{Top Down Induction of Decision Trees})]
Indukce stromu shora dolů metodou \emph{rozděl a~panuj}:
\begin{enumerate}
\item Vyber jeden atribut jako \emph{kořen} částečného (pod)stromu.
\item Rozděl data přicházející do tohoto uzlu na podmnožiny podle hodnot vybraného atributu
      a~přidej uzel pro každou podmnožinu.
\item Existuje-li uzel takový, že vzory do něj přicházející nepatří všechny do téže třídy,
      zopakuj proces pro něj od kroku 1; jinak skonči.
\end{enumerate}
\end{defbox}

\begin{defbox}[Entropie a~volba dělicího atributu]
Cílem je vybrat takový atribut, který nejlépe rozdělí vzory z~různých tříd. \emph{Entropie}
jako míra neuspořádanosti (\uv{překvapení}) v~systému:
\[ H=-\sum_{j=1}^{m}p_j\log_2 p_j, \]
kde $p_j$ je pravděpodobnost výskytu třídy $j$ (relativní frekvence na uvažované množině
vzorů) a~$m$ počet tříd. Konvence: $0\log_2 0=0$; přepočet
$\log_2 x=\log_{10}x/0{,}301$.

Pro atribut $A$: pro každou jeho možnou hodnotu $v$ se spočítá entropie $H(A(v))$ na množině
příkladů pokrytých kategorií $A(v)$,
\[ H\big(A(v)\big)=-\sum_{j=1}^{m}p_j^{A(v)}\log_2 p_j^{A(v)}, \]
a~\emph{vážená průměrná entropie} atributu jako vážený součet
\[ H(A)=\sum_{v}\frac{|A(v)|}{|D|}\,H\big(A(v)\big). \]
Vybírá se atribut s~\emph{minimální} $H(A)$.
\end{defbox}

\begin{defbox}[Příklad --- žádost o~úvěr]
\begin{center}
\small
\begin{tabular}{@{}lllllc@{}}
\toprule
Klient & Příjem & Konto & Pohlaví & Nezaměstnaný & Úvěr \\
\midrule
C1 & vysoký & vysoké & žena & ne & ano \\
C2 & vysoký & vysoké & muž & ne & ano \\
C3 & nízký & nízké & muž & ne & ne \\
C4 & nízký & vysoké & žena & ano & ano \\
C5 & nízký & vysoké & muž & ano & ano \\
C6 & nízký & nízké & žena & ano & ne \\
C7 & vysoký & nízké & muž & ne & ano \\
C8 & vysoký & nízké & žena & ano & ano \\
C9 & nízký & střední & muž & ano & ne \\
C10 & vysoký & střední & žena & ne & ano \\
C11 & nízký & střední & žena & ano & ne \\
C12 & nízký & střední & muž & ne & ano \\
\bottomrule
\end{tabular}
\end{center}

\textbf{Krok 1 --- výběr atributu podle minimální entropie.} Čtyřpolní tabulka pro
\emph{příjem} a~\emph{úvěr}: vysoký příjem 5 ano / 0 ne, nízký příjem 3 ano / 4 ne.
\begin{align*}
H(\text{příjem}(\text{vysoký}))&=-\tfrac55\log_2\tfrac55-\tfrac05\log_2\tfrac05=0,\\
H(\text{příjem}(\text{nízký}))&=-\tfrac37\log_2\tfrac37-\tfrac47\log_2\tfrac47=0{,}9852,\\
H(\text{příjem})&=\tfrac{5}{12}\cdot 0+\tfrac{7}{12}\cdot 0{,}9852=0{,}5747.
\end{align*}
Obdobně $H(\text{konto})=\tfrac13\cdot 0+\tfrac13\cdot 1+\tfrac13\cdot 1=0{,}6667$,
$H(\text{pohlaví})=\tfrac12\cdot 0{,}9183+\tfrac12\cdot 0{,}9183=0{,}9183$
a~$H(\text{nezaměstnaný})=\tfrac12\cdot 1+\tfrac12\cdot 0{,}65=0{,}825$.
Minimum má \textbf{příjem} $\Rightarrow$ první větvení.

\textbf{Krok 2.} Příjem(vysoký) $\to$ list úvěr(ano). Příjem(nízký), 7 klientů:
$H(\text{konto})=0{,}3935$, $H(\text{pohlaví})=0{,}965$, $H(\text{nezam.})=0{,}9792$
$\Rightarrow$ \textbf{konto}. \textbf{Krok 3.} Konto(vysoké) $\to$ ano, konto(nízké) $\to$ ne;
konto(střední), 3 klienti: $H(\text{pohlaví})=0{,}6667$, $H(\text{nezam.})=0$ $\Rightarrow$
\textbf{nezaměstnanost}.

\textbf{Pravidla ze stromu} (cesta kořen--list $=$ pravidlo): příjem(vysoký) $\to$ ano;
příjem(nízký) $\wedge$ konto(vysoké) $\to$ ano; příjem(nízký) $\wedge$ konto(nízké) $\to$ ne;
příjem(nízký) $\wedge$ konto(střední) $\wedge$ nezam.(ano) $\to$ ne; příjem(nízký) $\wedge$
konto(střední) $\wedge$ nezam.(ne) $\to$ ano.
\end{defbox}

\textbf{Převod stromu na pravidla} obecně: každá cesta od kořene k~listu odpovídá jednomu
pravidlu; vnitřní uzly (s~hodnotou pro odpovídající hranu) tvoří levou stranu (tělo),
list (cílová třída) pravou stranu (hlavu).

\textbf{Faktory důležité pro indukci stromu:}
\begin{itemize}
\item \emph{volba} rozhodovacích atributů --- ovlivňuje efektivitu; užitečný je informovaný
      vstup doménového experta;
\item \emph{pořadí} atributů --- eliminuje redundantní porovnání;
\item \emph{větvení} --- počet nutných dělení; náročnější pro spojité atributy či mnoho hodnot;
\item \emph{forma stromu} --- vhodnější vyvážené s~málo úrovněmi (za cenu složitějších
      porovnání); některé algoritmy tvoří jen binární stromy (často hluboké);
\item \emph{zastavovací kritéria} --- správná klasifikace trénovacích dat; \emph{předčasné
      zastavení} zabrání velkým stromům a~přeučení (přesnost vs.\ efektivita);
      \emph{Occamova břitva} (1320) --- preferuj nejjednodušší hypotézu pro $D$;
\item \emph{množství trénovacích dat} --- málo $\Rightarrow$ nespolehlivý strom, mnoho
      $\Rightarrow$ přeučení;
\item \emph{prořezávání} (\emph{pruning}) --- vyšší přesnost i~efektivita na testovacích datech;
      odstraní redundantní porovnání, podstromy a~irelevantní atributy; lze už během indukce,
      nebo až na hotovém stromu;
\item \emph{složitost} --- závisí na $p$ (data), $n$ (atributy) a~formě stromu: indukce
      $O(n\,p\log p)$, klasifikace $q$ vzorů stromem hloubky $O(\log p)$ je $O(q\log p)$.
\end{itemize}

\textbf{Algoritmus ID3} (\emph{Iterative Dichotomiser}, Ross Quinlan, 1986): učení funkcí
s~booleovskými hodnotami, \emph{hladová} metoda, staví strom shora dolů; v~každém uzlu volí
atribut, který nejlépe klasifikuje lokální trénovací data --- nejlepší atribut má
\emph{největší informační zisk}. Proces pokračuje, dokud nejsou všechna trénovací data
správně klasifikována, nebo dokud nejsou využity všechny atributy.

\begin{defbox}[Algoritmus ID3 a~informační zisk]
\textsc{ID3}(\emph{EXAMPLES}, \emph{TARGET\_ATTRIBUTE}, \emph{ATTRIBUTES}):
\begin{algorithmic}[1]
\State vytvoř kořen $\mathit{ROOT}$
\If{všechny \emph{EXAMPLES} jsou pozitivní} \Return $\mathit{ROOT}$ s~jediným uzlem
  označeným \uv{$+$} \EndIf
\If{všechny \emph{EXAMPLES} jsou negativní} \Return $\mathit{ROOT}$ s~jediným uzlem
  označeným \uv{$-$} \EndIf
\If{$\mathit{ATTRIBUTES}=\emptyset$} \Return $\mathit{ROOT}$ s~uzlem označeným
  nejčastější hodnotou \emph{TARGET\_ATTRIBUTE} v~\emph{EXAMPLES} \EndIf
\State $A\gets$ atribut z~\emph{ATTRIBUTES}, který nejlépe klasifikuje \emph{EXAMPLES}
\State rozhodovací atribut pro $\mathit{ROOT}\gets A$
\For{každou možnou hodnotu $v_i$ atributu $A$}
   \State připoj k~$\mathit{ROOT}$ novou hranu odpovídající predikátu $A=v_i$
   \State $\mathit{EXAMPLES}_{v_i}\gets$ podmnožina \emph{EXAMPLES} s~hodnotou $v_i$ pro $A$
   \If{$\mathit{EXAMPLES}_{v_i}=\emptyset$}
      \State připoj k~hraně list označený nejčastější hodnotou \emph{TARGET\_ATTRIBUTE}
   \Else
      \State připoj podstrom \textsc{ID3}$(\mathit{EXAMPLES}_{v_i},
             \mathit{TARGET\_ATTRIBUTE}, \mathit{ATTRIBUTES}\setminus\{A\})$
   \EndIf
\EndFor
\State \Return $\mathit{ROOT}$
\end{algorithmic}

\textbf{Informační zisk} (\emph{information gain}) --- očekávané snížení entropie po rozdělení
dat podle uvažovaného atributu:
\[ \mathrm{GAIN}(D,A)=H(D)-\sum_{v\in \mathrm{VALUES}(A)}\frac{|D_v|}{|D|}H(D_v). \]
Protože $H(D)$ je pro všechny atributy stejná, je maximalizace informačního zisku ekvivalentní
minimalizaci vážené entropie $H(A)$.
\end{defbox}

\textbf{Algoritmus C4.5} (Quinlan, 1993) --- modifikace ID3:
\begin{itemize}
\item \textbf{Chybějící data} --- při konstrukci ignorována (kritérium dělení se počítá jen ze
      známých hodnot), při klasifikaci odhadována z~hodnot známých u~jiných vzorů.
\item \textbf{Spojitá data} --- kategorizují se podle hodnot atributu na trénovacích vzorech.
\item \textbf{Prořezávání validací} (trénovací $\tfrac23$ / validační $\tfrac13$).
      \emph{Reduced-error pruning}: podstrom se nahradí listem s~nejčastější třídou, ale jen
      když nový strom není na validační množině horší. \emph{Subtree raising}: podstrom se
      nahradí svým nejčastěji používaným podstromem zvednutým výš; testuje se přesnost.
\item \textbf{Prořezávání pravidel} (\emph{rule post-pruning}) --- strom (přeučení dovoleno)
      $\to$ pravidla (jedno na cestu kořen--list) $\to$ vypouštěj předpoklady, dokud to
      nezhorší odhadovanou přesnost $\to$ seřaď podle očekávané přesnosti a~v~tom pořadí
      klasifikuj. Snazší a~radikálnější než prořezávání stromů, výsledek je průhledný.
\item \textbf{Odhad přesnosti pravidel} --- standardně validační množinou; C4.5 hodnotí
      \emph{na trénovacích datech} $\mathit{acc}=k/n$ (správně klasifikované ku pokrytým).
      Charakteristikou pravidla je \emph{dolní mez} z~binomického rozdělení, na 95\% hladině
      \[ \mathit{acc}-1{,}96\sqrt{\frac{\mathit{acc}(1-\mathit{acc})}{n}}. \]
\item \textbf{Poměrný informační zisk} (\emph{information gain ratio}) jako kritérium dělení
      (bere se \emph{největší}):
      \[ \mathrm{GAIN\_RATIO}(D,A)=\frac{\mathrm{GAIN}(D,A)}
         {-\sum_{v}\frac{|D_v|}{|D|}\log_2\frac{|D_v|}{|D|}}, \]
      jmenovatel penalizuje \emph{nadměrné větvení} (atributy s~mnoha hodnotami).
\end{itemize}

\textbf{Algoritmus C5.0} --- komerční C4.5 pro velké databáze: efektivnější generování pravidel
a~vyšší přesnost díky \emph{boostingu} (vzorům se přiřadí váhy podle vlivu při klasifikaci, pro
každou kombinaci vah vznikne samostatný klasifikátor, hlasuje se většinově).

\begin{defbox}[CART --- \emph{Classification And Regression Trees}]
Generuje \emph{binární} stromy --- při každém dělení dat se tvoří právě dvě hrany. Volba
nejlepšího rozhodovacího atributu podle entropie nebo \textbf{Gini indexu}: pro atribut
o~$r$ hodnotách $v_1,\dots,v_r$, $n$ vzorů a~$m$ tříd
\[ G=\sum_{i=1}^{r}\frac{n_i}{n}G(v_i),\qquad G(v_i)=1-\sum_{j=1}^{m}p_j^2,
   \qquad \sum_{i=1}^{r}n_i=n, \]
kde $n_i$ je počet vzorů s~hodnotou $v_i$ a~$p_j$ pravděpodobnost třídy $j$ v~těchto $n_i$
vzorech. \emph{Nižší} hodnoty $G$ znamenají lepší diskriminaci.

Dělení se provádí podle \uv{nejlepšího} nalezeného bodu pro daný uzel $t$; procházejí se
všechny možné hodnoty $v$ rozhodovacího atributu $A$ (při rovnosti se použije pravý
podstrom). Označme $P_L,P_R$ pravděpodobnosti zpracování levým, resp.\ pravým podstromem
a~$P(C_j\mid t_L),P(C_j\mid t_R)$ pravděpodobnosti, že vzor patří do třídy $C_j$ a~je v~levém,
resp.\ pravém podstromu uzlu $t$. Kritérium dělení:
\[ \Phi=2P_LP_R\sum_{j=1}^{m}\big|P(C_j\mid t_L)-P(C_j\mid t_R)\big|, \]
které nabývá maxima, jsou-li všechny vzory dané třídy seskupeny v~témže podstromu.

\textbf{Charakteristické vlastnosti CART:} pořadí atributů odpovídá jejich vlivu při
klasifikaci; chybějící data jsou ignorována a~nezapočítávají se do vyhodnocení rozhodovacích
kritérií; zastavovací kritérium --- žádné další dělení už nezlepší klasifikační přesnost;
vysoká přesnost na trénovací množině nemusí odrážet přesnost na testovacích datech.
\end{defbox}

\begin{defbox}[CHAID --- \emph{Chi-square Automatic Interaction Detection}]
Kritériem větvení je $\chi^2$. Algoritmus automaticky \emph{seskupuje hodnoty} kategoriálních
atributů --- hodnoty se postupně slučují od původního počtu, dokud nezůstanou jen dvě grupy;
následně se vybere atribut a~jeho kategorizace, které jsou pro dělení dat v~daném kroku
nejlepší. Při větvení tedy nevzniká tolik větví, kolik má atribut hodnot.

\textbf{Seskupování hodnot atributu:}
\begin{enumerate}
\item Opakuj, dokud nezůstanou jen dvě grupy hodnot:
      \emph{(a)}~vyber pár kategorií atributu, které jsou si z~hlediska $\chi^2$ nejvíce
      podobné a~lze je sloučit; \emph{(b)}~novou kategorii uvažuj jako možné seskupení
      v~daném kroku.
\item Pro každou možnou grupu hodnot testuj $\chi^2$ testem její relevanci k~predikované
      proměnné.
\item Vyber seskupení, které dává \uv{nejlepší} rozdělení s~nejvýznamnějším rozštěpením.
\item Urči, zda toto seskupení statisticky významně přispívá k~rozlišení vzorů z~různých tříd.
\end{enumerate}
\end{defbox}

\subsubsection{Bayesovská klasifikace}

\begin{defbox}[Bayesova formule a~hypotéza MAP]
Pravděpodobnost, že hypotéza $H$ platí, pozorujeme-li $E$:
\[ P(H\mid E)=\frac{P(E\mid H)\,P(H)}{P(E)}, \]
kde $P(H)$ je \emph{apriorní} pravděpodobnost hypotézy, $P(H\mid E)$ \emph{aposteriorní}
(podmíněná) pravděpodobnost a~$P(E)$ pravděpodobnost jevu $E$.

\emph{Nejpravděpodobnější hypotéza} $H_{MAP}$ (\emph{maximum a~posteriori}) má největší
aposteriorní pravděpodobnost. Protože jmenovatel $P(E)$ je pro všechny hypotézy stejný, lze
jej zanedbat:
\[ H_{MAP}=\arg\max_{H\in\mathcal H} P(H\mid E)=\arg\max_{H\in\mathcal H} P(E\mid H)\,P(H). \]
\end{defbox}

\begin{defbox}[Naivní bayesovský klasifikátor]
\textbf{Předpoklad:} jevy $E_1,\dots,E_k$ jsou \emph{podmíněně nezávislé} při platnosti
hypotézy $H$. V~reálných úlohách je tento předpoklad splněn jen zřídka --- odtud označení
\uv{naivní}. Pak
\[ P(H\mid E_1,\dots,E_k)=\frac{P(E_1,\dots,E_k\mid H)P(H)}{P(E_1,\dots,E_k)}
   =\frac{P(H)}{P(E_1,\dots,E_k)}\prod_{i=1}^{k}P(E_i\mid H), \]
a~klasifikujeme hypotézou s~největší aposteriorní pravděpodobností
\[ H_{MAP}=\arg\max_{H_t} P(H_t)\prod_{i=1}^{k}P(E_i\mid H_t). \]
Odhady z~dat: $P(H_t)=P(\text{třída }t)=\frac{\text{počet vzorů z~}t}{\text{počet všech
vzorů}}$ a~$P(E_k\mid H_t)=\frac{\text{počet vzorů z~}t\text{ splňujících }E_k}
{\text{počet vzorů z~}t}$.

\textbf{Výhoda:} možnost klasifikovat i~\emph{nekompletně popsané} vzory.
\textbf{Nevýhoda:} nulová aposteriorní pravděpodobnost $P(E_k\mid H_t)$, není-li
v~trénovací množině žádný odpovídající vzor, resp.\ podhodnocené členy při nízkém společném
výskytu $E_k$ a~$H_t$.
\end{defbox}

\begin{defbox}[Příklad --- bayesovská klasifikace úvěru]
\emph{Jednoduchý případ:} $P(\text{půjčit})=0{,}667$, $P(\text{zamítnout})=0{,}333$,
$P(\text{vysoký příjem}\mid\text{půjčit})=0{,}91$,
$P(\text{vysoký příjem}\mid\text{zamítnout})=0{,}12$. Pak
$0{,}91\cdot 0{,}667=0{,}607$ vs.\ $0{,}12\cdot 0{,}333=0{,}040$, tedy
$H_{MAP}=\text{půjčit}$.

\emph{Naivní klasifikátor} nad daty z~příkladu o~úvěru (12 klientů):
$P(\text{úvěr(ano)})=8/12=0{,}667$, $P(\text{úvěr(ne)})=4/12=0{,}333$;
$P(\text{konto(střední)}\mid\text{ano})=2/8=0{,}25$,
$P(\text{konto(střední)}\mid\text{ne})=2/4=0{,}5$;
$P(\text{nezam.(ne)}\mid\text{ano})=5/8=0{,}625$,
$P(\text{nezam.(ne)}\mid\text{ne})=1/4=0{,}25$.

Pro klienta se středním kontem, který není nezaměstnaný:
\[ 0{,}667\cdot 0{,}25\cdot 0{,}625=0{,}1042 \quad\text{vs.}\quad
   0{,}333\cdot 0{,}5\cdot 0{,}25=0{,}0416, \]
klient se tedy zařadí do třídy \textbf{úvěr(ano)}.
\end{defbox}

\subsubsection{Support Vector Machines}

SVM (V.~Vapnik, 70.~léta v~Rusku; na Západě známé od roku 1992) jsou \emph{lineární}
klasifikátory hledající hyperrovinu, která oddělí pozitivní a~negativní třídu; pro nelineární
separaci se používají \emph{kernelové} funkce. Mají důkladný teoretický základ i~vyšší přesnost
než většina metod, zejména pro vysokorozměrná data --- patří k~nejlepším klasifikátorům textu.

\begin{defbox}[Lineární SVM, separabilní případ]
Trénovací množina $D=\{(x_1,y_1),\dots,(x_r,y_r)\}$, $x_i\in X\subseteq\mathbb R^n$,
$y_i\in\{1,-1\}$. SVM hledá lineární funkci $f(x)=\langle w\cdot x\rangle+b$ tak, že
\[ y_i=\begin{cases}1 & \text{pro }\langle w\cdot x_i\rangle+b\ge 0,\\
-1 & \text{pro }\langle w\cdot x_i\rangle+b<0.\end{cases} \]
Hyperrovina $\langle w\cdot x\rangle+b=0$ se nazývá \emph{rozhodovací hranice}
(\emph{decision boundary}). Takových hyperrovin je mnoho --- SVM hledá tu s~\emph{největším
okrajem} (\emph{margin}); teorie strojového učení říká, že tato hyperrovina minimalizuje
mez chyby.

Vzdálenost bodu $x_i$ od hyperroviny je $\dfrac{|\langle w\cdot x_i\rangle+b|}{\|w\|}$, kde
$\|w\|=\sqrt{\langle w\cdot w\rangle}$. Pro okrajové hyperroviny $H_+$ a~$H_-$ platí
\[ d_+=d_-=\frac{1}{\|w\|},\qquad \text{margin}=d_++d_-=\frac{2}{\|w\|}. \]

\textbf{Optimalizační úloha:}
\[ \text{minimalizuj }\ \frac{\langle w\cdot w\rangle}{2}
   \quad\text{za podmínek}\quad y_i\big(\langle w\cdot x_i\rangle+b\big)\ge 1,\
   i=1,\dots,r, \]
což je souhrnný zápis podmínek $\langle w\cdot x_i\rangle+b\ge 1$ pro $y_i=1$
a~$\langle w\cdot x_i\rangle+b\le -1$ pro $y_i=-1$.
\end{defbox}

\begin{defbox}[Duální formulace a~podpůrné vektory]
Úloha se řeší standardní \emph{Lagrangeovou} metodou:
\[ L_P=\frac{\langle w\cdot w\rangle}{2}-\sum_{i=1}^{r}\alpha_i
   \big[y_i(\langle w\cdot x_i\rangle+b)-1\big],\qquad \alpha_i\ge 0. \]
Optimální řešení musí splňovat \emph{Kuhn--Tuckerovy podmínky}, které jsou pro \emph{konvexní}
účelovou funkci s~\emph{lineárními} omezeními nejen nutné, ale i~postačující. Podmínka
komplementarity ukazuje, že $\alpha_i>0$ mohou mít jen body \emph{na} okrajových
hyperrovinách (pro ně $y_i(\langle w\cdot x_i\rangle+b)-1=0$) --- to jsou
\textbf{podpůrné vektory} (\emph{support vectors}); pro všechny ostatní je $\alpha_i=0$.

Vynulováním parciálních derivací $L_P$ podle $w$ a~$b$ a~substitucí zpět dostaneme
\emph{Wolfeho duální} úlohu
\[ L_D=\sum_{i=1}^{r}\alpha_i-\frac12\sum_{i,j}y_iy_j\alpha_i\alpha_j
   \langle x_i\cdot x_j\rangle\ \longrightarrow\ \max, \]
jejíž maximum se nabývá pro tytéž hodnoty $w,b,\alpha_i$ jako minimum $L_P$. Duální úloha se
řeší snáz než primární (numerické techniky).

\textbf{Rozhodovací hranice} a~\textbf{testování} vzoru $z$:
\[ \langle w\cdot x\rangle+b=\sum_{i\in \mathit{sv}}y_i\alpha_i\langle x_i\cdot x\rangle+b=0,
\qquad
\mathrm{sign}\Big(\sum_{i\in \mathit{sv}}\alpha_iy_i\langle x_i\cdot z\rangle+b\Big). \]
Vrátí-li výraz 1, je $z$ klasifikován jako pozitivní, jinak jako negativní.
\end{defbox}

\begin{defbox}[Kernelový trik]
Reálná data mohou vyžadovat nelineární separaci, proto se \emph{transformují} do jiného (obvykle
mnohem vícerozměrného) \emph{prostoru příznaků} (\emph{feature space}) $F$, v~němž třídy oddělí
\emph{lineární} hranice: $\phi:X\to F$.

\textbf{Problém explicitní transformace:} dimenze $F$ může být pro užitečné transformace enormní
i~při rozumném počtu vstupních atributů --- \emph{prokletí dimenzionality} činí explicitní
výpočet nezvládnutelným.

\textbf{Klíčové pozorování:} duál potřebuje ke konstrukci optimální hyperroviny i~k~vyhodnocení
rozhodovací funkce \emph{pouze skalární součiny} $\langle\phi(x)\cdot\phi(z)\rangle$, nikdy
$\phi(x)$ samotné. Umíme-li tedy počítat skalární součin přímo ze vstupních vektorů, $\phi$ znát
nemusíme --- \emph{kernelová funkce}
\[ K(x,z)=\langle\phi(x)\cdot\phi(z)\rangle \]
nahradí skalární součin ve všech výrazech (\textbf{kernelový trik}).

\emph{Příklad (polynomiální kernel $d=2$, 2-D vstup):}
\begin{align*}
\langle x\cdot z\rangle^2&=(x_1z_1+x_2z_2)^2=x_1^2z_1^2+2x_1z_1x_2z_2+x_2^2z_2^2\\
&=\big\langle (x_1^2,x_2^2,\sqrt2 x_1x_2)\cdot(z_1^2,z_2^2,\sqrt2 z_1z_2)\big\rangle
 =\langle\phi(x)\cdot\phi(z)\rangle.
\end{align*}
Kernel $\langle x\cdot z\rangle^d$ je tedy skalární součin v~transformovaném prostoru, aniž
bychom $\phi$ kdy vyčíslili. Zda je daná funkce kernelem, říká \emph{Mercerova věta}. Běžné
kernely: lineární (identita), polynomiální, gaussovský (RBF), sigmoidální.
\end{defbox}

\textbf{Další vlastnosti SVM.} Pracuje jen v~reálném prostoru (kategoriální atributy je nutné
převést na numerické) a~jen \emph{dvoutřídně} --- pro více tříd \emph{jeden proti ostatním}
(\emph{one-against-rest}) nebo kódy opravující chyby. Hyperrovina je pro člověka obtížně
pochopitelná a~kernely to zhoršují, takže se SVM používá tam, kde se srozumitelnost
nevyžaduje. \mimo{} Pro nelineárně separabilní data se zavedou \emph{proměnné vůle}
(\emph{slack variables}) $\xi_i\ge 0$ a~minimalizuje se
$\frac{\langle w\cdot w\rangle}{2}+C\sum_i\xi_i$; $C$ řídí kompromis mezi šířkou okraje
a~počtem chyb.

\subsubsection{Logistická regrese}
\label{ssec:logreg}

Logistická regrese je nejdůležitější případ nelineární regrese (odd.~\ref{sec:priprava}) a~zároveň
klasifikátor: závislá proměnná $y$ je kategoriální, typicky dvouhodnotová, a~modeluje se
\emph{pravděpodobnost}, že $y$ nabude určité hodnoty, v~závislosti na kombinaci hodnot
nezávislých veličin.

\begin{defbox}[Logistická regrese]
Pro binární $y\in\{0,1\}$ se modeluje logaritmus \uv{podmíněné šance}
(\emph{conditional odds}):
\[ \ln\frac{P(y=1\mid x_1,\dots,x_m)}{1-P(y=1\mid x_1,\dots,x_m)}
   =\beta_0+\beta_1x_1+\cdots+\beta_mx_m=z, \]
odkud
\[ P(y=1\mid x)=\frac{e^{z}}{1+e^{z}}=\frac{1}{1+e^{-z}}=\sigma(z)
   \qquad\text{(sigmoida).} \]

\textbf{Učení} maximalizací věrohodnosti (\emph{maximum likelihood}):
\[ L=\prod_{i=1}^{p}P\big(y_i\mid x_{i,1},\dots,x_{i,m}\big)
   =\prod_{i=1}^{p}\Big(\frac{1}{1+e^{-z_i}}\Big)^{y_i}
     \Big(\frac{1}{1+e^{z_i}}\Big)^{1-y_i}. \]
Místo maximalizace $L$ se minimalizuje \emph{křížová entropie} $E=-\log L$:
\[ E=-\sum_i\Big[y_i\log\sigma(z_i)+(1-y_i)\log\big(1-\sigma(z_i)\big)\Big]. \]
Protože derivace sigmoidy je $\sigma'=\sigma(1-\sigma)$, dostaneme po zjednodušení velmi
jednoduchá adaptační pravidla
\[ \Delta\beta_j=-\eta\frac{\partial E}{\partial\beta_j}
   =\eta\sum_i\big(y_i-\sigma(z_i)\big)x_{i,j},\qquad
   \Delta\beta_0=\eta\sum_i\big(y_i-\sigma(z_i)\big), \]
kde $\eta$ je rychlost učení. Chyba tedy vstupuje do úpravy vah přímo jako rozdíl
\emph{pozorované} a~\emph{predikované} hodnoty.
\end{defbox}

Je běžnou volbou základního klasifikátoru tam, kde se požaduje pravděpodobnostní výstup
a~interpretovatelné koeficienty --- např.\ predikce vazeb v~sítích (kap.~\ref{sec:sna}).

\subsubsection{Diskriminační analýza}
\label{ssec:da}

\begin{defbox}[Diskriminační analýza]
\emph{Klasifikace vzorů do známých tříd}: hledá se závislost nominální veličiny (určující
příslušnost ke třídě) na ostatních numerických veličinách. Předpokládáme, že ke každé třídě
$C_j$, $j=1,\dots,m$ existuje \emph{diskriminační funkce} $g_j$; vzor
$x=(x_1,\dots,x_m)$ patří do třídy $C_j$, právě když
\[ g_j(x)=\max_{k} g_k(x). \]

\textbf{Lineární diskriminační analýza:} $g_j(x)=\beta_{j0}+\beta_{j1}x_1+\cdots+\beta_{jm}x_m$.
Pro dvě třídy lze místo $g_1,g_2$ hledat jedinou funkci $g=g_1-g_2$ a~klasifikovat podle
\emph{znaménka}: \uv{$+$} znamená třídu 1, \uv{$-$} třídu 2.

\textbf{Optimální klasifikace ve smyslu minimální chyby} nastane, jsou-li diskriminační
funkce $\propto$ aposteriorní (podmíněné) pravděpodobnosti zařazení pozorování $x$ do
třídy $C_j$: $g_j(x)=P(C_j\mid x)=\dfrac{p(x\mid C_j)P(C_j)}{p(x)}$; pro dvě třídy
$g(x)=p(x\mid C_1)P(C_1)-p(x\mid C_2)P(C_2)$.
\end{defbox}

Pro \emph{normální rozdělení} s~různými středními hodnotami $\mu_1,\mu_2$ a~obecnými
kovariančními maticemi $S_1,S_2$ vychází \emph{kvadratická} diskriminační funkce. Jsou-li
kovarianční matice \emph{shodné} ($S_1=S_2=S$), redukuje se na \emph{lineární}
\[ g(x)=(\mu_1-\mu_2)^{\!\top}S^{-1}x-\tfrac12(\mu_1-\mu_2)^{\!\top}S^{-1}(\mu_1+\mu_2)
   -\ln\frac{P(C_2)}{P(C_1)}, \]
a~jsou-li navíc obě třídy stejně pravděpodobné a~$S=I$, jde o~jednoduché srovnávání
vzdáleností od středních hodnot. Pro normální rozdělení se tedy hledání diskriminačních
funkcí redukuje na \textbf{odhad středních hodnot} (výběrovými průměry) a~\textbf{kovariančních
matic} (korigovanými výběrovými rozptyly).

\subsubsection{Neuronové sítě a~extrémní učicí stroje (ELM)}

\begin{defbox}[Síť ELM (\emph{Extreme Learning Machine}, G.-B.~Huang)]
Dopředná síť s~\emph{jednou} skrytou vrstvou: $n$ vstupních neuronů, $L$ skrytých, 1 výstupní.
Výstup skrytých neuronů $h(x)=\big[G(a_1,b_1,x),\dots,G(a_L,b_L,x)\big]$, kde $a_i$ je
váhový vektor a~$b_i$ bias $i$-tého neuronu; např.\
$G(a_i,b_i,x)=g(\langle a_i\cdot x\rangle+b_i)$ (sigmoida) nebo
$G(a_i,b_i,x)=g(b_i\|x-a_i\|)$ (RBF). Výstup sítě
$f_L(x)=\sum_{i=1}^{L}\beta_iG(a_i,b_i,x)$.

\textbf{Klíčová myšlenka:} parametry skrytých neuronů \emph{není nutné adaptovat} --- mohou být
vygenerovány \emph{náhodně}, dokonce bez znalosti trénovacích dat. Pro každou spojitou $f$
a~náhodnou posloupnost $\{(a_i,b_i)\}$ platí s~pravděpodobností 1
\[ \lim_{L\to\infty}\Big\|f-\textstyle\sum_{i=1}^{L}\beta_iG(a_i,b_i,\cdot)\Big\|=0 \]
při $\beta_i$ minimalizujících $\|f-f_L\|$.

\textbf{Maticový model:} soustava $\sum_{i=1}^{L}\beta_iG(a_i,b_i,x_j)=t_j$, $j=1,\dots,N$ je
ekvivalentní $H\beta=T$, kde $H$ je \emph{výstupní matice skryté vrstvy} ($j$-tý řádek $=$
výstupy skrytých neuronů pro vzor $x_j$).

\textbf{Algoritmus učení} (tři neiterativní kroky):
\begin{algorithmic}[1]
\State zvol \emph{náhodně} parametry skrytých neuronů $(a_i,b_i)$, $i=1,\dots,L$
\State spočítej výstupní matici skryté vrstvy $H$
\State spočítej výstupní váhy $\beta=H^{+}T$
\end{algorithmic}
kde $H^{+}$ je Mooreova--Penroseova pseudoinverze $H$. Stabilnější řešení s~lepší generalizací
dá regularizace --- přičtení $I/C$ na diagonálu $H^{\!\top}H$:
$f(x)=h(x)\big(I/C+H^{\!\top}H\big)^{-1}H^{\!\top}T$.
\end{defbox}

\textbf{Výhody ELM:} neiterativní procedura ze tří kroků; velmi rychlé trénování --- neřeší se
lokální minima, inicializace ani přeučení; parametry $a_i,b_i$ jsou nezávislé na datech i~na
sobě, takže je lze vygenerovat \emph{už před} předložením dat (konvenční techniky musí data
\uv{vidět}); funguje s~\emph{libovolnou} omezenou po částech spojitou nekonstantní přenosovou
funkcí (gradientní algoritmy vyžadují hladké).

\subsubsection{Ensemble metody}

\begin{defbox}[Bagging (\emph{bootstrapped aggregating})]
Snaží se snížit \emph{rozptyl} klasifikátoru: je-li rozptyl klasifikátoru $\sigma^2$, pak
rozptyl průměru $k$ nezávislých a~stejně rozdělených (i.i.d.) klasifikátorů klesne na
$\sigma^2/k$.

\textbf{Bootstrapping} aproximuje i.i.d.\ klasifikátory: body se vybírají \emph{s~opakováním}
uniformně z~původních dat; vzorek má přibližně stejnou velikost jako původní data, může
obsahovat duplikáty a~obsahuje podíl
\[ 1-\Big(1-\frac1n\Big)^{\!n}\approx 1-\frac1e=63{,}2\,\% \]
různých bodů z~původních dat (pravděpodobnost, že bod ve vzorku \emph{není}, je
$(1-1/n)^n$). Nakreslí se $k$ nezávislých bootstrapových vzorků velikosti $n$, na každém se
natrénuje klasifikátor a~pro testovací vzor se predikce určí \emph{většinovým hlasem}.

\textbf{Vhodnost:} rozhodovací stromy jsou pro bagging ideální --- mají \emph{nízké
zkreslení} a~\emph{vysoký rozptyl} (jsou-li dostatečně hluboké).
\textbf{Omezení:} bagging \emph{nesnižuje zkreslení} modelu; je-li zkreslení hlavním zdrojem
chyby, může přesnost i~mírně zhoršit kvůli korelacím mezi komponentami (a~tedy porušenému
i.i.d.\ předpokladu).
\end{defbox}

\begin{defbox}[Náhodné lesy (\emph{random forests})]
Mají-li $k$ klasifikátorů s~rozptylem $\sigma^2$ vzájemnou \emph{pozitivní korelaci} $\rho$,
je rozptyl zprůměrované predikce
\[ \rho\sigma^2+\frac{(1-\rho)\sigma^2}{k}. \]
Člen $\rho\sigma^2$ je \emph{nezávislý} na počtu komponent $k$ --- právě on omezuje přínos
baggingu. Bootstrapované stromy přitom bývají pozitivně korelované.

\emph{Náhodný les} je ensemble stromů, do jejichž \emph{konstrukce} je vložena náhodnost, aby
byly komponenty méně korelované. Dvě cesty: \emph{(1)}~\textbf{bagging} --- stromy se trénují na
bootstrapových vzorcích; \emph{(2)}~\textbf{omezený počet příznaků} --- v~každém uzlu se předloží
\emph{náhodná podmnožina} $m$ příznaků a~informační zisk se počítá jen na ní. Druhá cesta navíc
zrychluje trénování a~stromy není nutné prořezávat; snížený rozptyl nezhoršuje zkreslení. Výstup
lesa je většinový hlas.

\textbf{Algoritmus:} pro každý z~$N$ stromů vytvoř bootstrapový vzorek a~nad ním trénuj strom
tak, že v~každém uzlu vybereš náhodně $m$ příznaků, spočítáš informační zisk jen na nich
a~vybereš optimální.

\textbf{Parametry:} $m\approx\sqrt{\text{počet příznaků}}$ (lesy na něm nejsou příliš citlivé);
počet stromů --- lze pokračovat, dokud chyba klesá.
\end{defbox}

\begin{defbox}[Boosting a~AdaBoost (Freund \& Schapire, 1996)]
\textbf{Princip:} každý trénovací vzor má \emph{váhu}, klasifikátory se trénují s~těmito vahami
a~váhy se iterativně upravují podle výkonu --- budoucí modely tedy \emph{závisí na předchozích}
(tentýž algoritmus, jinak vážený trénovací soubor). \textbf{Idea:} zvýšit relativní váhu
\emph{nesprávně klasifikovaných} vzorů, a~tak se v~dalších iteracích soustředit na ně; vážená
kombinace takové posloupnosti má nižší \emph{zkreslení} (bagging naopak snižuje rozptyl).

\textsc{AdaBoost}(data $D$, bázový klasifikátor $A$, max.\ počet kol $T$):
\begin{algorithmic}[1]
\State $t\gets 0$; pro každý vzor $i$ inicializuj $W_1(i)=1/n$
\Repeat
  \State $t\gets t+1$
  \State urči váženou chybu $\epsilon_t$ na $D$, aplikuje-li se $A$ na data s~vahami
         $W_t(\cdot)$
  \State $\alpha_t\gets \tfrac12\ln\big((1-\epsilon_t)/\epsilon_t\big)$
  \For{každý \emph{nesprávně} klasifikovaný $X_i\in D$}
     \State $W_{t+1}(i)\gets W_t(i)e^{\alpha_t}$
  \EndFor
  \For{každý \emph{správně} klasifikovaný $X_i\in D$}
     \State $W_{t+1}(i)\gets W_t(i)e^{-\alpha_t}$
  \EndFor
  \State normalizuj $W_{t+1}(i)\gets W_{t+1}(i)/\sum_j W_{t+1}(j)$
\Until{$t\ge T$ \textbf{nebo} $\epsilon_t=0$ \textbf{nebo} $\epsilon_t\ge 0{,}5$}
\State klasifikuj testovací vzory ensemblem s~vahami komponent $\alpha_t$
\end{algorithmic}
$\epsilon_t$ je podíl nesprávně predikovaných trénovacích vzorů (po vážení $W_t$) modelem
v~$t$-té iteraci.

\textbf{Vlastnosti:} \emph{posloupnost} klasifikátorů téhož bázového algoritmu, každý závisí na
předchozím a~soustředí se na jeho chyby; při testování se výsledky kombinují. \textbf{Citlivost
na šum:} při mnoha odlehlých hodnotách může důraz na \uv{obtížné} příklady výkon zhoršit.
\end{defbox}

\begin{center}
\small
\begin{tabular}{@{}p{7.6cm}p{8.0cm}@{}}
\toprule
\textbf{Náhodné lesy} & \textbf{Boosting} \\
\midrule
rychlé --- mohou běžet \emph{paralelně} a~v~každém uzlu prohledávají malou množinu příznaků
  & vyčerpávající --- v~každé fázi prohledává \emph{celou} množinu příznaků a~každá fáze
  závisí na předchozí \\
stromy jsou levné na trénování, lze jich ve stejném čase vypěstovat více i~na velkých
  a~složitých datech & musí běžet \emph{sekvenčně}, což může být nákladné \\
překvapivě dobré výsledky vzhledem k~malé části prostoru problému, kterou každý strom vidí
  & při \emph{stejném počtu} stromů boosting náhodné lesy překonává (lesy v~každé fázi
  prohledávají jen podmnožinu) \\
\bottomrule
\end{tabular}
\end{center}

\subsubsection{Srovnání klasifikačních metod}

\begin{center}
\small
\begin{tabular}{@{}p{2.9cm}p{2.3cm}p{2.2cm}p{2.3cm}p{4.6cm}@{}}
\toprule
\textbf{Metoda} & \textbf{Srozumitel\-nost} & \textbf{Přesnost} & \textbf{Rychlost}
  & \textbf{Poznámka} \\
\midrule
Rozhodovací strom & vysoká & střední & vysoká & obdélníkové oblasti, nutné prořezávání,
  ignoruje korelace atributů \\
Pravidla z~C4.5 & nejvyšší & střední & vysoká & prořezávání pravidel je radikálnější než
  stromů \\
Naivní Bayes & střední & překvapivě dobrá & velmi vysoká & zvládne nekompletní vzory;
  problém s~nulovými pravděpodobnostmi \\
SVM & velmi nízká & vysoká & nízká (trénink) & nejlepší pro vysokorozměrná data (text);
  jen 2 třídy, jen reálný prostor \\
Logistická regrese & vysoká & střední & vysoká & pravděpodobnostní výstup,
  interpretovatelné koeficienty \\
Diskriminační analýza & střední & dobrá při normalitě & vysoká & stačí odhadnout
  $\mu_j$ a~$S_j$ \\
Neuronová síť / ELM & velmi nízká & vysoká & ELM velmi vysoká & ELM negeneruje lokální
  minima, netrénuje skrytou vrstvu \\
Náhodný les & nízká & velmi vysoká & vysoká (paralelně) & snižuje rozptyl, nezhoršuje
  zkreslení \\
Boosting / AdaBoost & nízká & velmi vysoká & nízká (sekvenčně) & snižuje zkreslení; citlivý
  na šum a~odlehlé hodnoty \\
\bottomrule
\end{tabular}
\end{center}

\subsection{Klastrová analýza}
\label{ssec:clustering}

\subsubsection{Formulace úlohy a~míry vzdálenosti}

\emph{Klastrová analýza} (shlukování, \emph{cluster analysis}) rozděluje pozorované vzory do
skupin (klastrů, shluků) vzájemně podobných vzorů. \textbf{Předpoklad:} umíme měřit
\emph{vzdálenost} mezi vzory. Každý vzor je charakterizován $m$ numerickými veličinami.

\begin{defbox}[Míry vzdálenosti]
Pro vzory $x_1=(x_{11},\dots,x_{1m})$ a~$x_2=(x_{21},\dots,x_{2m})$:
\begin{enumerate}
\item \textbf{Hammingova} (\emph{Manhattan}, $L_1$):
      $d_H(x_1,x_2)=\sum_{i=1}^{m}|x_{1i}-x_{2i}|$
\item \textbf{Euklidovská} ($L_2$):
      $d_E(x_1,x_2)=\sqrt{\sum_{i=1}^{m}(x_{1i}-x_{2i})^2}$
\item \textbf{Čebyševova} ($L_\infty$):
      $d_C(x_1,x_2)=\max_i |x_{1i}-x_{2i}|$
\item \textbf{Minkowského metrika} ($L_q$) --- předchozí tři jsou jejími speciálními případy:
      $d_q(x_1,x_2)=\Big(\sum_{i=1}^{m}|x_{1i}-x_{2i}|^{q}\Big)^{1/q}$
\end{enumerate}
Platí $d_H=d_1$, $d_E=d_2$, $d_C=\lim_{q\to\infty}d_q$; pro vzor $(0,\dots,0)$ a~vzor
jednotkové vzdálenosti dávají všechny tři metriky hodnotu 1.
\end{defbox}

Volba metriky závisí na zúčastněných veličinách a~jejich rozsahu $\Rightarrow$ veličiny je
třeba \textbf{normalizovat} (dělením průměrem, směrodatnou odchylkou, rozsahem
$\max-\min$, \dots; viz kap.~\ref{sec:priprava}). Uvedené metriky předpokládají \emph{stejný
rozptyl} všech veličin. Mají-li veličiny \emph{různý} rozptyl, používá se

\begin{defbox}[Mahalanobisova vzdálenost]
\[ d_M^2(x_1,x_2)=(x_1-x_2)^{\!\top}S^{-1}(x_1-x_2), \]
kde $S$ je kovarianční matice. Zohledňuje jak různé rozptyly, tak korelace mezi veličinami
(pro $S=I$ se redukuje na kvadrát euklidovské vzdálenosti).
\end{defbox}

\begin{defbox}[Vzdálenost mezi dvěma klastry $C_k$ a~$C_l$]
\textbf{Metoda nejbližšího souseda} (\emph{single linkage}) --- minimum ze vzdáleností jejich
prvků:
$D(C_k,C_l)=\min_{x_i,x_j} d(x_i,x_j)$, $x_i\in C_k$, $x_j\in C_l$.

\textbf{Metoda nejvzdálenějšího souseda} (\emph{complete linkage}) --- maximum:
$D(C_k,C_l)=\max_{x_i,x_j} d(x_i,x_j)$.

\textbf{Metoda průměrné vzdálenosti} (\emph{average linkage}) --- průměr vzdáleností mezi
vzory ($n_k$, $n_l$ jsou počty vzorů v~klastrech):
\[ D(C_k,C_l)=\frac{1}{n_kn_l}\sum_{i=1}^{n_k}\sum_{j=1}^{n_l} d(x_i,x_j). \]

\textbf{Centroidní metoda} --- vzdálenost středů klastrů: $D(C_k,C_l)=d(c_k,c_l)$, kde $c_k$,
$c_l$ jsou středy (centroidy) klastrů.
\end{defbox}

\emph{Centroid} je střed klastru --- \uv{prototyp}, který jej reprezentuje. Podle tvaru klastru
a~metriky může mít tentýž klastr i~více centroidů současně.

\subsubsection{Metoda $k$-means}

\begin{defbox}[Algoritmus $k$-means (Lloyd, 1982)]
\begin{algorithmic}[1]
\State rozděl vzory \emph{náhodně} do $k$ klastrů $C_1,\dots,C_k$ (disjunktních, v~sjednocení
       všechny vzory)
\State \label{alg:km2} urči centroidy všech klastrů v~aktuálním rozdělení,
       např.\ $c_k=\frac{1}{|C_k|}\sum_{x\in C_k}x$
\For{každý vzor $x$}
   \State urči vzdálenosti $d(x,c_k)$ pro $1\le k\le K$
   \State nechť $d(x,c_l)=\min_k d(x,c_k)$
   \If{$x$ nepatří do klastru $l$} \State přesuň $x$ do klastru $l$ \EndIf
\EndFor
\If{došlo k~přesunu} \State goto krok~\ref{alg:km2} \Else \State \textbf{konec} \EndIf
\end{algorithmic}
Po skončení jsou klastry reprezentovány svými centroidy.

\textbf{Varianty:} při počátečním rozdělení prohlásit za centroidy \emph{prvních $k$ vzorů}
(eliminuje krok~\ref{alg:km2} v~první iteraci); upravovat centroidy \emph{po každém}
přesunu (v~cyklu, nikoli po jeho skončení).
\end{defbox}

\textbf{Rozšíření Lloydova algoritmu.} Střídavé minimalizační kroky neposkytují záruku
aproximace, přesto Lloydův algoritmus \emph{nikdy nezvýší} cenu počátečního shlukování
$\Rightarrow$ vyplatí se dodat algoritmu \emph{prokazatelně dobrou inicializaci}:
\begin{itemize}
\item \textbf{$k$-means++} (Arthur, Vassilvitskii, 2007) --- vyber $k$ z~$P$ vzorů
      s~pravděpodobností proporcionální kvadrátu vzdálenosti k~nejbližšímu už existujícímu
      centroidu:
      $\dfrac{d^2(x_i,c)}{\sum_{j=1}^{P}d^2(x_j,c)}$.
\item \textbf{LocalSearch++} (Lattanzi, Sohler, 2019) --- nejprve $k$-means++ pro volbu $k$
      počátečních centroidů; poté se vzorkují nové vzory s~analogickou pravděpodobností
      a~nahrazují existující centroidy, pokud tím sníží celkovou kvadratickou vzdálenost
      vzorů k~jejich nejbližším centroidům.
\end{itemize}

\textbf{Metoda $k$-medians} používá jako účelovou funkci Manhattanovu vzdálenost,
$d(x_i,c)=\|x_i-c\|_1$. Lze ukázat, že optimálním reprezentantem $c$ je pak \emph{medián}
datových bodů v~klastru \emph{podél každé dimenze}. Protože se medián volí v~každé dimenzi
nezávisle, výsledný $m$-rozměrný reprezentant obvykle \emph{nepatří} do původní datové
množiny. $k$-medians volí reprezentanty \emph{robustněji} než $k$-means --- medián není tak
citlivý na odlehlé hodnoty jako průměr.

\subsubsection{Metoda $k$-medoids, CLARA a~CLARANS}

\begin{defbox}[Algoritmus $k$-medoids]
Pro datovou množinu $D$ o~$n$ $d$-rozměrných bodech $x_1,\dots,x_n$ se hledá $k$
reprezentantů $y_1,\dots,y_k$ \emph{z~původní databáze} $D$ (počet $k$ zadává uživatel), které
minimalizují účelovou funkci danou součtem vzdáleností bodů k~jejich nejbližším
reprezentantům:
\[ O=\sum_{i=1}^{n}\min_{j} d(x_i,y_j). \]
Časová složitost jedné iterace je $O(k\,n\,d)$; obvykle stačí malý konstantní počet iterací.

\textbf{Proč reprezentanty z~$D$?} \emph{(1)}~Centroidy $k$-means mohou být zkresleny odlehlými
hodnotami a~ocitnout se v~prázdných oblastech nereprezentativních pro klastr, což může vést
k~částečnému splynutí klastrů. \emph{(2)}~Pro složitá data (množina časových řad různé délky)
může být optimální centrální reprezentant obtížně spočitatelný. $k$-medoids stačí vhodná
funkce podobnosti, lze jej tedy definovat pro \emph{libovolný typ dat}.

\textbf{Postup:} \emph{hill-climbing} --- množina reprezentantů $S$ se inicializuje body
z~$D$ a~iterativně se zlepšuje výměnou bodu z~$S$ za datový bod z~$D$. Strategie výměny:
\begin{enumerate}
\item vyzkoušet všech $|S|\cdot|D|$ možností a~vybrat nejlepší --- \emph{extrémně drahé};
\item použít \emph{náhodně vybranou} množinu $r$ párů $(x,y)$ ($x\in D$, $y\in S$) a~vyměnit
      nejlepší z~nich. Vyžaduje čas proporcionální $r\cdot|D|$; algoritmus konverguje, když
      se účelová funkce nezlepší, nebo je-li průměrné zlepšení pod uživatelem zadaným prahem.
\end{enumerate}
$k$-medoids je \emph{pomalejší} než $k$-means, ale použitelný na různé typy dat.
\end{defbox}

\begin{defbox}[Škálovatelné varianty: CLARA a~CLARANS]
Škálovatelné implementace $k$-medoids pro data, která se nevejdou do hlavní paměti. Dědí výhody
i~nevýhody bázového algoritmu (generalizovatelnost na různé typy dat, ale vysokou výpočetní
složitost).

\textbf{CLARA} (\emph{Clustering LARge Applications}) vychází z~varianty \emph{PAM}
(\emph{Partitioning Around Medoids}), která testuje na výměnu \emph{všech} $k(n-k)$ párů
medoid--nemedoid a~vymění nejlepší z~nich, až do konvergence k~lokálnímu optimu --- $O(kn^2)$
výpočtů vzdálenosti, tj.\ $O(kn^2d)$ času na iteraci, \emph{značně drahé}. CLARA proto PAM
aplikuje na \emph{vzorek} velikosti $f\cdot n$ ($f\ll 1$) a~zbylé body přiřadí k~nalezeným
medoidům; postup se opakuje nad nezávisle volenými vzorky a~vybere se nejlepší shlukování.
Iterace pak stojí $O(kf^2n^2d+k(n-k))$ --- pro malá $f$ o~řády rychleji. \textbf{Hlavní
problém:} vzorek nemusí obsahovat dobrou volbu medoidů.

\textbf{CLARANS} (\emph{\dots\ based on RANdomized Search}) pracuje s~\emph{plnými} daty, čímž
se problému špatných vzorků vyhýbá. Testuje výměny \emph{náhodných} medoidů za \emph{náhodné}
nemedoidy; zlepší-li se kvalita, výměna se provede definitivně, jinak roste počítadlo
neúspěchů --- lokální optimum je nalezeno po \emph{MaxAttempt} neúspěšných pokusech. Celé se
opakuje \emph{MaxLocal}-krát. CLARANS prozkoumá \emph{větší rozmanitost} prostoru řešení než
CLARA.
\end{defbox}

\subsubsection{Hierarchické shlukování}

\begin{defbox}[Aglomerativní hierarchické shlukování]
Přístup \uv{zdola nahoru} (\emph{bottom-up}).

\textbf{Inicializace:} \emph{(1)}~urči vzájemné vzdálenosti mezi všemi vzory;
\emph{(2)}~přiřaď každý vzor do samostatného klastru.

\textbf{Hlavní cyklus:} dokud je klastrů více než jeden: \emph{(1)}~najdi dva vzájemně
nejbližší klastry a~slouč je; \emph{(2)}~pro nový klastr spočítej jeho vzdálenost od
ostatních klastrů.

Výsledek se zobrazuje jako \textbf{dendrogram} --- ukazuje (zleva doprava) postupné slučování
klastrů; výška spojení odpovídá vzdálenosti, při níž ke sloučení došlo. \emph{Optimální počet
klastrů není znám předem} --- volí se řezem dendrogramu ve vhodné výšce.
\end{defbox}

\textbf{Divizivní} hierarchické shlukování (\uv{shora dolů}) postupuje obráceně: začíná
jedním klastrem obsahujícím všechny vzory a~ten postupně dělí. Nejznámějším příkladem
v~kontextu sítí je Girvanové--Newmanův algoritmus (kap.~\ref{sec:sna}).

\begin{defbox}[CURE (\emph{Clustering Using REpresentatives})]
Aglomerativní hierarchický algoritmus, \emph{velmi rychlý} a~schopný odhalit klastry
\emph{libovolného tvaru} (na rozdíl např.\ od CLARANS). Postup:
\begin{enumerate}
\item Vyber vzorek $s$ bodů z~databáze $D$ o~$n$ bodech.
\item Rozděl $s$ bodů na $p$ částí velikosti $s/p$.
\item Každou část shlukuj \emph{nezávisle} hierarchickým slučováním do $k'$ klastrů;
      celkový počet $k'\cdot p$ klastrů přes všechny části je stále větší než uživatelem
      požadované $k$.
\item Nad $k'\cdot p$ klastry proveď hierarchické shlukování až na požadovaných $k$.
\item Každý z~$(n-s)$ nevzorkovaných bodů přiřaď ke klastru s~nejbližším reprezentantem.
\end{enumerate}
\end{defbox}

\subsubsection{Vektorová kvantizace: algoritmus LVQ}

\begin{defbox}[LVQ --- \emph{Learning Vector Quantization}]
Slouží pro \textbf{shlukování s~učitelem} (třídy jsou označeny $C$):
\begin{algorithmic}[1]
\State inicializuj všechny váhové vektory $w_j(0)$, rychlosti učení $\mu(0)$, $k\gets 0$
\While{není splněna zastavovací podmínka}
  \For{každý trénovací vzor $x$}
     \State urči index $j=q$ váhového vektoru s~\emph{minimální} euklidovskou vzdáleností,
            tj.\ $\min_j\|x-w_j\|_2^2$
     \If{třída $x$ souhlasí s~třídou $w_q$}
        \State $w_q(k+1)\gets w_q(k)+\mu(k)\big(x-w_q(k)\big)$ \Comment{přitáhni}
     \Else
        \State $w_q(k+1)\gets w_q(k)-\mu(k)\big(x-w_q(k)\big)$ \Comment{odtlač}
     \EndIf
  \EndFor
  \State $k\gets k+1$; sniž rychlost učení, např.\ $\mu(k)=\mu(k-1)/(k+1)$ pro $k>0$
\EndWhile
\end{algorithmic}
Váhové vektory se obvykle inicializují \emph{prvními $m$ vektory} trénovací množiny (musí
zahrnovat vzory ze všech tříd) --- tím se neurony automaticky \uv{kalibrují} s~třídami.
\end{defbox}

\subsubsection{Mřížkové a~hustotní metody}

\begin{defbox}[Mřížkové metody (např.\ MAFIA)]
\begin{itemize}
\item Data se diskretizují do $p$ (ekvidistantních) intervalů, což pro $d$-rozměrná data dá
      $p^d$ hyperkrychlí.
\item Výsledné hyperkrychle jsou stavebními bloky pro shlukování.
\item Pomocí prahu hustoty $\tau$ se identifikují \emph{husté} hyperkrychle (alespoň $\tau$
      bodů).
\item Klastr libovolného tvaru se skládá z~více hustých oblastí, které jsou \emph{propojeny};
      dvě mřížkové oblasti jsou propojené, sdílejí-li stěnu.
\end{itemize}
\textbf{Cíl:} určit klastry jako husté vzájemně propojené oblasti. Hustým mřížkovým buňkám
odpovídají uzly grafu, hrany reprezentují sousedskou propojenost; data pokrytá buňkami
z~komponent souvislosti tohoto grafu tvoří výsledné klastry.
\end{defbox}

\begin{defbox}[DBSCAN --- \emph{Density-Based Spatial Clustering of Applications with Noise}]
Stavebními bloky jsou místo mřížkových buněk \emph{jednotlivé datové body}, resp.\ kruhové
oblasti okolo nich; hustota je dána počtem bodů ležících v~jejich radiu $\varepsilon$
(včetně bodu samotného):
\begin{enumerate}
\item \textbf{Jádrový bod} (\emph{core point}) --- kulová oblast o~radiu $\varepsilon$ okolo
      něj obsahuje alespoň $\tau$ datových bodů.
\item \textbf{Hraniční bod} (\emph{border point}) --- oblast obsahuje \emph{méně} než $\tau$
      bodů, ale zároveň alespoň jeden jádrový bod ve vzdálenosti $\varepsilon$.
\item \textbf{Šumový bod} (\emph{noise point}) --- není ani jádrový, ani hraniční.
\end{enumerate}

\textbf{Algoritmus:}
\begin{enumerate}
\item Urči jádrové, hraniční a~šumové body.
\item Zkonstruuj graf propojenosti, kde každý uzel odpovídá jádrovému bodu; hrana mezi dvěma
      jádrovými body existuje \emph{právě když} leží od sebe ve vzdálenosti nejvýše
      $\varepsilon$.
\item Najdi všechny komponenty souvislosti tohoto grafu --- odpovídají klastrům (postaveným
      na jádrových bodech). Hraniční body se přiřadí ke klastru, s~nímž mají nejvyšší úroveň
      propojenosti; šumové body se ohlásí jako \emph{odlehlé hodnoty}.
\end{enumerate}
\end{defbox}

\begin{center}
\small
\begin{tabular}{@{}p{2.6cm}p{2.3cm}p{2.4cm}p{2.2cm}p{5.0cm}@{}}
\toprule
\textbf{Metoda} & \textbf{Tvar klastrů} & \textbf{Počet klastrů} & \textbf{Odlehlé hodnoty}
  & \textbf{Poznámka} \\
\midrule
$k$-means & kulovité & zadán & citlivá & rychlá; závislá na inicializaci $\to$
  $k$-means++ \\
$k$-medians & kulovité & zadán & robustnější & medián po dimenzích; reprezentant není
  z~dat \\
$k$-medoids & kulovité & zadán & robustní & reprezentanti \emph{z~dat}; libovolný typ dat;
  pomalejší \\
CLARA / CLARANS & kulovité & zadán & robustní & škálovatelné $k$-medoids (vzorkování /
  náhodné prohledávání) \\
Hierarchické & podle metriky spojování & určí se z~dendrogramu & citlivé & bez nutnosti
  zadat $k$ předem; $O(n^2)$--$O(n^3)$ \\
CURE & libovolný & zadán & robustní & vzorkování $+$ hierarchické slučování po částech \\
LVQ & kulovité & dán počtem neuronů & --- & shlukování \emph{s~učitelem} \\
Mřížkové (MAFIA) & libovolný & vyplyne & šum vypadne & citlivé na granularitu mřížky \\
DBSCAN & libovolný & vyplyne & \emph{detekuje je} & parametry $\varepsilon$, $\tau$;
  problém s~různou hustotou klastrů \\
\bottomrule
\end{tabular}
\end{center}

\subsection{Shrnutí ke~zkoušce}

\textbf{Asociační pravidla.}
\begin{itemize}
\item \textbf{Support} $=$ podíl transakcí s~$i$ i~$j$; \textbf{confidence} $=$ podíl
      transakcí s~$i$ i~$j$ mezi transakcemi s~$i$; \textbf{lift}
      $=p(i\wedge j)/\big(p(i)p(j)\big)$. Lift $<1$ $\Rightarrow$ pravidlo horší než náhoda,
      negace závěru může být lepší. Umět spočítat všechny tři z~malé tabulky transakcí.
\item \textbf{Apriori} využívá \emph{uzavřenost dolů} (každá podmnožina frekventované množiny
      je frekventovaná), prohledává \emph{do šířky}; \textsc{Apriori-Gen} $=$ krok spojení
      (kombinace shodné v~$k-2$ kategoriích) $+$ krok prořezání (chybí-li podkombinace délky
      $k-1$). Prostor pravidel je $O(2^m)$, MBA produkuje \emph{nadměrné množství pravidel}.
\item \textbf{Taxonomie a~virtuální položky} řeší volbu úrovně položek; vzácné položky výše
      v~taxonomii, časté níže; virtuální položky nesou informaci o~transakci (den, čas,
      zákazník) --- pozor na redundanci pravidel.
\item \textbf{MS-Apriori}: problém vzácných položek $\to$ každá položka má vlastní
      $\mathrm{MIS}$, minsup pravidla je \emph{minimum} MIS jeho položek, plus omezení
      rozdílu supportů $\varphi$. \emph{Uzavřenost dolů přestává platit} $\Rightarrow$ položky
      se setřídí vzestupně podle MIS, používá se množina semen $L$ (ne $F_1$)
      a~\texttt{tailCount} pro generování pravidel. Umět uvést protipříklad k~uzavřenosti dolů.
\item \textbf{CAR}: transakce označkované třídou, pravidlo $X\to y$; dolují se
      \emph{přímo v~jednom kroku} přes pravidlové položky $(\mathit{condset},y)$; lze zadat
      různé minsupy různým třídám (101~\% $=$ zakázat třídu).
\item \textbf{Sekvenční vzory}: sekvence $=$ uspořádaný seznam množin položek;
      \emph{velikost} $=$ počet prvků, \emph{délka} $=$ počet položek; podsekvence přes
      vnořené inkluze. Algoritmus \textbf{GSP} pracuje obdobně jako Apriori. Pozor:
      $\langle\{3\}\{8\}\rangle$ a~$\langle\{3,8\}\rangle$ se navzájem \emph{neobsahují}.
\item \textbf{FP-Growth}: 2 průchody daty, žádné generování kandidátů; FP-strom sdílí
      \emph{prefixy} (pořadí položek podle klesajícího supportu --- heuristika), spojové
      seznamy položek; extrakce zdola nahoru \emph{rozděl a~panuj} přes \emph{prefixové
      podstromy} a~\emph{podmíněné FP-stromy} (aktualizovat počítadla $\to$ odstranit uzly
      položky $\to$ odstranit nefrekventované položky). Nejhorší případ: unikátní transakce
      $\Rightarrow$ strom větší než data.
\end{itemize}

\textbf{Učení s~učitelem.}
\begin{itemize}
\item \textbf{Rozhodovací stromy}: TDIDT (rozděl a~panuj), volba atributu podle
      \emph{minimální vážené entropie} $H(A)=\sum_v\frac{|A(v)|}{|D|}H(A(v))$, ekvivalentně
      maximálního \emph{informačního zisku}. Umět projít příklad s~úvěrem (příjem 0,5747;
      konto 0,6667; pohlaví 0,9183; nezaměstnaný 0,825) a~převést strom na pravidla.
      Složitost indukce $O(np\log p)$, klasifikace $O(q\log p)$.
\item \textbf{ID3} (Quinlan 1986) --- hladový, informační zisk. \textbf{C4.5} (1993) ---
      chybějící a~spojitá data, prořezávání validací (\emph{reduced-error pruning},
      \emph{subtree raising}), \emph{rule post-pruning}, odhad přesnosti dolní mezí
      $\mathit{acc}-1{,}96\sqrt{\mathit{acc}(1-\mathit{acc})/n}$, \emph{poměrný} informační
      zisk (penalizuje nadměrné větvení). \textbf{C5.0} --- komerční, boosting.
      \textbf{CART} --- \emph{binární} stromy, Gini index $G(v_i)=1-\sum_j p_j^2$ (nižší $=$
      lepší), kritérium $\Phi=2P_LP_R\sum_j|P(C_j|t_L)-P(C_j|t_R)|$. \textbf{CHAID} ---
      $\chi^2$, automatické seskupování hodnot kategoriálních atributů.
\item \textbf{Bayes}: $H_{MAP}=\arg\max_H P(E\mid H)P(H)$ (jmenovatel lze zanedbat);
      \emph{naivní} předpoklad podmíněné nezávislosti $\Rightarrow$
      $\arg\max_{H_t}P(H_t)\prod_i P(E_i\mid H_t)$. Umí nekompletní vzory; problém nulových
      pravděpodobností. Umět spočítat příklad (0,1042 vs.\ 0,0416).
\item \textbf{SVM}: maximální okraj $2/\|w\|$, primární úloha $\min\frac{\langle w\cdot
      w\rangle}{2}$ za podmínek $y_i(\langle w\cdot x_i\rangle+b)\ge 1$; Kuhn--Tuckerovy
      podmínky jsou pro konvexní úlohu s~lineárními omezeními nutné \emph{i~postačující};
      $\alpha_i>0$ jen pro \emph{podpůrné vektory}; Wolfeho duál. \textbf{Kernelový trik}:
      duál potřebuje jen skalární součiny $\Rightarrow$ $K(x,z)=\langle\phi(x)\cdot\phi(z)\rangle$,
      $\phi$ nikdy nevyčíslíme (Mercerova věta). Jen 2 třídy, jen reálný prostor,
      neinterpretovatelné.
\item \textbf{Logistická regrese}: $\ln\frac{p}{1-p}=\beta^{\!\top}x$, tedy
      $p=\sigma(z)=1/(1+e^{-z})$; učení maximální věrohodností / minimalizací křížové
      entropie; $\sigma'=\sigma(1-\sigma)$ $\Rightarrow$ adaptační pravidlo
      $\Delta\beta_j=\eta\sum_i(y_i-\sigma(z_i))x_{ij}$.
\item \textbf{Diskriminační analýza}: $g_j(x)$ pro každou třídu, klasifikuje se maximem; pro
      dvě třídy stačí znaménko $g=g_1-g_2$; optimální je $g_j\propto P(C_j\mid x)$; normální
      rozdělení $\to$ kvadratická funkce, při shodných kovariančních maticích \emph{lineární}
      $\Rightarrow$ stačí odhadnout $\mu_j$ a~$S$.
\item \textbf{ELM}: náhodné parametry skrytých neuronů (i~bez znalosti dat), výstupní váhy
      \emph{jednou} pseudoinverzí $\beta=H^{+}T$; regularizace $I/C$ na diagonále. Rychlé,
      bez lokálních minim, funguje s~libovolnou omezenou po částech spojitou přenosovou funkcí.
\item \textbf{Ensembly}: \emph{bagging} snižuje \emph{rozptyl} ($\sigma^2\to\sigma^2/k$ pro
      i.i.d.), bootstrap obsahuje $\approx 63{,}2\,\%$ různých bodů ($1-1/e$), stromy jsou
      ideální komponenty (nízké zkreslení, vysoký rozptyl); zkreslení nesnižuje.
      \emph{Náhodné lesy} řeší korelaci komponent --- rozptyl je
      $\rho\sigma^2+(1-\rho)\sigma^2/k$, první člen na $k$ nezávisí; náhodnost přes bootstrap
      \emph{a} náhodnou podmnožinu $m\approx\sqrt{\#\text{příznaků}}$ příznaků v~každém uzlu.
      \emph{Boosting}/\textbf{AdaBoost} snižuje \emph{zkreslení}, váhy chybně klasifikovaných
      vzorů rostou, $\alpha_t=\frac12\ln\frac{1-\epsilon_t}{\epsilon_t}$, běží
      \emph{sekvenčně}, je citlivý na šum. Umět srovnat lesy vs.\ boosting (paralelní
      a~levné vs.\ vyčerpávající a~přesnější při stejném počtu stromů).
\end{itemize}

\textbf{Klastrová analýza.}
\begin{itemize}
\item \textbf{Metriky}: Hammingova/Manhattan ($L_1$), euklidovská ($L_2$), Čebyševova
      ($L_\infty$), Minkowského ($L_q$, ostatní jsou její speciální případy);
      \textbf{Mahalanobisova} $d_M^2=(x_1-x_2)^{\!\top}S^{-1}(x_1-x_2)$ pro veličiny
      s~různým rozptylem. Před shlukováním \emph{normalizovat}.
\item \textbf{Vzdálenost klastrů}: nejbližší soused (min), nejvzdálenější soused (max),
      průměrná vzdálenost, centroidní metoda.
\item \textbf{$k$-means} (Lloyd 1982): náhodné rozdělení $\to$ centroidy $\to$ přesuny
      $\to$ opakuj, dokud dochází k~přesunům. Nikdy nezvýší cenu, ale bez záruky aproximace
      $\Rightarrow$ \textbf{$k$-means++} (pravděpodobnost $\propto$ kvadrát vzdálenosti
      k~nejbližšímu centroidu), \textbf{LocalSearch++}.
\item \textbf{$k$-medians} (Manhattan, medián po dimenzích, robustnější) vs.\
      \textbf{$k$-medoids} (reprezentanti \emph{z~dat}, hill-climbing s~výměnami,
      $O(knd)$ na iteraci, libovolný typ dat). \textbf{PAM} testuje všech $k(n-k)$ výměn
      ($O(kn^2d)$); \textbf{CLARA} vzorkuje ($O(kf^2n^2d+k(n-k))$, riziko špatného vzorku),
      \textbf{CLARANS} prohledává náhodně nad plnými daty (\emph{MaxAttempt},
      \emph{MaxLocal}).
\item \textbf{Hierarchické shlukování} zdola nahoru: každý vzor svůj klastr $\to$ slučuj dva
      nejbližší; výstupem \textbf{dendrogram}, počet klastrů se volí až řezem.
      \textbf{CURE}: vzorek $\to$ $p$ částí $\to$ nezávisle hierarchicky na $k'$ klastrů
      $\to$ hierarchicky na $k$ $\to$ přiřaď zbytek; rychlé, libovolný tvar klastrů.
\item \textbf{LVQ}: shlukování \emph{s~učitelem} --- vítězný váhový vektor se ke vzoru
      přitáhne při souhlasu tříd a~odtlačí při nesouhlasu; $\mu$ klesá jako $\mu(k-1)/(k+1)$.
\item \textbf{Mřížkové metody} (MAFIA): $p^d$ hyperkrychlí, prah hustoty $\tau$, klastry
      jako propojené husté oblasti (komponenty souvislosti grafu buněk).
      \textbf{DBSCAN}: \emph{jádrový} bod ($\ge\tau$ bodů v~radiu $\varepsilon$),
      \emph{hraniční} ($<\tau$, ale jádrový bod do $\varepsilon$), \emph{šumový};
      klastry $=$ komponenty souvislosti grafu jádrových bodů, šumové body se hlásí jako
      odlehlé. Umí klastry \emph{libovolného tvaru} a~sama detekuje odlehlé hodnoty.
\end{itemize}

%=====================================================================
\newpage
\section{Charakteristická a~diskriminační pravidla a~jejich zajímavost}
\label{sec:pravidla}
%=====================================================================

Věta okruhu žádá \emph{metody pro extrakci charakteristických diskriminačních pravidel
a~měření jejich zajímavosti}. Kapitola je postavena především na slidech kurzu
\emph{Internet a~klasifikační metody} (Holeňa), hlavně 5.~přednášky \emph{Kdy je
klasifikace srozumitelná uživateli?}: proč pravidla (odd.~\ref{ssec:srozumitelnost}),
extrakce pravidel z~klasifikačních stromů (odd.~\ref{ssec:extrakce}) a~evolucí
(odd.~\ref{ssec:evoluce}), observační pravidla (odd.~\ref{ssec:observacni}), statistický
pohled (odd.~\ref{ssec:stat-pohled}), měření zajímavosti (odd.~\ref{ssec:zajimavost})
a~aplikace (odd.~\ref{ssec:aplikace-pravidel}). Pojmy \emph{charakteristické}
a~\emph{diskriminační} pravidlo s~mírami t-weight a~d-weight
(odd.~\ref{ssec:char-diskr}) jsou doplněny podle Han--Kamber, protože je okruh jmenuje
výslovně.

\begin{poznbox}[Zdroje kapitoly]
\small
Ze šesti přednášek kurzu \emph{Internet a~klasifikační metody}
(\texttt{ikm1}--\texttt{ikm6}) je pro tuto kapitolu klíčová \texttt{ikm5}
(srozumitelnost klasifikace: pravidla, jejich extrakce a~observační kalkul); dále se zde
využívá \texttt{ikm1} (aplikace: spam, doporučování, hrozby v~síti), \texttt{ikm2} (míry
kvality binární klasifikace) a~\texttt{ikm3} (diskriminační analýza). Látka \texttt{ikm4}
(SVM: maximalizace pásu, opěrné vektory, kernel trick) a~\texttt{ikm6} (týmy
klasifikátorů: bagging, boosting, náhodné lesy) se kryje s~kap.~\ref{sec:metody}, matice
záměn a~ROC/AUC z~\texttt{ikm2} s~kap.~\ref{sec:reprez} --- zpracovány jsou tam. Slidy
\texttt{nmr1}--\texttt{nmr5} z~téže složky patří jinému (volitelnému) Holeňovu kurzu
a~látku kapitoly nepokrývají. Podle Han--Kamber je jen odd.~\ref{ssec:char-diskr};
regresní a~diskriminační analýzu dokládají i~slidy NDBI023 (\emph{Data Analysis}).
\end{poznbox}

\subsection{Pravidla jako srozumitelná forma klasifikace}
\label{ssec:srozumitelnost}

Klasifikátor je matematicky zobrazení $f\colon X\to\{C_1,\dots,C_m\}$ z~prostoru příznaků
do konečné množiny tříd; přesněji $f\colon X\times\Theta\to\{C_1,\dots,C_m\}$, kde učení
znamená volbu parametrů $\theta\in\Theta$ podle zkušenosti s~kvalitou klasifikace při
předchozích volbách (odd.~\ref{ssec:supervised}). Konkrétní naučený předpis --- třeba
znaménko $\langle w,x\rangle+b$ u~SVM nebo složená nelineární funkce u~neuronové sítě ---
je ale pro běžného uživatele málo srozumitelný. Popisuje-li binární klasifikátor
$f\colon X\to\{1,0\}$ pravdivost nějaké implikace $\varphi\Rightarrow\psi$ či ekvivalence
$\varphi\Leftrightarrow\psi$ mezi vlastnostmi objektu, stačí místo celého $f$ znát toto
\emph{pravidlo}: pravidlo logiky je srozumitelnější než matematická formule.

\begin{defbox}[Booleovská a~fuzzy pravidla]
\textbf{Booleovská} pravidla pracují s~pravdivostními hodnotami $t(\varphi)\in\{0,1\}$:
\emph{ekvivalence} $\varphi\Leftrightarrow\psi$ je pravdivá, právě když
$t(\varphi)=t(\psi)$; \emph{implikace} $\varphi\Rightarrow\psi$ je pravdivá, právě když
$t(\varphi)=0$, nebo $t(\varphi)=t(\psi)=1$. Pravidla se typicky uplatňují na objektech,
kde antecedent platí ($t(\varphi)=1$) --- tam implikace splývá s~ekvivalencí.

\textbf{Fuzzy} pravidla připouštějí stupně pravdivosti $t(\varphi)\in[0,1]$. Fuzzy
implikace je \emph{reziduum} t-normy $T$:
\[ t(\varphi\Rightarrow\psi)=\sup\{\,c\in[0,1]: T(t(\varphi),c)\le t(\psi)\,\}, \]
fuzzy ekvivalence $t(\varphi\Leftrightarrow\psi)
=T\big(t(\varphi\Rightarrow\psi),\,t(\psi\Rightarrow\varphi)\big)$. Důležité speciální
případy: je-li $t(\varphi)\le t(\psi)$, je $t(\varphi\Rightarrow\psi)=1$; je-li
$t(\varphi)=t(\psi)$, je $t(\varphi\Leftrightarrow\psi)=1$.
\end{defbox}

Srozumitelnost je důvod, proč se pravidla používají tam, kde výsledek klasifikace
posuzuje člověk: analytik malwaru chce vědět, \emph{proč} je program podezřelý,
a~doporučovací systémy vysvětlením zvyšují důvěru uživatele v~doporučení
(odd.~\ref{ssec:aplikace-pravidel}).

\subsection{Charakteristická a~diskriminační pravidla: t-weight a~d-weight}
\label{ssec:char-diskr}

\mimo{} Oddíl je zpracován podle Han--Kamber --- okruh oba druhy pravidel jmenuje
výslovně.

\begin{defbox}[Charakterizace vs.\ diskriminace]
\textbf{Charakterizace dat} (\emph{data characterization}) --- shrnutí obecných
vlastností \emph{cílové třídy} (\uv{jak vypadá typický klient s~hypotékou?}); výsledkem
jsou \emph{charakteristická pravidla}. \textbf{Diskriminace dat} (\emph{data
discrimination}) --- porovnání cílové třídy s~jednou či více \emph{kontrastními třídami}
(\uv{čím se liší od klienta bez hypotéky?}); výsledkem jsou \emph{diskriminační
pravidla}. Vlastnost může být pro cílovou třídu \emph{typická}, a~přesto nemít žádnou
\emph{rozlišovací} sílu --- je-li stejně typická i~pro třídu kontrastní.
\end{defbox}

Pravidla tohoto druhu se získávají \emph{generalizací dat}: \emph{indukce orientovaná na
atributy} (\emph{attribute-oriented induction}, AOI) nahrazuje hodnoty atributů
obecnějšími podle koncepčních hierarchií (město $\to$ kraj, datum narození $\to$ věková
kategorie), atributy s~mnoha hodnotami bez hierarchie odstraňuje (jméno, telefon)
a~shodné zobecněné záznamy slučuje s~čítačem \texttt{count}; totéž provádí operace
\emph{roll-up} / \emph{drill-down} nad OLAP kostkou (kap.~\ref{sec:paradigmata}). Nad
zobecněnými záznamy se pak definují dvě váhy:

\begin{defbox}[t-weight, d-weight a~kvantitativní pravidla]
Nechť $q_1,\dots,q_N$ jsou zobecněné záznamy cílové třídy.

\textbf{Typičnost} (t-weight) záznamu $q_a$:
\[ t\text{-}\mathrm{weight}(q_a)=\frac{\mathrm{count}(q_a)}
   {\sum_{i=1}^{N}\mathrm{count}(q_i)}\in[0,1], \]
tj.\ podíl objektů cílové třídy, které $q_a$ pokrývá (součet přes všechny záznamy je
100~\%). Vystupuje v~\textbf{kvantitativním charakteristickém pravidle} (nutná podmínka
--- implikace \emph{z}~třídy):
\[ \forall X\ \ \mathrm{target}(X)\Rightarrow \mathrm{cond}_1(X)\,[t{:}\,w_1]\ \vee\
   \cdots\ \vee\ \mathrm{cond}_N(X)\,[t{:}\,w_N]. \]

\textbf{Rozlišovací váha} (d-weight) záznamu $q_a$ vůči cílové třídě $C_t$ a~kontrastním
třídám $C_1,\dots,C_m$:
\[ d\text{-}\mathrm{weight}(q_a)=\frac{\mathrm{count}(q_a\in C_t)}
   {\mathrm{count}(q_a\in C_t)+\sum_{j=1}^{m}\mathrm{count}(q_a\in C_j)}\in[0,1], \]
tj.\ podíl objektů pokrytých záznamem $q_a$, které patří do cílové třídy. Vystupuje
v~\textbf{kvantitativním diskriminačním pravidle} (postačující podmínka --- implikace
\emph{do} třídy):
\[ \forall X\ \ \mathrm{target}(X)\Leftarrow \mathrm{cond}(X)\,[d{:}\,w]. \]
\emph{Kvantitativní popisné pravidlo} kombinuje obojí ($\Leftrightarrow$,
$[t{:}\,w,\ d{:}\,w']$).
\end{defbox}

\begin{defbox}[Příklad --- klienti s~hypotékou (200) vs.\ bez hypotéky (800)]
Primární relace po generalizaci:
\begin{center}
\small
\begin{tabular}{@{}lllrrrr@{}}
\toprule
Věk & Příjem & Region & \multicolumn{2}{c}{počet} & $t$-weight & $d$-weight \\
 & & & cíl & kontrast & (cíl) & \\
\midrule
30--40 & vysoký & Praha & 90 & 10 & \textbf{45~\%} & \textbf{90~\%} \\
30--40 & střední & Praha & 50 & 150 & 25~\% & 25~\% \\
40--50 & vysoký & mimo Prahu & 40 & 160 & 20~\% & 20~\% \\
20--30 & nízký & mimo Prahu & 20 & 480 & 10~\% & 4~\% \\
\bottomrule
\end{tabular}
\end{center}
Výpočet pro první záznam: $t=90/200=45\,\%$, $d=90/(90+10)=90\,\%$.

\textbf{Klasická past:} 3.~řádek má $t=20\,\%$, ale $d=40/200=20\,\%$ je \emph{přesně}
apriorní podíl cílové třídy ($200/1000=20\,\%$) --- rozlišovací síla je nulová;
4.~řádek s~$d=4\,\%$ hluboko pod apriorním podílem hypotéku spíše \emph{vylučuje}.
Vysoká typičnost tedy sama o~sobě neznamená rozlišovací sílu: d-weight je vždy třeba
srovnávat s~apriorním podílem třídy.
\end{defbox}

\subsection{Extrakce pravidel z~klasifikačních stromů}
\label{ssec:extrakce}

Nejpřímější metoda extrakce pravidel vychází z~klasifikačního stromu: uzly binárního
stromu testují booleovské podmínky (typicky $x_i\le c$, tedy řezy kolmé na souřadnicové
osy), vstup podle výsledku testu propadá doleva či doprava a~každý list odpovídá jediné
třídě.

\begin{defbox}[Pravidla ze stromu]
\begin{itemize}
\item Každé cestě kořen $\to$ list odpovídá \emph{konjunkce} podmínek testovaných na
      cestě; třída listu dává konsekvent:
      $\big(\mathrm{cond}_1\wedge\cdots\wedge\mathrm{cond}_k\big)\Rightarrow C_j$.
\item Cesty do listů \emph{téže} třídy se spojí \emph{disjunkcí} --- třída je sjednocením
      obdélníkových oblastí, které jednotlivé cesty vymezují.
\end{itemize}
\end{defbox}

\textbf{Učení stromu} je optimalizace dělení v~uzlu vůči volbě proměnné a~dělicího bodu
(podprostory spojitých vstupů, podmnožiny diskrétních); dvě nejčastěji optimalizovaná
kritéria jsou \emph{informační zisk}
$IG(\text{rodič};L,P)=H_{\text{rodič}}-\big(\tfrac{n_L}{n}H_L+\tfrac{n_P}{n}H_P\big)$
a~\emph{Giniho index nečistoty} $1-\sum_j p_j^2$; podrobně (TDIDT, ID3, C4.5, CART)
odd.~\ref{ssec:supervised}. Nové vstupy se klasifikují průchodem uzlových testů;
v~listech se předpokládá, že rozdělení nových vstupů odpovídá rozdělení učicích dat.

\textbf{Prořezávání} (\emph{pruning}) snižuje přeučení a~současně extrahovaná pravidla
\emph{zkracuje} (srov.\ prořezávání pravidel v~C4.5, odd.~\ref{ssec:supervised}).
\textbf{Zobecněná dělení} --- obecná nadrovina či kvadratická plocha místo řezu kolmého
na osu --- zlepšují aproximaci hranice tříd, ale snižují srozumitelnost pravidel
získaných konjunkcí podmínek.

\subsection{Evoluční konstrukce pravidel}
\label{ssec:evoluce}

Pravidla lze konstruovat i~genetickým algoritmem --- informatickou aplikací biologického
vývoje druhů: křížením a~mutací řetězce (genomu) kódujícího pravidlo se výrazněji
posilují žádoucí vlastnosti, zde míry kvality klasifikace (přesnost AC, precision PR,
TPr, TNr či AUC --- odd.~\ref{ssec:zajimavost}). Postup je vhodný, jsou-li antecedent
i~konsekvent konjunkcemi malého počtu atomů. Všechny evoluční algoritmy jsou
\emph{populační}; podle toho, co je jedincem populace, se rozlišují dva přístupy:

\begin{center}
\small
\begin{tabular}{@{}lp{5.2cm}p{6.8cm}@{}}
\toprule
\textbf{Přístup} & \textbf{Jedinec} & \textbf{Vlastnosti} \\
\midrule
michiganský & jedno pravidlo; hledá se celá populace & jednoduchý, ale populace zbytečně
  konverguje k~podobným pravidlům \\
pittsburghský & celý hledaný soubor pravidel & výsledkem rovnou konzistentní sada;
  evoluce populace souborů je složitější a~výpočetně náročnější \\
\bottomrule
\end{tabular}
\end{center}

\subsection{Observační pravidla: konfidence, podpora a~testy hypotéz}
\label{ssec:observacni}

Booleovská logika umí o~datech vypovědět jen dvojím způsobem: \uv{\emph{všechny} objekty
splňují $\varphi\Rightarrow\psi$} (nastává zřídka), nebo \uv{\emph{existuje} objekt
splňující $\varphi\wedge\psi$} (slabá informace). Pro užitečnější popis dat byl navržen
\emph{observační kalkul} (Hájek--Havránek 1978, základ české metody GUHA).

\begin{defbox}[Observační pravidla a~zobecněné kvantifikátory]
\emph{Observační pravidlo} má tvar $\varphi\sim\psi$, kde $\sim$ je \emph{zobecněný
kvantifikátor}: pravdivost pravidla zachycuje výsledek statistické analýzy dat
odpovídajících $\varphi$ a~$\psi$, shrnutých čtyřpolní tabulkou --- $a$ objektů splňuje
$\varphi\wedge\psi$, $b$ objektů $\varphi\wedge\neg\psi$, $c$ objektů
$\neg\varphi\wedge\psi$, $d$ objektů $\neg\varphi\wedge\neg\psi$, celkem $n=a+b+c+d$.

Dva běžné typy observačních pravidel:
\begin{itemize}
\item \textbf{založená na odhadech pravděpodobností} (případně podmíněných) ---
      nejběžnější jsou asociační pravidla (odd.~\ref{ssec:ar}) s~konfidencí
      $c_0\in(0,1)$ a~podporou $s_0\in(0,1)$:
      \[ (\varphi\sim\psi)=1\quad\text{právě když}\quad
         \frac{a}{a+b}\ge c_0\ \wedge\ \frac{a}{n}\ge s_0, \]
      kde $\tfrac{a}{a+b}$ je nestranný odhad $P(\psi\mid\varphi)$;
\item \textbf{založená na testech hypotéz} ($H_0$ nulová, $H_1$ alternativní) --- např.\
      $(\varphi\sim\psi)=1$, právě když jednostranný Fisherův test
      (odd.~\ref{ssec:kontingencni}) na hladině $\alpha\in(0;0{,}5\rangle$ zamítá $H_0$:
      nezávislost $\varphi$ a~$\psi$ ve prospěch $H_1$: pozitivní závislost $\varphi$
      a~$\psi$.
\end{itemize}
\end{defbox}

\subsection{Statistický pohled: regresní a~diskriminační analýza}
\label{ssec:stat-pohled}

Charakteristický a~diskriminační popis tříd nemusí mít podobu logických pravidel ---
klasická statistika řeší tytéž dvě úlohy \emph{regresní} a~\emph{diskriminační analýzou}.
Kontext: klasifikační metody se dělí podle toho, zda hranici mezi třídami primárně
\emph{hledají} (neuronové sítě, rozhodovací stromy, SVM --- moderní, výpočetně náročné),
nebo ji \emph{nehledají} a~vycházejí ze zkušenosti podobných případů či
z~pravděpodobností tříd ($k$-NN, bayesovská klasifikace, diskriminační analýza ---
tradiční, výpočetně nenáročné).

\textbf{Regresní analýza} určuje vztah veličiny $Y$ k~jedné či více veličinám
$X_1,\dots,X_n$: lineární regrese jedné proměnné metodou nejmenších čtverců,
vícerozměrná lineární regrese a~logistická regrese pro kategoriální výstup --- podrobně
odd.~\ref{ssec:regrese} a~\ref{ssec:logreg}. Regresní model je kvantitativní obdobou
\emph{charakteristického} popisu: říká, jaké hodnoty výstupu jsou pro dané vstupy
typické.

\textbf{Diskriminační analýza} (DA) klasifikuje vzory do \emph{známých} tříd, ale hranici
nehledá: každé třídě $C_j$ \emph{prokládá} pravděpodobnostní rozdělení z~předem zvolené
rodiny hustot $F=\{f(\cdot\mid\theta):\theta\in\Theta\}$ --- v~původní Fisherově DA
vícerozměrné normální rozdělení. Učení je maximálně věrohodný odhad parametrů $\theta_j$
z~učicích dat třídy $C_j$; nový vzor $x$ se klasifikuje podle
$\arg\max_j p_j\,f(x\mid\theta_j)$, kde $p_j$ je apriorní pravděpodobnost třídy.

\begin{defbox}[Lineární vs.\ kvadratická diskriminační analýza]
\textbf{LDA}: kovarianční matice $\Sigma$ je dána předem a~společná všem třídám;
diskriminační funkce
\[ g_j(x)=\ln p_j-\tfrac12\,(x-\mu_j)^{\!\top}\Sigma^{-1}(x-\mu_j), \]
odhaduje se jen $\mu_j$ ($q$ parametrů na třídu pro $x\in\mathbb{R}^q$); množina
$\{x: g_j(x)=g_k(x)\}$ mezi dvěma třídami je \emph{nadrovina}.

\textbf{QDA}: navíc vlastní kovarianční matice $\Sigma_j$ každé třídy (dalších
$\tfrac{q(q+1)}{2}$ parametrů na třídu); hranice mezi třídami je \emph{kvadratická
plocha}.
\end{defbox}

Pro dvě třídy lze místo $g_1,g_2$ hledat jedinou funkci $g=g_1-g_2$ a~klasifikovat podle
\emph{znaménka} --- podmínka $g(x)>0$ je kvantitativní obdobou \emph{diskriminačního
pravidla} (postačující podmínka příslušnosti k~první třídě). Odvození lineární
a~kvadratické diskriminační funkce z~normálních rozdělení viz odd.~\ref{ssec:da}.

\begin{center}
\small
\begin{tabular}{@{}lll@{}}
\toprule
\textbf{Úloha} & \textbf{Pravidlová podoba} & \textbf{Statistická podoba} \\
\midrule
charakterizace & charakteristické pravidlo ($\Rightarrow$, t-weight) & regresní analýza \\
diskriminace & diskriminační pravidlo ($\Leftarrow$, d-weight) & diskriminační analýza \\
\bottomrule
\end{tabular}
\end{center}

\subsection{Měření zajímavosti pravidel}
\label{ssec:zajimavost}

Extrakce typicky vydá víc pravidel, než kolik jich uživatel unese --- pravidla je třeba
měřit, třídit a~filtrovat. Pravidlo $\varphi\Rightarrow\psi$ je přitom samo binární
klasifikátor, takže se na ně vztahují míry kvality binární klasifikace
(odd.~\ref{ssec:eval}): přesnost (\emph{accuracy}) AC, správnost predikce
(\emph{precision}) PR, citlivost (\emph{recall}) TPr, specificita TNr $=1-{}$FPr, F-míra
$FM=\frac{2\,PR\cdot TPr}{PR+TPr}$ a~AUC. U~observačních pravidel jim odpovídají
\emph{konfidence} $\tfrac{a}{a+b}$ (precision pravidla) a~\emph{podpora} $\tfrac{a}{n}$;
u~charakteristických a~diskriminačních pravidel je d-weight precision pravidla
a~t-weight jeho recall vůči cílové třídě.

Role měr při práci s~pravidly:
\begin{itemize}
\item \emph{fitness} při evoluční konstrukci pravidel (odd.~\ref{ssec:evoluce});
\item \emph{třídění}: pravidla pro rodiny malwaru se řadí podle precision
      (odd.~\ref{ssec:aplikace-pravidel});
\item \emph{filtrace}: odstraňování redundantních a~nepřesných pravidel;
\item \emph{prahování}: observační pravidlo je \uv{zajímavé}, překročí-li konfidence
      a~podpora zvolené prahy, nebo zamítne-li test nezávislost
      (odd.~\ref{ssec:observacni}) --- každá volba míry a~prahu je jiný zobecněný
      kvantifikátor, a~tedy jiná množina pravidel nalezených v~datech.
\end{itemize}

\textbf{Různá cena různých chyb.} Tradiční přesnost a~chybovost předpokládají, že každá
chyba vadí stejně --- nerealisticky: nedoručený ham vadí víc než nerozpoznaný spam,
nerozpoznané napadení sítě víc než falešný poplach. Zobecněním chybovosti je
\emph{střední hodnota ceny} $\frac1n\sum_{i,j}c(i,j)\,n_{ij}$, kde $c(i,j)$ je cena
klasifikace třídy $i$ jako $j$ (tradiční chybovost: $c(i,i)=0$, jinak 1); zajímavost
pravidla se pak měří cenou chyb, kterých se dopouští.

\mimoq{subjektivní míry podle Han--Kamber} Vedle objektivních měr rozhodují
\emph{subjektivní}: srozumitelnost (celý smysl extrakce pravidel), jednoduchost (délka
konjunkce, počet pravidel), neočekávanost (rozpor s~očekáváním uživatele)
a~akceschopnost (lze na základě pravidla jednat). Prakticky: pravidlo je zajímavé, je-li
současně srozumitelné, platné na nových datech, potenciálně užitečné a~nové.

\subsection{Aplikace: spam, malware a~síťové hrozby, doporučování}
\label{ssec:aplikace-pravidel}

Tři internetové aplikace klasifikačních metod z~1.~přednášky kurzu ilustrují roli
pravidel:

\textbf{Filtrace spamu} je klasifikace \{kombinace hodnot příznaků\} $\to$ \{spam, ham\}:
příznaky z~obsahu zprávy (slova, kombinace slov, obrázky) a~z~metainformací (doména,
kódování, podvržená adresa odesilatele; anomálie průběhu SMTP; stejná zpráva mnoha
adresátům). Není-li dosažitelná současně nízká FP i~FN, pomáhá třetí třída mezi spamem
a~hamem --- karanténa/nejistota (např.\ SpamBayes). Pravidlové filtry pracují hlavně
s~metainformacemi; prostá booleovská disjunkce pravidel vede k~vysoké falešné pozitivitě
(\uv{grafické tělo zprávy $\Rightarrow$ automaticky spam}), lepší je Łukasiewiczova fuzzy
disjunkce pravidel se stejným konsekventem. FP (nedoručený ham) přitom vadí víc než FN.

\textbf{Detekce malwaru a~síťových hrozeb.} Podezřelé programy nakonec vyhodnocuje
analytik --- chce rozumět, \emph{proč} jsou podezřelé, proto pravidla. Vzniká jich mnoho,
takže se odstraňují redundantní a~nepřesná; typická je hierarchická klasifikace se dvěma
vrstvami pravidel (první přesná pro normální software a~citlivá pro malware, druhá třídí
rodiny malwaru, setříděná podle precision). U~síťových hrozeb (útoky R2L, U2R, DoS,
probe) jsou obdobou pravidel \emph{signatury} známých útoků --- vzory se stoprocentně
jistým závěrem.

\textbf{Doporučovací systémy.} Při kolaborativním filtrování podle chování uživatele se
používají pravidla \uv{uživatel prošel stránky $S_1,\dots,S_n$ $\Rightarrow$ doporuč
objekt $O$ (obecněji: nabídni další stránku $S_d$)}. Získávají se statistickou analýzou
clickstreamových dat --- odhady pravděpodobností a~testováním platnosti modelů
(hypotéz), z~jejichž výsledků se pravidla konstruují heuristicky; jde tedy o~observační
pravidla obou typů z~odd.~\ref{ssec:observacni}. Doporučovací systémy často vysvětlují,
\emph{proč} objekt doporučují --- vysvětlení zvyšuje srozumitelnost doporučení pro
uživatele.

\subsection{Shrnutí ke~zkoušce}

\begin{itemize}
\item \textbf{Proč pravidla}: klasifikátor $f\colon X\to\{C_1,\dots,C_m\}$ je
      matematický předpis, málo srozumitelný; pravidlo logiky
      ($\varphi\Rightarrow\psi$, $\varphi\Leftrightarrow\psi$) je srozumitelnější.
      Booleovská implikace pravdivá při $t(\varphi)=0$ nebo $t(\varphi)=t(\psi)=1$;
      fuzzy implikace $=$ reziduum t-normy, $t(\varphi)\le t(\psi)\Rightarrow$
      implikace pravdivá se stupněm 1.
\item \textbf{Charakteristické vs.\ diskriminační pravidlo} (Han--Kamber):
      charakterizace popisuje cílovou třídu ($\Rightarrow$, nutná podmínka,
      \textbf{t-weight} $=\frac{\mathrm{count}(q_a)}{\sum_i\mathrm{count}(q_i)}$);
      diskriminace ji odlišuje od kontrastních ($\Leftarrow$, postačující podmínka,
      \textbf{d-weight} $=$ podíl pokrytých objektů patřících do cílové třídy). Umět
      spočítat z~tabulky; d-weight vždy srovnat s~\emph{apriorním} podílem třídy.
\item \textbf{Extrakce ze stromu}: cesta kořen $\to$ list $=$ konjunkce podmínek, cesty
      do listů téže třídy $=$ disjunkce; učení dělením podle informačního zisku či
      Giniho indexu; prořezávání pravidla zkracuje; zobecněná (šikmá) dělení snižují
      srozumitelnost.
\item \textbf{Evoluce pravidel}: genetický algoritmus, fitness $=$ AC/PR/TPr/TNr/AUC;
      michiganský přístup (jedinec $=$ pravidlo) vs.\ pittsburghský (jedinec $=$ soubor
      pravidel).
\item \textbf{Observační pravidla} (Hájek--Havránek 1978, GUHA): $\varphi\sim\psi$ se
      zobecněným kvantifikátorem nad čtyřpolní tabulkou $(a,b,c,d)$; (1)~odhady
      pravděpodobností --- asociační pravidla, $\frac{a}{a+b}\ge c_0$ (konfidence)
      $\wedge$ $\frac{a}{n}\ge s_0$ (podpora); (2)~testy hypotéz --- jednostranný
      Fisherův test zamítá nezávislost ve prospěch pozitivní závislosti.
\item \textbf{Statisticky}: regresní analýza $=$ charakteristický popis vztahu $Y$ na
      $X_1,\dots,X_n$ (nejmenší čtverce, vícerozměrná, logistická); diskriminační
      analýza $=$ klasifikace do známých tříd prokládáním rozdělení: ML odhad
      $\theta_j$, predikce $\arg\max_j p_j f(x\mid\theta_j)$; \textbf{LDA} (společná
      $\Sigma$, hranice nadrovina, $q$ parametrů) vs.\ \textbf{QDA} (vlastní $\Sigma_j$,
      hranice kvadratická plocha, $+\frac{q(q+1)}{2}$ parametrů); dvě třídy $\to$
      znaménko $g=g_1-g_2$.
\item \textbf{Zajímavost}: pravidlo $=$ binární klasifikátor $\to$ AC, PR
      ($=$ konfidence), TPr (recall), TNr, F-míra, AUC; role: fitness, třídění
      (malware podle precision), filtrace redundantních, prahování (konfidence
      a~podpora $=$ zobecněný kvantifikátor); různé ceny chyb (střední hodnota ceny).
      Subjektivně: srozumitelnost, jednoduchost, neočekávanost, akceschopnost.
\item \textbf{Aplikace}: spam (metainformace, 3.~třída karanténa, Łukasiewiczova
      disjunkce proti FP), malware (dvouvrstvá hierarchie pravidel), síťové hrozby
      (signatury), doporučování (pravidla z~clickstreamu).
\end{itemize}
%=====================================================================
\newpage
\section{Reprezentace, vyhodnocování a~vizualizace znalostí}
\label{sec:reprez}
%=====================================================================

Tři části podle tří slov ve větě okruhu: \emph{v~jaké formě} znalost vyjádříme, \emph{jak
ověříme}, že je dobrá, a~\emph{jak ji zobrazíme}. Prostřední část --- vyhodnocování
klasifikátorů --- je ve slidech pokrytá celá a~ptá se na ni u~zkoušky nejčastěji.

\subsection{Formy reprezentace znalostí}

\mimoq{tabulka jako přehled; jednotlivé metody v~ní jsou ve slidech doloženy}
\begin{center}
\small
\begin{tabular}{@{}p{3.4cm}p{4.4cm}p{7.0cm}@{}}
\toprule
\textbf{Forma} & \textbf{Typický zdroj} & \textbf{Vlastnosti} \\
\midrule
IF-THEN pravidla & asociační pravidla, C4.5 & nejsrozumitelnější; lokální (každé pravidlo se
  čte samostatně); riziko konfliktů a~nadprodukce \\
Rozhodovací strom & ID3/C4.5/CART & globální model, snadno čitelný do hloubky 4--5; nad ni
  nepřehledný \\
Rozhodovací tabulka & diskretizovaná data & úplný výčet kombinací; přehledná jen pro málo
  atributů \\
Matematická funkce & regrese, SVM, NN & kompaktní a~přesná; obtížně interpretovatelná (kromě
  lineárního modelu s~normalizovanými vahami) \\
Prototypy / centroidy & $k$-means, $k$-medoids, LVQ & \uv{typický zástupce} skupiny; velmi
  intuitivní pro byznys \\
Pravděpodobnostní model & naivní Bayes, bayesovská síť & explicitní nejistota; graf sítě je
  čitelný \\
Grafy a~sítě & analýza sociálních sítí & vztahy mezi entitami; vizuálně silné \\
Instance (paměť případů) & $k$-NN, CBR & \uv{model} je samotná data; vysvětlení odkazem na
  podobné případy \\
\bottomrule
\end{tabular}
\end{center}

\begin{defbox}[Interpretovatelnost a~vysvětlitelnost]
\mimoq{vysvětlitelnost a~post-hoc metody; interpretovatelnost jako kritérium hodnocení modelu
je ve slidech --- viz odd.~\ref{ssec:eval}}
\emph{Interpretovatelnost} (\emph{interpretability}) --- schopnost člověka porozumět tomu, jak
model dospěl k~výstupu. \emph{Vysvětlitelnost} (\emph{explainability}) --- schopnost podat
srozumitelné odůvodnění konkrétního rozhodnutí.

Rozlišujeme \emph{modely inherentně interpretovatelné} (lineární modely, rozhodovací stromy,
sady pravidel, prototypy) a~\emph{post-hoc vysvětlení černých skříněk} --- důležitost atributů
(permutační, SHAP), lokální surogátní modely (LIME), grafy částečné závislosti (PDP), extrakce
pravidel z~neuronové sítě, analýza citlivosti, kontrafaktuální vysvětlení (\uv{kdyby byl váš
příjem o~5000 vyšší, úvěr by byl schválen}).

Interpretovatelnost není jen pohodlí: v~regulovaných doménách (úvěry, pojištění, medicína) je
právním požadavkem (GDPR čl.~22 --- právo na vysvětlení automatizovaného rozhodnutí; AI Act
pro vysoce rizikové systémy).
\end{defbox}

\subsection{Vyhodnocování modelů}
\label{ssec:eval}

\subsubsection{Kritéria hodnocení klasifikačních metod}

\begin{itemize}
\item \textbf{Prediktivní přesnost} (\emph{accuracy}) --- podíl správně klasifikovaných vzorů:
      \[ \mathit{Accuracy}=\frac{\text{počet správných klasifikací}}
         {\text{počet všech testovaných případů}}. \]
\item \textbf{Efektivita} --- čas potřebný pro konstrukci modelu a~čas pro jeho použití.
\item \textbf{Robustnost} --- zacházení se šumem a~chybějícími hodnotami.
\item \textbf{Škálovatelnost} --- efektivita pro databáze uložené na disku.
\item \textbf{Interpretovatelnost} --- jasnost a~vhled, který model poskytuje.
\item \textbf{Kompaktnost modelu} --- velikost stromu, počet pravidel, \dots
\end{itemize}

\subsubsection{Metody odhadu chyby}

\begin{defbox}[Holdout, křížová validace, leave-one-out, validační množina]
\textbf{Holdout}: $D$ se rozdělí na disjunktní \emph{trénovací} $D_{train}$ a~\emph{testovací}
$D_{test}$. \textbf{Zásadní pravidlo:} trénovací množina se nesmí použít při testování
a~testovací při učení --- jen neviděná testovací množina dá \emph{nezkreslený} odhad. Vhodné pro
\emph{velká} $D$.

\textbf{$k$-násobná křížová validace} (\emph{$k$-fold cross-validation}): data se uniformně
rozdělí na $k$ stejně velkých disjunktních podmnožin, každá se jednou použije jako testovací
a~zbylých $k-1$ jako trénovací. Z~$k$ přesností $\mathit{Accuracy}_i$ je výsledkem průměr
\[ \mathit{Accuracy}_{CV}=\frac1k\sum_{i=1}^{k}\mathit{Accuracy}_i. \]
Běžné je 5 a~10; metoda se používá spíše pro \emph{menší} data.

\textbf{Leave-one-out}: křížová validace, kde každý fold má jediný testovací vzor --- pro $m$
vzorů $m$-násobná CV. Pro \emph{velmi malá} data.

\textbf{Validační množina}: data se rozdělí na \emph{tři} části --- trénovací, validační
a~testovací. Validační slouží k~odhadu \emph{parametrů} učicích algoritmů (vezmou se hodnoty
s~nejlepší přesností na ní); místo ní lze použít i~křížovou validaci.
\end{defbox}

\begin{defbox}[Interval spolehlivosti pro křížovou validaci]
Přibližný $N\%$ interval spolehlivosti:
\[ \mathit{Accuracy}_{CV}\pm t_{N,k-1}\cdot s_{\mathit{Accuracy}},\qquad
   s_{\mathit{Accuracy}}=\sqrt{\frac{1}{k(k-1)}\sum_{i=1}^{k}
   \big(\mathit{Accuracy}_i-\mathit{Accuracy}_{CV}\big)^2}. \]
Hodnoty $t_{N,k-1}$ pro dvoustranné intervaly:
\begin{center}
\small
\begin{tabular}{@{}lrrr@{}}
\toprule
& \multicolumn{3}{c}{hladina spolehlivosti $N$} \\
$k$ & 90~\% & 95~\% & 99~\% \\
\midrule
2 & 2,92 & 4,30 & 9,92 \\
10 & 1,81 & 2,23 & 3,17 \\
30 & 1,70 & 2,04 & 2,75 \\
$\infty$ & 1,64 & 1,96 & 2,58 \\
\bottomrule
\end{tabular}
\end{center}
Často se používá $t=1{,}96$ (pro $k\to\infty$ a~$N=95\,\%$).
\end{defbox}

\subsubsection{Matice záměn a~odvozené míry}

Přesnost je jen \emph{jedna} míra ($\mathit{error}=1-\mathit{accuracy}$) a~v~některých
aplikacích není vhodná. Při dolování z~textu nás mohou zajímat jen dokumenty určitého tématu,
což je malá část velké kolekce. Při klasifikaci silně \emph{nevyvážených} dat (detekce
průniků do sítě, finanční podvody) nás zajímá jen minoritní třída: při 1~\% průniků dosáhneme
99\% přesnosti tím, že \emph{neuděláme nic} --- a~neodhalíme žádný průnik. Zajímavá třída se
nazývá \emph{pozitivní}, ostatní \emph{negativní}.

\begin{defbox}[Matice záměn, precision, recall, $F_1$]
\begin{center}
\small
\begin{tabular}{@{}llcc@{}}
\toprule
& & \multicolumn{2}{c}{\textbf{klasifikováno jako}} \\
& & pozitivní & negativní \\
\midrule
\textbf{skutečnost} & pozitivní & TP & FN \\
                    & negativní & FP & TN \\
\bottomrule
\end{tabular}
\end{center}
$TP$ --- správné klasifikace pozitivních vzorů (\emph{true positive}); $FN$ --- nesprávné
klasifikace pozitivních vzorů (\emph{false negative}); $FP$ --- nesprávné klasifikace
negativních vzorů (\emph{false positive}); $TN$ --- správné klasifikace negativních vzorů.

\[ p=\mathit{precision}=\frac{TP}{TP+FP},\qquad r=\mathit{recall}=\frac{TP}{TP+FN}. \]
\emph{Precision} je podíl správných mezi vzory \emph{označenými} za pozitivní, \emph{recall}
podíl nalezených mezi \emph{skutečně} pozitivními. \textbf{Obě míry hodnotí klasifikaci pouze na
pozitivní třídě.}

\textbf{$F_1$-míra} kombinuje obě do jedné --- je jejich \emph{harmonickým průměrem}:
\[ F_1=\frac{2pr}{p+r}=\frac{2}{\frac1p+\frac1r}. \]
Harmonický průměr dvou čísel je blíž tomu \emph{menšímu} z~nich, takže $F_1$ je velká jen
tehdy, jsou-li \emph{obě} míry velké.
\end{defbox}

\emph{Ilustrace pasti:} klasifikátor, který správně označí jediný pozitivní vzor a~žádný
negativní neoznačí chybně, má $p=100\,\%$ a~$r=1\,\%$. Vysoká precision sama o~sobě tedy
neříká nic o~užitečnosti.

\subsubsection{Skórování, řazení a~lift analýza}

\emph{Skórování} (\emph{scoring}) je příbuzné klasifikaci, zajímá nás ale jediná (pozitivní)
třída --- např.\ kupující v~marketingové databázi. Místo definitivního zařazení dostane každý
vzor \emph{odhad pravděpodobnosti} (\emph{probability estimate}, PE), že do pozitivní třídy
patří; klasifikátor jej musí umět dodat (u~rozhodovacího stromu lze použít confidence v~listu).
Vzory se pak podle PE \emph{seřadí} a~rozdělí do $n$ (např.\ 10) přihrádek (\emph{bins}); podle
počtu pozitivních vzorů v~přihrádkách se nakreslí \textbf{lift křivka} --- \emph{lift analýza}.

\begin{defbox}[Příklad --- lift analýza pro cílenou kampaň]
Chceme rozeslat propagační materiály potenciálním zákazníkům (hodinky). Každá zásilka stojí
\$0,50, prodej jedněch hodinek přinese \$5 zisku. Kolik zásilek poslat a~komu?

Testovací množina má 10\,000 vzorů, z~nichž 500 je pozitivních. Po oskórování a~seřazení do
10 přihrádek (po 1000 vzorech):

\begin{center}
\small
\begin{tabular}{@{}lrrrrrrrrrr@{}}
\toprule
Přihrádka & 1 & 2 & 3 & 4 & 5 & 6 & 7 & 8 & 9 & 10 \\
\midrule
Pozitivních & 210 & 120 & 60 & 40 & 22 & 18 & 12 & 7 & 6 & 5 \\
\% z~celku & 42 & 24 & 12 & 8 & 4,4 & 3,6 & 2,4 & 1,4 & 1,2 & 1,0 \\
Kumulativně & 42 & 66 & 78 & 86 & 90,4 & 94 & 96,4 & 97,8 & 99 & 100 \\
\bottomrule
\end{tabular}
\end{center}

První přihrádka obsahuje 42~\% všech pozitivních případů, první tři 78~\%. Náhodné rozeslání
by v~každé přihrádce dalo 10~\% --- rozdíl mezi lift křivkou a~touto diagonálou je přínos
modelu. Rozhodnutí \uv{kolik zásilek} pak vychází z~ekonomiky: zásilkou do první přihrádky
získáme $210\cdot\$5-1000\cdot\$0{,}50=\$550$, do desáté
$5\cdot\$5-1000\cdot\$0{,}50=-\$475$.
\end{defbox}

\subsubsection{ROC křivka a~AUC}

\begin{defbox}[ROC křivka (\emph{Receiver Operating Characteristic})]
Graf \emph{míry správně pozitivních} (TPR) proti \emph{míře falešně pozitivních} (FPR):
\[ \mathit{TPR}=\frac{TP}{TP+FN},\qquad \mathit{FPR}=\frac{FP}{FP+TN},\qquad
   \mathit{TNR}=\frac{TN}{FP+TN}. \]
$\mathit{TPR}$ je \emph{recall} pozitivní třídy, nazývaný též \emph{senzitivita};
$\mathit{FPR}$ je podíl skutečně negativních případů klasifikovaných jako pozitivní;
$\mathit{TNR}$ (recall negativní třídy) je \emph{specificita} a~platí
$\mathit{FPR}=1-\text{specificita}$.

Každá ROC křivka vede z~$(0,0)$ (vše negativní) do $(1,1)$ (vše pozitivní). \textbf{Hlavní
diagonála} odpovídá náhodnému hádání --- predikci pozitivní třídy s~fixní pravděpodobností, kdy
$\mathit{TPR}=\mathit{FPR}$.
\end{defbox}

\textbf{Konstrukce ROC křivky} z~klasifikace ve formě žebříčku testovacích případů:
\begin{enumerate}
\item Získej žebříček setříděním testovacích případů \emph{klesajícím} výstupem klasifikátoru
      (např.\ aposteriorní pravděpodobností).
\item Rozděl žebříček \emph{na každém} testovacím případě na dvě části a~považuj každý případ
      v~horní části za pozitivní a~každý v~dolní za negativní.
\item Pro každé takové rozdělení spočítej dvojici $(\mathit{FPR},\mathit{TPR})$. Je-li horní
      část prázdná, dostaneme bod $(0,0)$; je-li prázdná dolní část, bod $(1,1)$.
\item Spoj sousední body.
\end{enumerate}

\begin{defbox}[AUC --- plocha pod ROC křivkou]
Dva klasifikátory se mohou v~ROC prostoru \emph{křížit}: v~příkladu ze slidů je pro
$\mathit{FPR}<0{,}43$ lepší $C_1$ a~pro $\mathit{FPR}>0{,}43$ lepší $C_2$ --- jednoznačná
odpověď \uv{který je lepší} tedy neexistuje. Celkově lze použít \emph{plochu pod ROC křivkou}
(\emph{area under the curve}, AUC): je-li $\mathrm{AUC}(C_i)>\mathrm{AUC}(C_j)$, říkáme, že
$C_i$ je lepší než $C_j$. Perfektní klasifikátor má $\mathrm{AUC}=1$, náhodné hádání
$\mathrm{AUC}=0{,}5$.
\end{defbox}

\subsection{Vizualizační techniky}

\mimoq{tabulka jako přehled; krabicový graf, dendrogram a~ROC/lift křivka samy ve slidech jsou}
\begin{center}
\small
\begin{tabular}{@{}p{3.9cm}p{5.0cm}p{6.0cm}@{}}
\toprule
\textbf{Technika} & \textbf{K~čemu} & \textbf{Omezení / poznámka} \\
\midrule
Histogram & rozdělení jedné numerické veličiny & citlivý na volbu šířky košů \\
Krabicový graf (\emph{box plot}) & medián, kvartily, odlehlé hodnoty; srovnání skupin &
  neukáže bimodalitu (řeší \emph{violin plot}) \\
Bodový graf (\emph{scatter plot}) & vztah dvou numerických veličin & při mnoha bodech
  překryv $\to$ průhlednost, hexbin \\
Matice bodových grafů (SPLOM) & všechny dvojice atributů najednou & použitelné do zhruba
  10 atributů \\
Paralelní souřadnice & vysokorozměrná data --- objekt je lomená čára přes svislé osy
  atributů & silně závisí na pořadí os; nutné škálování; křížení čar $=$ negativní korelace \\
Teplotní mapa (\emph{heatmap}) & korelační matice, matice záměn & volba barevné škály
  (divergentní pro korelace) \\
Dendrogram & výsledek hierarchického shlukování & pořadí listů není jednoznačné \\
Projekce (PCA, MDS, t-SNE, UMAP) & 2-D pohled na vysokorozměrná data & t-SNE/UMAP zkreslují
  globální vzdálenosti --- \emph{velikosti} shluků a~mezery mezi nimi nelze interpretovat \\
ROC / lift křivka & kvalita klasifikátoru & viz odd.~\ref{ssec:eval} \\
Mozaikový graf & kontingenční tabulka kategoriálních atributů & vizuální alternativa
  k~$\chi^2$ \\
Vizualizace sítě & sociální sítě, grafy vztahů & \uv{chlupatá koule} u~velkých grafů $\to$
  agregace komunit \\
\bottomrule
\end{tabular}
\end{center}

\textbf{Tři role vizualizace v~KDD:} \emph{explorativní analýza} (fáze porozumění datům) ---
odhalí rozdělení, odlehlé a~chybějící hodnoty, nečekané vztahy; Anscombeho kvartet jsou čtyři
datové sady s~identickými průměry, rozptyly i~korelacemi, ale zcela odlišnou strukturou, kterou
odhalí jedině graf; \emph{vizualizace modelu} --- strom, dendrogram, síť, hranice rozhodování;
\emph{prezentace výsledků} --- agregace, jednoduchost, jeden graf $=$ jedno sdělení.

\textbf{Zásady:} vysoký poměr datové a~inkoustové složky (Tufte); nezkreslovat (nulová osa
u~sloupcových grafů, lineární nebo zřetelně označené logaritmické osy); barva podle typu dat
(kategorie $\to$ kvalitativní paleta, pořadí $\to$ sekvenční, odchylka od středu $\to$
divergentní) a~s~ohledem na barvosleposti; vždy popisky a~jednotky; žádné 3-D efekty ani
koláčové grafy s~mnoha výsečemi.

\subsection{Vyhodnocení znalostí z~pohledu uživatele}

Pojem \uv{vyhodnocování získaných znalostí} má tři různé roviny, které je třeba rozlišovat:
\begin{itemize}
\item \emph{statistické vyhodnocení modelu} --- přesnost, matice záměn, ROC/AUC, odhad chyby
      křížovou validací (odd.~\ref{ssec:eval});
\item \emph{vyhodnocení zajímavosti jednotlivých vzorů} --- support, confidence, lift,
      $\chi^2$, t-weight a~d-weight (odd.~\ref{ssec:ar} a~kap.~\ref{sec:pravidla});
\item \emph{vyhodnocení z~pohledu uživatele} --- zda je nalezená znalost srozumitelná,
      platná, užitečná a~nová.
\end{itemize}
První dvě roviny jsou objektivní a~měřitelné, třetí je nutně subjektivní --- tentýž vzor může
být pro jednoho objevem a~pro druhého notorietou. Část výsledků bude zřejmá, část použitelná
přímo a~část je nutné prezentovat srozumitelněji; obvykle se shrnou do \emph{analytické
zprávy}, výstupem může být i~provedení akce.

\subsection{Shrnutí ke~zkoušce}

\begin{itemize}
\item \textbf{Formy reprezentace znalostí}: IF-THEN pravidla (nejsrozumitelnější, lokální),
      rozhodovací strom (globální, čitelný do hloubky 4--5), rozhodovací tabulka, matematická
      funkce (kompaktní, neinterpretovatelná), prototypy/centroidy, pravděpodobnostní model,
      grafy a~sítě, instance ($k$-NN, CBR). Umět ke~každé formě říct, z~jaké metody vzniká
      a~čím se platí za její srozumitelnost.
\item \textbf{Kritéria hodnocení metod}: prediktivní přesnost, efektivita (čas konstrukce
      i~použití), robustnost, škálovatelnost, interpretovatelnost, kompaktnost modelu.
\item \textbf{Odhad chyby}: \emph{holdout} (velká data; trénovací a~testovací množina se
      nesmí míchat), \emph{$k$-násobná křížová validace} (menší data; 5 a~10 je běžné;
      přesnost jako průměr $k$ přesností $+$ interval spolehlivosti s~$t_{N,k-1}$),
      \emph{leave-one-out} (velmi malá data, $m$-násobná CV), \emph{validační množina} (třetí
      část dat, na ní se ladí \emph{parametry} algoritmu).
\item \textbf{Matice záměn}: $TP$, $FN$, $FP$, $TN$; $p=\frac{TP}{TP+FP}$,
      $r=\frac{TP}{TP+FN}$, $F_1=\frac{2pr}{p+r}$ (harmonický průměr --- blíž menšímu
      z~čísel). Umět vysvětlit, \emph{proč} accuracy nestačí u~nevyvážených dat (1~\% průniků
      $\Rightarrow$ 99~\% přesnosti \uv{nedělat nic}).
\item \textbf{ROC}: $\mathit{TPR}=\frac{TP}{TP+FN}$ (senzitivita $=$ recall pozitivní třídy)
      proti $\mathit{FPR}=\frac{FP}{FP+TN}$; $\mathit{TNR}$ $=$ specificita
      $=1-\mathit{FPR}$. Křivka vede z~$(0,0)$ do $(1,1)$, diagonála $=$ náhodné hádání.
      Umět popsat konstrukci ze žebříčku (dělení na každém případu). \textbf{AUC}: 1 pro
      perfektní, 0,5 pro náhodný klasifikátor; potřebná proto, že se ROC křivky dvou
      klasifikátorů mohou \emph{křížit}.
\item \textbf{Skórování a~lift}: místo třídy se přiřadí odhad pravděpodobnosti (u~stromu
      confidence v~listu), vzory se seřadí a~rozdělí do přihrádek; lift křivka proti
      diagonále náhodného výběru. Umět provést ekonomickou úvahu (kolik zásilek se vyplatí).
\item \textbf{Tři roviny \uv{vyhodnocování}}: statistická kvalita modelu, zajímavost vzorů,
      vyhodnocení uživatelem. První dvě objektivní, třetí subjektivní. \textbf{Čtyři kritéria
      zajímavosti znalosti}: srozumitelnost, platnost, užitečnost, novost.
\item \textbf{Vizualizace}: technika podle úlohy (histogram, box plot, scatter, SPLOM,
      paralelní souřadnice, heatmapa, dendrogram, projekce, ROC/lift); tři role v~KDD
      (explorace --- Anscombeho kvartet, vizualizace modelu, prezentace); zásady (Tufte,
      nezkreslovat osami, barva podle typu dat). Znát omezení: t-SNE a~UMAP zkreslují
      globální vzdálenosti.
\end{itemize}

%=====================================================================
\newpage
\section{Modely pro analýzu sociálních sítí, centralita a~detekce komunit}
\label{sec:sna}
%=====================================================================

Okruh jmenuje \emph{modely}, \emph{míry centrality} a~\emph{detekci komunit}; podle toho je
kapitola členěna: reprezentace sítí a~míry na úrovni aktéra (centralita) a~sítě (koheze), pak
modely (vlastnosti reálných sítí, šíření vlivu, homofilie, segregace, strukturní balance,
analýza vazeb --- HITS a~PageRank) a~nakonec detekce komunit a~predikce vazeb.

\subsection{Vymezení oboru a~motivace}

\begin{defbox}[Analýza sociálních sítí]
\emph{Analýza sociálních sítí} (\emph{social network analysis}, SNA) je interdisciplinární
obor s~původem ve třech oblastech:
\begin{itemize}
\item \emph{sociální vědy} --- otázky a~problémy vysvětlující sociální chování analýzou
      společností;
\item \emph{analýza sítí} --- formulace a~řešení problémů, které mají síťovou strukturu;
\item \emph{teorie grafů} --- matematický aparát pro reprezentaci a~analýzu grafů.
\end{itemize}
Zásadní posun perspektivy: SNA se soustředí na \textbf{vztahy} mezi sociálními entitami
(jednotlivci, skupinami, institucemi) \emph{spíše než na entity samotné} a~jejich atributy.
Studium společnosti ze síťové perspektivy chápe jednotlivce jako \emph{zasazené} do sítě
vztahů a~hledá vysvětlení sociálního chování ve struktuře těchto sítí, nikoli v~jednotlivcích
samotných.
\end{defbox}

\textbf{Analýza vazeb} (\emph{link analysis}) jako motivace: najít a~vizualizovat vztahy mezi
daty. Aplikace: telekomunikace; kriminalistika a~právo (propojené skupiny zločinců); marketing
(vztahy mezi zákazníky).

\textbf{Šest stupňů odloučení} (Milgram, 1967): náhodně vybraní lidé z~Nebrasky měli poslat
dopis burzovnímu makléři v~Bostonu, a~to jen přes osobně známé prostředníky. Mezi dopisy, které
cíl našly, byl průměrný počet vazeb \emph{šest} (mnoho dopisů ale nedorazilo). Experiment
předznamenal pozdější objevy --- huby, tendenci postupovat podle geografie či profese. Totéž
dokládají \emph{Erdősova čísla}: Erdős napsal přes 1500 článků s~507 spoluautory, průměrné
Erdősovo číslo mezi matematiky je 4,65, maximum 13.

\textbf{Bezškálové sítě} (\emph{scale-free}, SF) jako druhá motivace: několik uzlů má extrémně
mnoho spojení (\emph{huby}), většina jen několik. Důsledkem je \emph{odolnost proti náhodné
poruše} a~zároveň \emph{zranitelnost koordinovaným útokem}: může selhat až 80~\% náhodných uzlů
bez fragmentace, ale eliminace 5--15~\% \emph{hubů} rozpadne celý systém. Aplikace: ochrana
proti šíření virů internetem (vymýcení je v~SF architektuře v~podstatě nemožné), medicína
(vakcinace cílená na huby, léky mířící na klíčové molekuly), business (kaskádové finanční
krachy, influencer marketing). Příklady SF sítí: \emph{sociální} (spoluautorství, Hollywood),
\emph{biologické} (metabolismus buňky, proteinové regulační sítě), \emph{sociotechnické}
(internet --- routery a~spoje; WWW --- stránky a~URL).

\subsection{Reprezentace sítí}

\begin{defbox}[Základní reprezentace]
Síť je graf $G=(V,E)$: \emph{vrcholy} (uzly) jsou aktéři, \emph{hrany} (vazby) vztahy mezi
nimi. Tři ekvivalentní zápisy:
\begin{itemize}
\item \textbf{obrázek grafu};
\item \textbf{seznam hran} (\emph{edge list}) --- dvojice vrcholů, u~ohodnocených grafů plus
      sloupec vah;
\item \textbf{matice sousednosti} (\emph{adjacency matrix}) $A=(a_{ij})$, $a_{ij}=1$, je-li
      vrchol $i$ spojen s~vrcholem $j$, jinak 0; u~vážených grafů obsahuje váhy.
\end{itemize}
U~\textbf{orientovaných} grafů (kdo koho kontaktuje) matice sousednosti \emph{nemusí} být
symetrická; u~\textbf{neorientovaných} (kdo koho zná) symetrická je. Tentýž seznam hran má
v~obou případech jinou interpretaci.
\end{defbox}

\textbf{Váhy hran} mohou vyjadřovat frekvenci interakce v~období pozorování, počet vyměněných
předmětů, individuální percepci silné vazby, náklady komunikace či výměny (vzdálenost), nebo
jejich kombinaci. Hrany samy reprezentují interakce, toky informací či zboží, podobnosti
a~afiliace, nebo sociální vztahy.

\textbf{Ego sítě a~\uv{celé} sítě.} \emph{Ego síť} je osobní síť jednoho aktéra (\emph{ego})
a~jeho sousedů (\emph{alterů}); pro další analýzu ego \emph{není} potřebné, protože všichni
alteři jsou s~ním spojeni z~definice. \uv{Celá} síť obsahuje i~izolované uzly
(\emph{isolates}), ale \textbf{žádná studovaná síť není v~praxi \uv{celá}} --- je to jen
částečný obraz (\emph{problém specifikace hranic}, \emph{boundary specification problem}).

\subsection{Síla vazeb}

Uživatel může mít online tisíce \uv{přátel}, ale vazby nejsou stejně silné a~hrají různou roli
při formování komunit a~difuzi informací. \emph{Silné vazby} jsou blízcí přátelé, \emph{slabé}
známí. \textbf{Síla slabých vazeb} (Granovetter, 1973): občasná setkání se vzdálenými známými
mohou poskytnout důležitou informaci o~nových příležitostech (hledání práce).

\begin{defbox}[Homofilie, tranzitivita, přemostění]
\textbf{Homofilie} (\emph{homophily}) --- tendence vztahovat se k~lidem s~podobnými
charakteristikami (status, přesvědčení, \dots). Vede k~tvorbě homogenních skupin (klastrů),
kde je snazší vytvářet vztahy; \emph{extrémní} homogenizace ale může působit proti inovaci
a~generování idejí, takže \emph{heterofilie} je v~některých kontextech žádoucí. Homofilní
vazby mohou být silné i~slabé.

\textbf{Tranzitivita} (\emph{transitivity}) --- \uv{existuje-li vazba mezi $A$ a~$B$ a~mezi
$B$ a~$C$, budou v~tranzitivní síti spojeni i~$A$ a~$C$.} Silné vazby jsou tranzitivní
častěji než slabé, takže tranzitivita \emph{indikuje} existenci silných vazeb (není to ale
nutná ani postačující podmínka). Tranzitivita a~homofilie společně vedou ke vzniku
\emph{klik} (plně propojených klastrů).

\textbf{Přemostění} (\emph{bridging}) --- \emph{most} (\emph{bridge}) je uzel či hrana
spojující různé skupiny; obvykle jde o~slabé vazby, ale ne každá slabá vazba je most. Mosty
usnadňují mezigrupovou komunikaci, zvyšují sociální kohezi a~podporují inovaci.
\end{defbox}

\textbf{Určování síly vazby z~topologie sítě.} Mosty spojující dvě různé komunity jsou slabé
vazby; hrana je \emph{most}, jestliže její odstranění rozpojí její koncové uzly. Mosty jsou
v~reálných sítích poměrně vzácné $\Rightarrow$ definice se \emph{uvolní} na \textbf{zkratkový
most} (\emph{shortcut bridge}): kontroluje se, zda se po odstranění hrany \emph{zvětší}
vzdálenost jejích koncových uzlů (na více než 2). Čím větší vzdálenost, tím \emph{slabší}
vazba. Např.\ $d(2,5)=4$ po odstranění hrany $(2,5)$, ale $d(5,6)=2$ po odstranění $(5,6)$
$\Rightarrow$ $(5,6)$ je silnější vazba než $(2,5)$.

\begin{defbox}[Překryv sousedství (\emph{neighborhood overlap})]
Pro dva uzly $i,j$ spojené vazbou --- čím větší překryv, tím \emph{silnější} vazba:
\[ O(i,j)=\frac{\text{počet společných přátel } i \text{ a } j}
   {\text{počet přátel, kteří jsou sousedy alespoň }i\text{ nebo }j}
   =\frac{|N_i\cap N_j|}{|N_i\cup N_j|-2}, \]
kde $-2$ ve jmenovateli vylučuje $i$ a~$j$ samotné.

\emph{Příklady:} $O(2,5)=0$; $O(5,6)=\frac{|\{4\}|}{|\{2,4,5,6,7,10\}|-2}=\frac14$;
pro graf $G=(\{1,2\},\{(1,2)\})$ je $O(1,2)=0$.
\end{defbox}

\textbf{Určování síly vazby z~profilů a~interakcí.} Na Twitteru je skutečná síť přátelství
určena \emph{frekvencí}, s~jakou spolu dva uživatelé mluví, nikoli sítí follower--followee ---
a~právě tato skutečná síť ovlivňuje používání Twitteru silněji. Na Facebooku lze sílu vazeb
predikovat z~příspěvků iniciovaných přítelem, zpráv na zdi a~počtu společných přátel.
Numerickou sílu lze učit maximální věrohodností: podobnost profilů určuje sílu vazby a~ta
určuje interakci, takže se maximalizuje věrohodnost pozorovaných profilů a~interakcí.

\subsection{Klíčoví hráči: míry centrality}

\begin{defbox}[Stupňová centralita (\emph{degree centrality})]
Vstupní a~výstupní stupeň uzlu $=$ počet vazeb do uzlu / z~uzlu (v~neorientovaných grafech
shodné). Měří \emph{propojenost}, \emph{vliv} či \emph{popularitu} --- které uzly jsou
centrální pro šíření informace a~ovlivňování ostatních ve svém \uv{sousedství}.
\end{defbox}

\begin{defbox}[Cesty a~nejkratší cesty]
\emph{Cesta} mezi dvěma uzly je libovolná posloupnost neopakujících se uzlů, která je spojuje.
\emph{Nejkratší cesta} spojuje dva uzly nejmenším počtem hran (\emph{vzdálenost} uzlů).
Kratší cesty jsou preferované pro rychlou komunikaci či výměnu --- ale ne vždy, například při
šíření nemoci.

\textbf{Algoritmy nejkratších cest.} \emph{Jeden zdroj:} obecné grafy (orientované, záporné
váhy dovoleny) --- Bellmanův--Fordův algoritmus, $O(|V||E|)$; grafy bez záporných vah ---
Dijkstrův algoritmus, $O(|V|^2)$; hledá-li se cesta jen mezi dvěma vrcholy, lze algoritmus
zastavit po jejím nalezení. \emph{Všechny páry:} záporné váhy dovoleny ---
Floydův--Warshallův, $O(|V|^3)$; řidší grafy se zápornými vahami --- Johnsonův,
$O(|E||V|+|V|^2\log_2|V|)$; bez záporných vah --- Dijkstra, $O(|V|^3)$.
\end{defbox}

\begin{defbox}[Mezilehlá centralita (\emph{betweenness centrality})]
Pro uzel $v$: pro každou dvojici $s,t$ se počet nejkratších cest mezi $s$ a~$t$ procházejících
$v$ vydělí všemi nejkratšími cestami mezi $s$ a~$t$ a~hodnoty se sečtou přes všechny dvojice
(někdy se normalizuje tak, aby maximum bylo 1). Indikuje uzly na \emph{komunikačních cestách}
mezi jinými uzly --- tedy místa, kde by se síť mohla \emph{rozpadnout}.
\end{defbox}

\begin{defbox}[Blízkostní centralita (\emph{closeness centrality})]
Míra \emph{dosahu} --- rychlost, s~jakou se informace z~daného uzlu dostane k~ostatním:
spočítá se \emph{průměrná} délka všech nejkratších cest z~uzlu ke všem ostatním uzlům sítě
(kolik \uv{skoků} je v~průměru potřeba), a~vezme se \emph{převrácená hodnota}. Vyšší hodnoty
znamenají vyšší blízkost (jsou \uv{lepší}). Užitečná, jde-li především o~rychlost šíření.
\end{defbox}

\begin{defbox}[Vlastní centralita (\emph{eigenvector centrality})]
Je proporcionální součtu vlastních centralit přímých sousedů: moc a~status aktéra jsou
\emph{rekurzivně} definovány mocí a~statusem jeho sousedů. Měří tedy, jak dobře je aktér spojen
s~jinými \emph{dobře propojenými} aktéry.

Formálně: pro $G=(V,E)$ s~maticí sousednosti $A=(a_{i,j})$
\[ x_i=\frac1\lambda\sum_{j\in M(i)}x_j=\frac1\lambda\sum_{j\in V}a_{i,j}x_j
   \quad\Longleftrightarrow\quad Ax=\lambda x, \]
kde $M(i)$ je množina sousedů $i$. Existuje mnoho vlastních čísel $\lambda$, pro něž
nenulové řešení existuje; \emph{požadavek nezápornosti} všech složek vlastního vektoru však
znamená, že žádanou míru centrality dává \emph{jen největší} vlastní číslo. Vlastní vektor je
určen až na společný násobek, takže je nutná normalizace (např.\ součet přes všechny vrcholy
$=1$ nebo $=|V|$).

Blízce souvisí s~Google PageRank (odkazy ze silně odkazovaných stránek váží více,
odd.~\ref{ssec:pagerank}).
\end{defbox}

\begin{center}
\small
\begin{tabular}{@{}p{2.7cm}p{12.6cm}@{}}
\toprule
\textbf{Míra} & \textbf{Interpretace v~sociálních sítích} \\
\midrule
Stupňová & Kolik lidí může tato osoba dosáhnout \emph{přímo}? --- \emph{Síť hudebních
  spoluprací:} s~kolika lidmi spolupracoval? \\
Mezilehlá & Jak pravděpodobně leží tato osoba na nejpřímější cestě v~síti? --- \emph{Síť
  špionů:} přes koho nejspíš protéká většina informací? \\
Blízkostní & Jak rychle může tato osoba dosáhnout všech v~síti? --- \emph{Síť sexuálních
  vztahů:} jak rychle se do zbytku sítě rozšíří přenosná nemoc? \\
Vlastní & Jak dobře je tato osoba spojena s~jinými dobře propojenými lidmi? --- \emph{Citační
  síť:} který autor je nejvíce citován jinými dobře citovanými autory? \\
\bottomrule
\end{tabular}
\end{center}

\textbf{Množiny klíčových hráčů.} Jednotlivé míry nestačí. V~příkladu ze slidů je uzel 10
nejcentrálnější podle stupňové centrality, ale uzly 3 a~5 \emph{společně} dosáhnou více uzlů
a~vazba mezi nimi je \emph{kritická} --- při jejím přerušení se síť rozpadne na dvě izolované
podsítě. Uzly 3 a~5 společně jsou tedy pro síť důležitější než uzel 10 sám.

\begin{defbox}[Strukturní ekvivalence]
Vyjadřuje \emph{podobnost} mezi aktéry v~sociální síti; je založena na sousedech, které
sdílejí (ekvivalentně na počtu identických vazeb). Dva aktéři jsou \emph{strukturně
ekvivalentní}, lze-li jejich pozice vzájemně vyměnit \emph{bez změny celkové struktury} sítě.
Přibližnou strukturní ekvivalenci lze použít pro hierarchické shlukování sociálních sítí.
\end{defbox}

\subsection{Koheze: míry na úrovni sítě}

\begin{defbox}[Reciprocita, hustota, klastrovací koeficient]
\textbf{Reciprocita} (\emph{reciprocity}) --- míra vzájemnosti a~reciproční výměny v~síti,
související se sociální kohezí: podíl počtu \emph{reciprokých} vztahů (hrana v~obou směrech)
na celkovém počtu vztahů v~síti, kde dva vrcholy jsou \uv{ve vztahu}, existuje-li mezi nimi
alespoň jedna hrana. \textbf{Má smysl jen v~orientovaných grafech.}

\textbf{Hustota} (\emph{density}) --- podíl počtu hran v~síti na celkovém počtu \emph{možných}
hran:
\[ \text{density}=\frac{|E|}{|V|(|V|-1)/2}\quad\text{(neorientovaný graf)},\qquad
   \frac{|E|}{|V|(|V|-1)}\quad\text{(orientovaný)}. \]
Běžná míra toho, jak dobře je síť propojená (jak \uv{těsně upletená} je). Perfektně propojená
síť se nazývá \emph{klika} a~má hustotu 1. Orientované grafy mají poloviční hustotu než jejich
neorientované protějšky, protože možných hran je dvakrát více.

\textbf{Klastrovací koeficient} (\emph{clustering coefficient}) uzlu $i$ měří hustotu jeho
\emph{sousedství} --- sítě tvořené pouze uzly přímo spojenými s~$i$. Pro $G=(V,E)$ je
sousedství $N_i=\{v_j: e_{ij}\in E\ \vee\ e_{ji}\in E\}$, $k_i=|N_i|$, a
\[ C_i=\frac{2\,|\{e_{jk}: v_j,v_k\in N_i,\ e_{jk}\in E\}|}{k_i(k_i-1)}
   \ \text{(neorientovaný)},\qquad
   C_i=\frac{|\{e_{jk}\}|}{k_i(k_i-1)}\ \text{(orientovaný)}. \]
\emph{Klastrovací koeficient sítě} $C$ je průměr koeficientů všech uzlů,
$C=\frac{1}{|V|}\sum_{i=1}^{|V|}C_i$. Indikuje přítomnost (sub)komunit; shlukovací algoritmy
se snaží maximalizovat počet hran padnoucích do téhož klastru.
\end{defbox}

\begin{defbox}[Průměr, průměrná vzdálenost, excentricita, malé světy]
\textbf{Průměr} sítě (\emph{diameter}) --- nejdelší nejkratší cesta (vzdálenost) mezi
libovolnými dvěma uzly. Užitečná míra \emph{dosahu} sítě: udává, jak dlouho nejdéle trvá
dosáhnout libovolného uzlu. Řidší sítě mají obecně větší průměr.

\textbf{Průměrná vzdálenost} --- průměr všech nejkratších cest v~síti; udává, jak daleko od
sebe budou dva uzly v~průměru.

\textbf{Excentricita} vrcholu --- největší geodetická vzdálenost mezi daným vrcholem
a~libovolným jiným vrcholem grafu.

\textbf{Malý svět} (\emph{small world}) --- síť, která vypadá téměř náhodně, ale má
\emph{významně vysoký klastrovací koeficient} a~\emph{relativně krátkou průměrnou délku
cesty}: uzly se lokálně klastrují a~zároveň jsou dosažitelné v~několika krocích. V~sociálních
sítích je to velmi častý jev, a~to kvůli \emph{tranzitivitě} silných vazeb a~\emph{schopnosti
slabých vazeb dosáhnout přes klastry}. Sítě malého světa mají mnoho klastrů, ale i~mnoho mostů
mezi nimi, které zkracují průměrnou vzdálenost.
\end{defbox}

\subsection{Model reálných komplexních sítí a~jejich vlastnosti}

\begin{defbox}[Šest charakteristických vlastností reálných komplexních sítí]
Reálné komplexní sítě (sociální, informační a~znalostní, technologické, biologické) jsou
\emph{nenáhodné a~neregulární} grafy s~unikátními rysy:
\begin{enumerate}
\item \textbf{Efekt malého světa} --- první na něj poukázal Milgram (experiment z~60.~let, cca
      300 účastníků, mediánová délka úspěšných cest 6). Vychází z~\uv{šesti stupňů odloučení}:
      každý je od každého v~průměru šest kroků. Důsledky pro rychlost šíření informací i~nemocí.
\item \textbf{Tranzitivita / klastrování} --- vysoká pravděpodobnost \emph{úplných podsítí}
      uvnitř větších sítí (každý zná každého), zejména v~sociálních sítích.
\item \textbf{Mocninné rozdělení stupňů} --- \emph{rozdělení stupňů} je histogram podílu uzlů
      se stupněm $k$; u~reálných sítí je zcela odlišné od náhodných grafů.
\item \textbf{Odolnost sítě} (\emph{resilience}) --- vliv odstranění uzlů na propojenost.
      Většina sítí je \emph{robustní} proti náhodnému odstranění, ale podstatně méně proti
      \emph{cílenému} odstranění uzlů s~nejvyšším stupněm: odstraněním \uv{vrátných}
      (\emph{gatekeepers}) se výrazně změní komunikační schopnost sítě, protože se některé uzly
      odpojí. Jediný uzel strukturu obvykle neohrozí --- testuje se odstraněním \emph{procenta}
      uzlů.
\item \textbf{Vzorce míchání} (\emph{mixing patterns}) --- při koexistenci různých typů uzlů je
      běžná \emph{selektivita} ve tvorbě spojení, tzv.\ \emph{asortativní míchání}
      (\emph{assortative mixing}) čili homofilie (klasicky míchání podle rasy). Reálné sítě
      k~němu mají vyšší tendenci.
\item \textbf{Komunitní struktura} --- vzniká z~\emph{globální i~lokální heterogenity} rozložení
      hran: v~některých oblastech grafu je vysoká koncentrace hran (\emph{komunity}), mezi nimi
      nízká.
\end{enumerate}
\end{defbox}

\begin{defbox}[Bezškálové sítě a~preferenční připojování]
\textbf{Bezškálové sítě} (Barabási, Albert, 1999) mají rozdělení stupňů podle \emph{mocninného
zákona} (WWW, citační a~biologické sítě, některé sociální). Mocninná rozdělení vznikají tehdy,
\emph{závisí-li množství, které něčeho dostanete, na množství, které už máte}; topologicky to
dá síť s~několika silně propojenými uzly (\emph{huby}) a~velkým počtem uzlů s~nízkým stupněm.

Mechanismus se nazývá \emph{rich-get-richer / poor-get-poorer}, ekvivalentně \emph{Matoušův
efekt}; Price jej nazval \emph{kumulativní výhoda}, Barabási a~Albert \textbf{preferenční
připojování} (\emph{preferential attachment}). Nové uzly se připojují k~dobře propojeným uzlům,
které bývají spojené s~centrálními a~prestižními pozicemi (více statusu, popularity, znalostí,
peněz).

Dva mechanismy vzniku SF sítě jsou tedy \textbf{růst} a~\textbf{preferenční připojování}.
Výsledná síť má \emph{dlouhý chvost} rozdělení stupňů a~\emph{strukturu malého světa};
tranzitivita a~vlastnosti silných/slabých vazeb ke struktuře malého světa vést mohou, nejsou
k~jejímu vysvětlení ale \emph{nutné}.

\textbf{Důvody preferenčního připojování:} \emph{popularita} (\uv{bohatí bohatnou}) --- chceme
být spojeni s~populárními lidmi a~věcmi, což jejich popularitu dále zvyšuje, bez ohledu na
objektivní charakteristiky; \emph{kvalita} (\uv{dobří se zlepšují}) --- hodnotíme podle
objektivních kritérií, takže kvalitnější uzly přitáhnou pozornost rychleji; \emph{smíšený
model} --- kdo se stane hvězdou, nelze předpovědět, i~když kvalita hraje roli: mezi uzly
s~podobnými atributy se \uv{hvězdami} stanou ty, které první dosáhnou kritického množství
(\emph{halo efekt}).
\end{defbox}

\begin{defbox}[Struktura jádro--periferie a~centralizace]
\textbf{Centralizace} měří, jak je síť centralizovaná; vychází z~\emph{rozdílů ve stupních}
mezi uzly a~normalizuje se do $[0,1]$. Síť závislá na 1--2 vysoce propojených uzlech (důsledek
preferenčního připojování) má větší rozdíly ve stupňové centralitě. Centralizované struktury
mohou být lepší v~úlohách vyžadujících koordinaci (týmové řešení problémů), ale jsou
\emph{náchylnější k~selhání}, odpojí-li se klíčoví hráči.

Mnoho velkých komunit má \emph{jádro} hustě propojených uživatelů, kritické pro propojení
mnohem větší \emph{periferie}. Jádra lze identifikovat vizuálně, zkoumáním polohy uzlů
s~vysokým stupněm a~jejich společných rozdělení stupňů, nebo \emph{bow-tie analýzou} (použitou
pro analýzu struktury webu).
\end{defbox}

\subsection{Modelování vlivu}

\begin{defbox}[Společné vlastnosti metod modelování vlivu]
\begin{itemize}
\item sociální síť je reprezentována \emph{orientovaným} grafem, každý aktér je jeden uzel;
\item každý uzel začíná jako \emph{aktivní} nebo \emph{neaktivní};
\item aktivovaný uzel aktivuje své sousedy;
\item \textbf{jednou aktivovaný uzel už nelze deaktivovat.}
\end{itemize}
\end{defbox}

\begin{defbox}[Model lineárního prahu (LTM)]
Aktér (uzel) podnikne akci, pokud počet jeho přátel, kteří akci podnikli, překročí určitý
\emph{prah}.
\begin{itemize}
\item Každý uzel $v$ si zvolí prah $\theta_v$ \emph{náhodně} z~uniformního rozdělení na
      $[0,1]$.
\item V~každém diskrétním kroku zůstávají všechny uzly aktivní z~předchozího kroku aktivní.
\item Aktivují se uzly splňující
      \[ \sum_{u\in N_v^{\text{aktivní}}} w_{u,v}\ \ge\ \theta_v
         \sum_{u\in N_v} w_{u,v}. \]
\end{itemize}
LTM se soustředí na \emph{příjemce} vlivu.
\end{defbox}

\begin{defbox}[Model nezávislých kaskád (ICM)]
Soustředí se na \emph{vysílající} uzly.
\begin{itemize}
\item Uzel $v$ aktivovaný v~kroku $t$ má \emph{jedinou} příležitost aktivovat každého ze svých
      sousedů, a~to náhodně.
\item Pro sousední uzel $u$ uspěje aktivace s~pravděpodobností $p_{v,u}$ (např.\ $p=0{,}5$).
\item Uspěje-li, $u$ se stane aktivním v~kroku $t+1$.
\item V~dalších kolech se už $v$ o~aktivaci $u$ \emph{nepokouší}.
\item Difuze začíná počáteční aktivovanou množinou uzlů a~pokračuje, dokud není možná žádná
      další aktivace.
\end{itemize}
\end{defbox}

\begin{defbox}[Maximalizace vlivu]
\textbf{Úloha:} pro danou síť a~parametr $k$ vybrat $k$ uzlů do aktivační množiny $B$ tak, aby
se maximalizoval vliv ve smyslu počtu aktivních uzlů na konci. Označme $\sigma(B)$ očekávaný
počet uzlů ovlivněných množinou $B$:
\[ \max_{|B|\le k}\ \sigma(B). \]
Úloha je \textbf{NP-těžká}, ale lze dokázat, že \emph{hladový} přístup dává řešení na úrovni
\textbf{63~\%} optima.

\textbf{Hladový přístup:} začni s~$B=\emptyset$; vyhodnoť $\sigma(\{v\})$ pro každý uzel
a~vyber uzel s~maximálním $\sigma$ jako první prvek $B$; poté opakovaně vybírej uzel, který
$\sigma(B)$ zvýší nejvíce, tj.\ $\arg\max_{v\in V\setminus B}\sigma(B\cup\{v\})$.
\end{defbox}

\subsection{Vliv versus korelace}

Atributy a~chování uživatelů často korelují se sítí. Pro binární atribut u~každého uzlu (zda je
kuřák): koreluje-li se sítí, budou aktéři sdílet stejné hodnoty pozitivně korelovaně se
sociálními spoji (kuřáci interagují spíše s~kuřáky).

\begin{defbox}[Test korelace]
Jestliže podíl hran spojujících uzly s~\emph{různými} hodnotami atributu je významně
\emph{nižší} než očekávaná pravděpodobnost, jde o~důkaz korelace.

\emph{Příklad:} jsou-li spoje nezávislé na kuřáctví, je podíl kuřáků $p=4/9$ a~nekuřáků
$1-p=5/9$. Jedna hrana spojuje dva kuřáky s~pravděpodobností $p\cdot p$, dva nekuřáky
$(1-p)(1-p)$ a~kuřáka s~nekuřákem $2p(1-p)=2\cdot\frac49\cdot\frac59=49\,\%$. Ve
skutečném grafu je tento podíl $2/14=14\,\%<49\,\%$, takže síť \textbf{vykazuje určitý stupeň
korelace} vzhledem ke kuřáctví.
\end{defbox}

\begin{defbox}[Tři sociální procesy vysvětlující korelaci]
\begin{itemize}
\item \textbf{Homofilie} --- tendence propojovat se s~těmi, kdo s~námi sdílejí určitou
      podobnost (\uv{vrána k~vráně sedá}).
\item \textbf{Zmatení} (\emph{confounding}) --- korelace může být vytvořena i~\emph{vnějšími}
      vlivy prostředí (\uv{jednotlivci žijící ve stejném městě se spíše stanou přáteli než
      náhodní jednotlivci}).
\item \textbf{Vliv} (\emph{influence}) --- proces způsobující korelaci chování mezi sousedními
      aktéry (\uv{přejde-li většina mých přátel k~jinému mobilnímu operátorovi, mohu být jimi
      ovlivněn a~přejít také}).
\end{itemize}
\textbf{Proč vliv rozlišovat od homofilie a~zmatení?} \emph{Analýza:} predikce dynamiky
systému --- zda se nová norma chování, technologie nebo idea může šířit jako epidemie.
\emph{Návrh:} konstrukce systému vyvolávajícího konkrétní chování, např.\ vakcinační strategie
(náhodné, cílené na demografickou skupinu, náhodní známí) nebo virální marketingové kampaně.
\end{defbox}

\begin{defbox}[Logistický model vlivu a~shuffle test]
Vliv se určuje \emph{časovými značkami}: v~diskrétních krocích $t=0,1,2,\dots$ se každý neaktivní
uzel stane aktivním s~pravděpodobností $p(a)$, kde $a$ je počet aktivních kontaktů. $p$ se
modeluje \emph{logistickou} funkcí počtu aktivních přátel:
\[ p(a)=\frac{e^{\alpha\ln(a+1)+\beta}}{1+e^{\alpha\ln(a+1)+\beta}},\qquad\text{ekvivalentně}\quad
   \ln\frac{p(a)}{1-p(a)}=\alpha\ln(a+1)+\beta, \]
kde $\alpha$ je \emph{koeficient sociální korelace}, $\beta$ vrozená tendence k~aktivaci
a~$\ln(a+1)$ vysvětlující proměnná.

\textbf{Věrohodnost aktivace.} Stane-li se v~čase $t$ aktivními $y_{a,t}$ uživatelů s~$a$
aktivními přáteli a~$n_{a,t}$ jich zůstane neaktivních, je věrohodnost celé posloupnosti
$\prod_t\prod_a p(a)^{y_{a,t}}\big(1-p(a)\big)^{n_{a,t}}$. Z~logu aktivit se spočítají
$\alpha,\beta$, které ji maximalizují; \textbf{$\alpha$ měří korelaci mezi uzly}.

\textbf{Shuffle test (test na vliv).} Myšlenka: \emph{nehraje-li vliv roli, je časování aktivace
nezávislé na časování ostatních} --- po náhodném zpřeházení časových značek bychom měli dostat
podobnou sociální korelaci. Postup: spočítej $y_{a,t}$, $n_{a,t}$ a~z~nich $\alpha,\beta$;
vyber náhodnou permutaci $\pi$ časů aktivace ($t'_i=t_{\pi(i)}$), z~nových časů spočítej
$y'_{a,t}$, $n'_{a,t}$ a~z~nich $\alpha',\beta'$. Jsou-li $\alpha$ a~$\alpha'$ blízko sebe,
model \emph{nevykazuje} sociální vliv; \emph{významný rozdíl} je důkazem \textbf{vlivu}.
\end{defbox}

\subsection{Homofilie, výběr a~sociální vliv}

\begin{defbox}[Test homofilie]
Motivace: \uv{vrána k~vráně sedá} --- vaši přátelé jsou vám podobnější věkem, rasou, zájmy
a~názory než náhodně vybraná skupina jednotlivců. \emph{Homofilie} je princip, že máme tendenci
být podobní svým přátelům.

\textbf{Test:} kdyby byl každému uzlu \emph{náhodně} přiřazen atribut podle poměru v~reálné
síti, počet \uv{různorodých} hran by se od reálné sítě významně nelišil. Je-li podíl $p$ mužů
a~$q$ žen, je pravděpodobnost hrany muž--muž $p^2$, žena--žena $q^2$ a~\emph{různopohlavní}
hrany $2pq$.

\textbf{Je-li podíl heterogenních (různopohlavních) hran významně nižší než $2pq$, jde
o~důkaz homofilie.}

\emph{Příklad:} různopohlavních hran je 5 z~18; $p=6/9=2/3$, $q=3/9=1/3$; bez homofilie by
jich mělo být $2pq=4/9$, tj.\ 8 z~18 $\Rightarrow$ \emph{důkaz homofilie}.
\end{defbox}

\textbf{Poznámky k~testu homofilie.} \uv{Významně nižší} --- vhodné je posuzovat odchylku pod
střední hodnotou. Co když má síť \emph{významně více} než $2pq$ různorodých hran? Pak jde
o~\emph{inverzní homofilii} (příklad: partnerské vztahy muž--žena). Test lze použít na
libovolnou charakteristiku (rasa, věk, rodný jazyk, preference); pro charakteristiky s~více než
dvěma hodnotami se počítá analogicky --- počet heterogenních hran se porovná s~tím, jak by
vypadal náhodně generovaný graf s~reálnými daty jako pravděpodobnostmi.

\begin{defbox}[Výběr vs.\ sociální vliv]
\textbf{Tranzitivní uzávěr} (\emph{triadic closure}): mají-li dva jednotlivci společného
přítele, je přátelství mezi nimi pravděpodobnější. Homofilie naznačuje, že dva jednotlivci
jsou si podobnější \emph{kvůli} společnému příteli, takže vazba může vzniknout, i~když ani
jeden o~společném příteli neví. Obvykle je obtížné přisoudit vznik vazby jedinému faktoru.

\textbf{Výběr} (\emph{selection}) --- lidé mají tendenci vybírat si přátele, kteří jsou jim
podobní. Působí na různých úrovních intencionality: aktivně si vybíráte podobné přátele
v~malé skupině; \emph{nebo} je populace vaší školy relativně homogenní vůči celkové populaci,
takže vás prostředí \emph{nutí} vybírat si podobné přátele.

\textbf{Mutabilní vs.\ imutabilní charakteristiky:} \emph{imutabilní} se nemění (pohlaví, rasa)
nebo se mění konzistentně s~populací (věk, generace); \emph{mutabilní} se mění v~čase (chování,
přesvědčení, zájmy, názory).

\textbf{Sociální vliv} je \emph{opakem} výběru: \emph{výběr} --- individuální charakteristiky
řídí tvorbu vazeb; \emph{sociální vliv} --- existující vazby formují mutabilní charakteristiky
lidí.
\end{defbox}

\textbf{Longitudinální studie} jsou nutné proto, že z~jediného snímku sítě nelze rozhodnout,
zda působí výběr, nebo vliv: sledují spoje i~chování v~čase a~odpovídají na otázku \emph{mění se
chování po změnách v~síti, nebo síť po změnách v~chování?} Na tom závisí návrh intervencí ---
vykazuje-li užívání drog homofilii \emph{vlivem}, funguje program, v~němž přátelé přesvědčí
ostatní přestat; vznikla-li homofilie \emph{výběrem}, vyberou si bývalí uživatelé nové přátele
a~chování ostatních to neovlivní.

\emph{Christakis a~Fowler} sledovali obezitu a~sociální síť 12\,000 lidí po 32 let a~našli
homofilii podle obezity. Testovali tři vysvětlení --- \emph{výběr} (vybírají si lidé přátele
s~podobným stavem?), \emph{zmatení} (korelují se stavem i~jiné faktory?) a~\emph{sociální vliv}
(ovlivnila změna stavu přátel budoucí stav osoby?) --- a~\textbf{našli významný důkaz pro
třetí}: obezita je typ \uv{nákazy}, který se může šířit sociálním vlivem.

\subsection{Afiliační sítě a~uzávěrové procesy}

\begin{defbox}[Afiliační a~sociálně-afiliační síť]
\textbf{Afiliační síť} vkládá do sítě kontext tím, že zobrazuje spojení k~aktivitám,
společnostem, organizacím, čtvrtím atd. Je to \emph{bipartitní graf}: každá hrana spojuje dva
uzly patřící do \emph{různých} množin (osoby vs.\ \emph{afiliace} čili \emph{fokusy}).

\textbf{Sociálně-afiliační síť} má \emph{dva typy hran}: \emph{osoba--osoba} (přátelství nebo
jiný sociální vztah) a~\emph{osoba--fokus} (účast v~daném fokusu).
\end{defbox}

\begin{defbox}[Tři uzávěrové procesy]
\begin{itemize}
\item \textbf{Tranzitivní uzávěr} (\emph{triadic closure}) --- dva lidé se spojí, protože mají
      společného přítele.
\item \textbf{Fokální uzávěr} (\emph{focal closure}) --- dva lidé se spojí, protože mají
      společný fokus; uzávěr \emph{výběrem}.
\item \textbf{Členský uzávěr} (\emph{membership closure}) --- osoba se zapojí do fokusu,
      protože je v~něm zapojen její přítel; uzávěr \emph{sociálním vlivem}.
\end{itemize}
\end{defbox}

\textbf{Otázky o~uzávěru.} \emph{Jaká je pravděpodobnost, že dva lidé vytvoří vazbu, jako funkce
počtu společných přátel?} Analogicky pro fokální uzávěr (funkce počtu společných fokusů)
a~členský uzávěr (funkce počtu přátel už v~daném fokusu zapojených).

\begin{defbox}[Metodika měření tranzitivního uzávěru]
\begin{enumerate}
\item Pořiď dva snímky sítě v~časech $t_1$ a~$t_2$.
\item Pro každé $k$ identifikuj všechny \emph{páry} uzlů, které mají v~$t_1$ přesně $k$
      společných přátel, ale \emph{nejsou} přímo spojeny hranou.
\item Definuj $T(k)$ jako \emph{podíl} těchto párů, které do $t_2$ vytvořily hranu. To je
      pravděpodobnost, že mezi dvěma lidmi s~$k$ společnými přáteli vznikne vazba.
\item Vynes $T(k)$ jako funkci $k$.
\end{enumerate}

\textbf{Bázový (nulový) model:} je-li pravděpodobnost vytvoření vazby s~jedním společným
přítelem $p$ a~jednotlivé společné přátele považujeme za \emph{nezávislé} příležitosti, pak
\[ T_{\text{base}}(k)=1-(1-p)^{k}\qquad\text{(resp.\ }1-(1-p)^{k-1}\text{ podle konvence).} \]
\end{defbox}

\textbf{Empirické výsledky.} \emph{Tranzitivní uzávěr --- e-mailová síť} (Kossinets, Watts,
2006): komunikace 22\,000 studentů po jeden rok, vazba $=$ e-mail za posledních 60 dní, snímky
den po dni, $T(k)$ zprůměrováno přes páry snímků. Při \emph{žádném} společném příteli se
e-maily téměř nevyměňují; \emph{významný nárůst} z~1 na 2 společné přátele; další nárůst je
významný, ale na menší subpopulaci --- tedy \emph{klesající výnosy}.
\emph{Fokální uzávěr} tamtéž (fokusy $=$ vyučované předměty z~rozvrhů): \emph{sdílení předmětu
má téměř stejný efekt jako sdílení přítele}, opět s~klesajícími výnosy.
\emph{Členský uzávěr}: Backstrom et al.\ (2006) na LiveJournalu (fokusy $=$ členství
v~komunitách) a~Crandall et al.\ (2008) na Wikipedii (vazba $=$ komunikace přes diskusní
stránku, fokusy $=$ editované články) --- tentýž vzorec, nárůst od 1 ke 2 a~pak klesající
výnosy.

\textbf{Výběr a~sociální vliv na datech Wikipedie.} Podobnost dvou editorů se měřila jako
\[ \frac{\text{počet článků editovaných oběma editory}}
   {\text{počet článků editovaných alespoň jedním editorem}} \]
(obdoba překryvu sousedství). Zjištění: \emph{páry editorů, které komunikovaly, jsou si
v~chování významně podobnější než ty, které nekomunikovaly.} Vzniká homofilie tím, že editoři
mluví s~editory téhož článku (\emph{výběr}), nebo tím, že je k~článkům dovedou ti, s~nimiž mluví
(\emph{vliv})? Průměrováním historií interakcí (0 $=$ vznik vazby, jen editoři s~$\ge 100$
akcemi před i~po první interakci, baseline $=$ náhodný pár) vyšlo, že podobnost roste
\emph{už před} vznikem vazby (výběr) a~roste dál \emph{i~po} něm (vliv).

\subsection{Schellingův model segregace}

\begin{defbox}[Schellingův model (T.~Schelling, 70.~léta)]
Jednoduchý model, který ukazuje, jak mohou \emph{globální} vzorce dobrovolné segregace vznikat
z~\emph{lokální} homofilie.
\begin{itemize}
\item Populaci tvoří dva typy jednotlivců --- agenti $X$ a~$O$; oba typy reprezentují
      \emph{imutabilní} charakteristiku (rasa, etnicita, země původu, rodný jazyk).
\item Agenti jsou náhodně rozmístěni v~mřížce.
\item Agent je se svou polohou \emph{spokojen}, je-li alespoň $t$ jeho sousedů agentem téhož
      typu ($t$ je předem nastavený prah).
\item V~každém kole se nejprve identifikují všichni \emph{nespokojení} agenti; poté se všichni
      nespokojení přesunou na místo, kde budou spokojeni.
\end{itemize}
Někdy agent nové uspokojující místo nenajde --- pak zůstane, nebo se přesune na náhodné místo.
Přesuny mohou způsobit, že \emph{dříve spokojený agent přestane být spokojen} --- proto proces
pokračuje.

\textbf{Výsledek: mřížka je segregovanější než na začátku.}

\textbf{Varianty modelu:} pořadí přesunů (náhodné vs.\ \uv{zametací}), cíl přesunu (nejbližší
uspokojující vs.\ náhodné místo), prah (procentní, různý mezi skupinami i~v~nich), \uv{obtočený}
svět. Variant je mnoho, ale \emph{výsledek zůstává stejný: samovolně vzniklá segregace.}

\textbf{Závěr:} Schellingův model je příkladem, jak se \emph{imutabilní} charakteristiky mohou
stát silně korelovanými s~\emph{mutabilními}: volba místa k~životu (mutabilní) se v~čase
přizpůsobí typu agenta (imutabilní), čímž vzniká segregace (homofilie).
\end{defbox}

\subsection{Pozitivní a~negativní vztahy: strukturní balance}

Dosud měl vztah pozitivní konotaci; existují ale sítě s~\emph{negativními} (nepřátelskými)
vztahy. Otázka: \textbf{kdy je síť v~balancovaném stavu?} Opět uvidíme, jak lokální vlastnosti
vynucují globální (srov.\ Schellingův model). Zjednodušení (později uvolníme): \emph{všechny}
páry uzlů jsou spojeny pozitivním či negativním vztahem, tj.\ úplný graf --- vhodné pro malé
skupiny lidí nebo pro státy.

\begin{defbox}[Vlastnost strukturní balance]
Pro tři uzly existují (až na symetrii) čtyři konfigurace znamének:
\begin{center}
\small
\begin{tabular}{@{}p{3.0cm}p{9.5cm}p{2.4cm}@{}}
\toprule
\textbf{Znaménka} & \textbf{Interpretace} & \textbf{Stav} \\
\midrule
$+,+,+$ & $A$, $B$, $C$ jsou vzájemní přátelé & \emph{balancovaný} \\
$+,-,-$ & $A$ a~$B$ jsou přátelé se společným nepřítelem $C$ & \emph{balancovaný} \\
$+,+,-$ & $B$ a~$C$ jsou přátelé $A$, ale spolu nevycházejí & nebalancovaný \\
$-,-,-$ & $A$, $B$, $C$ jsou vzájemní nepřátelé & nebalancovaný \\
\bottomrule
\end{tabular}
\end{center}
\textbf{Vlastnost strukturní balance:} pro každý trojúhelník jsou všechny tři hrany pozitivní,
nebo je pozitivní \emph{právě jedna}.
\end{defbox}

\begin{defbox}[Věta o~balanci]
\textbf{Věta.} Je-li označkovaný \emph{úplný} graf balancovaný, pak buď
\begin{itemize}
\item jsou všechny páry uzlů přátelé, \emph{nebo}
\item lze uzly rozdělit do dvou skupin $X$ a~$Y$ tak, že každý pár uzlů v~$X$ se má rád,
      každý pár v~$Y$ se má rád, a~každý v~$X$ je nepřítelem každého v~$Y$.
\end{itemize}

\textbf{Důkaz.} Nechť $G$ je balancovaná síť. Jsou-li všechny hrany pozitivní, jsme hotovi.
Existuje-li alespoň jedna negativní hrana, vyber uzel $A$ a~polož
$X=$ všichni přátelé $A$ (včetně $A$), $Y=$ všichni nepřátelé $A$. Pak platí:
\emph{(1)}~každé dva uzly v~$X$ jsou přátelé --- kdyby ne, měli bychom trojúhelník
$A$ + dva jeho přátelé s~jedinou negativní hranou, tedy dvě kladné a~jednu zápornou, což je
nebalancované; \emph{(2)}~každé dva uzly v~$Y$ jsou přátelé --- kdyby ne, měli bychom
trojúhelník $A$ + dva jeho nepřátelé se třemi negativními hranami, což je nebalancované;
\emph{(3)}~každý uzel v~$X$ je nepřítelem každého uzlu v~$Y$ --- jinak by trojúhelník
$A$, přítel z~$X$, nepřítel z~$Y$ měl dvě kladné a~jednu zápornou hranu. $\square$
\end{defbox}

\textbf{Aplikace strukturní balance.} \emph{Vývoj aliancí v~Evropě 1872--1907}: Liga tří císařů
1872--81, Trojspolek 1882, přerušení německo-ruské smlouvy 1890, francouzsko-ruská aliance
1891--94, Entente Cordiale 1904, britsko-ruská aliance 1907 --- síť postupně \emph{sklouzává
k~balancovanému označkování}, a~tím do první světové války. \emph{Bangladéš 1972:} podpora
Pákistánu ze strany USA je méně překvapivá, uvědomíme-li si, že SSSR byl nepřítelem Číny, Čína
nepřítelem Indie a~Indie měla špatné vztahy s~Pákistánem --- USA tehdy zlepšovaly vztahy
s~Čínou, a~tak podpořily nepřátele jejích nepřátel.

\begin{defbox}[Slabá strukturní balance a~generalizace]
\textbf{Vlastnost slabé strukturní balance:} neexistuje trojice uzlů, jejíž hrany tvoří
\emph{právě dvě pozitivní a~jednu negativní}. (Konfigurace \uv{tři negativní} je tedy dovolena
--- ta odpovídá situaci, kdy pár vytvoří alianci proti třetímu; \uv{dvě pozitivní a~jedna
negativní} je silnější nebalancovanost, protože $B$ a~$C$ se snaží svou animozitu vyřešit.)

\textbf{Věta.} Je-li označkovaný úplný graf slabě balancovaný, lze jeho uzly rozdělit do
\emph{skupin vzájemných přátel} tak, že každé dva uzly patřící do různých skupin jsou
nepřátelé.

\textbf{Důkaz.} \emph{(1)}~Nechť $A$ je uzel a~$X$ množina přátel $A$ včetně $A$.
\emph{(2)}~Všichni přátelé $A$ jsou přátelé mezi sebou. \emph{(3)}~$A$ a~všichni jeho přátelé
jsou nepřáteli všech ostatních. \emph{(4)}~Odstraň $X$ ze sítě a~pokračuj krokem 1. $\square$

\textbf{Generalizace na neúplné grafy.} \emph{(a)}~Ne všechny páry v~síti jsou spojeny ---
(neúplný) graf je \emph{balancovaný}, lze-li jej \uv{doplnit} přidáním hran na
označkovaný úplný graf, který je balancovaný. \emph{(b)}~Ne všechny trojúhelníky musí být
balancované --- je-li \emph{většina} trojúhelníků balancovaná, lze síť \emph{přibližně}
rozdělit na dvě frakce.
\end{defbox}

\subsection{Analýza vazeb: huby, autority a~algoritmus HITS}

Oba algoritmy analýzy vazeb, PageRank i~HITS, byly publikovány v~roce 1998. PageRank se stal
dominantním díky \emph{nezávislosti na dotazu}, odolnosti proti spamu a~obchodnímu úspěchu
Googlu.

\begin{defbox}[Autority a~huby]
\textbf{Autorita} (\emph{authority}) --- stránka s~mnoha \emph{vstupními} odkazy: má
autoritativní obsah k~tématu, a~proto na ni mnoho lidí odkazuje. \textbf{Hub} --- stránka
s~mnoha \emph{výstupními} odkazy: organizátor informací k~tématu odkazující na mnoho dobrých
autorit.

\textbf{Klíčová myšlenka HITS:} \emph{dobrý hub odkazuje na mnoho dobrých autorit} a~\emph{dobrá
autorita je odkazována mnoha dobrými huby} --- jsou tedy ve vztahu \emph{vzájemného zesilování}
(bipartitní podgraf hustě propojených hubů a~autorit).

\textbf{Proč huby oddělovat od autorit?} Konkurenční firmy na sebe neodkazují, takže se nelze
chápat jako přímá vzájemná podpora. V~PageRanku naopak podpora přechází přímo z~jedné prominentní
stránky na druhou --- to je dominantní způsob podpory u~akademických a~vládních stránek, mezi
blogery a~ve vědecké literatuře.
\end{defbox}

\begin{defbox}[Algoritmus HITS (\emph{Hypertext Induced Topic Search})]
Na rozdíl od PageRanku, který je \emph{statický}, je HITS \emph{závislý na dotazu}. Po zadání
dotazu HITS nejprve rozšíří seznam relevantních stránek vrácených vyhledávačem a~pak vytvoří
\emph{dvě} pořadí rozšířené množiny: pořadí autorit a~pořadí hubů.

\textbf{Sběr stránek.} Pro obecný dotaz $q$: pošli $q$ vyhledávači; posbírej $t$
nejvýše hodnocených stránek (v~původním článku $t=200$) --- tato množina se nazývá
\emph{kořenová množina} $S$; rozšiř $S$ o~každou stránku, na kterou odkazuje nějaká stránka
z~$S$, a~o~každou stránku, která odkazuje na stránku z~$S$ --- vznikne větší
\emph{bázová množina} $B$.

\textbf{Skóre.} Nechť $G=(V,E)$ je hypertextový graf $B$ o~$n$ stránkách s~maticí sousednosti
$L$ ($L_{ij}=1$, existuje-li hrana $(i,j)$, jinak 0). Označme $a(i)$ skóre autority a~$h(i)$
skóre hubu stránky $i$. Vzájemné zesilování:
\[ a(i)=\sum_{(j,i)\in E}h(j),\qquad h(i)=\sum_{(i,j)\in E}a(j), \]
maticově, s~$a=(a(1),\dots,a(n))^{\!\top}$ a~$h=(h(1),\dots,h(n))^{\!\top}$:
\[ a=L^{\!\top}h,\qquad h=La. \]

\textbf{Výpočet mocninnou iterací:}
\begin{algorithmic}[1]
\State $a_0\gets h_0\gets (1,1,\dots,1)^{\!\top}$; $k\gets 1$
\Repeat
  \State $a_k\gets L^{\!\top}h_{k-1}$; \quad $h_k\gets La_k$
  \State $a_k\gets a_k/\|a_k\|_1$; \quad $h_k\gets h_k/\|h_k\|_1$ \Comment{normalizace}
  \State $k\gets k+1$
\Until{$\|a_k-a_{k-1}\|_1<\varepsilon$ \textbf{a} $\|h_k-h_{k-1}\|_1<\varepsilon$}
\State \Return $a$ a~$h$
\end{algorithmic}
\textbf{Normalizace je nezbytná}, protože velikost hodnot autorit a~hubů s~každou aktualizací
roste; konvergence nastane jen s~normalizací (dělí se součtem všech skóre).
\end{defbox}

\begin{defbox}[Spektrální analýza HITS]
Po $k$ krocích platí (spektrum matice $=$ množina jejích vlastních čísel)
\[ a_k=\big(L^{\!\top}L\big)^{k-1}a_0,\qquad h_k=\big(LL^{\!\top}\big)^{k}h_0. \]
Vektor $a_k$ tedy musí konvergovat k~\emph{vlastnímu vektoru} matice $L^{\!\top}L$: existují
normalizační konstanty $c_k$ takové, že $a_k/c_k$ konverguje k~limitě, a~směr se při násobení
maticí nezmění právě tehdy, je-li vektor vlastním vektorem.

Matice $L^{\!\top}L$ je \emph{symetrická}, má tedy $n$ vlastních vektorů $z_1,\dots,z_n$
s~vlastními čísly $c_1,\dots,c_n$, $|c_1|\ge|c_2|\ge\cdots\ge|c_n|$. Rozvineme-li
$a_0=x_1z_1+\cdots+x_nz_n$, pak
\[ a_k=c_1^{k}x_1z_1+\cdots+c_n^{k}x_nz_n,\qquad
   \frac{a_k}{c_1^{k}}=x_1z_1+x_2\Big(\frac{c_2}{c_1}\Big)^{\!k}z_2+\cdots \]
Za předpokladu $c_1>c_2\ge\cdots$ jde pro $k\to\infty$ každý člen kromě prvního k~nule, takže
$a_k/c_1^k$ konverguje k~$x_1z_1$ --- k~\emph{hlavnímu} vlastnímu vektoru. (Konvergenci
nezávisle na počátečním vektoru a~uvolnění předpokladu $c_1>c_2$ zde nedokazujeme.)
\end{defbox}

\begin{defbox}[Vztah ke ko-citaci a~bibliografickému párování]
\textbf{Ko-citace} stránek $i$ a~$j$ (např.\ počet článků, které citují \emph{oba} články $i$
a~$j$):
\[ C_{ij}=\sum_{k=1}^{n}L_{ki}L_{kj}=\big(L^{\!\top}L\big)_{ij}. \]
Matice autorit $L^{\!\top}L$ algoritmu HITS je tedy \emph{ko-citační matice}.

\textbf{Bibliografické párování} stránek $i$ a~$j$ (počet článků, které jsou citovány
\emph{oběma} články $i$ a~$j$):
\[ B_{ij}=\sum_{k=1}^{n}L_{ik}L_{jk}=\big(LL^{\!\top}\big)_{ij}. \]
Matice hubů $LL^{\!\top}$ je tedy matice bibliografického párování.
\end{defbox}

\textbf{Silné a~slabé stránky HITS.} \emph{Silná:} řadí stránky \emph{podle tématu dotazu}, což
dá relevantnější autority a~huby. \emph{Slabé:} snadné spamování (přidat výstupní odkazy na
vlastní stránku je triviální); \emph{drift tématu} (mnoho stránek v~rozšířené množině nemusí být
k~tématu); neefektivnost v~době dotazu (sběr kořenové množiny, rozšíření i~výpočet vlastních
vektorů jsou drahé).

\subsection{PageRank}
\label{ssec:pagerank}

\begin{defbox}[Prestiž pořadí (\emph{rank prestige})]
Míry prestiže uvažují \emph{prominenci aktérů, kteří \uv{hlasují}}: osoba vybraná důležitou
osobou je prestižnější než osoba vybraná méně důležitou (hlas ředitele váží více než hlas
zaměstnance). Je-li váš okruh vlivu plný prestižních aktérů, je vysoká i~vaše prestiž.

Prestiž pořadí $P(i)$ je definována jako lineární kombinace pořadí těch, kteří odkazují na $i$:
\[ P(i)=A_{1i}P(1)+A_{2i}P(2)+\cdots+A_{ni}P(n),\qquad\text{maticově}\quad P=A^{\!\top}P, \]
kde $A_{ji}=1$, odkazuje-li $j$ na $i$, jinak 0. To je \emph{charakteristická rovnice}
vlastního systému matice $A^{\!\top}$; $P$ je vlastní vektor $A^{\!\top}$.
\end{defbox}

\begin{defbox}[Základní PageRank]
PageRank interpretuje hyperodkaz ze stránky $u$ na stránku $v$ jako \emph{hlas} stránky $u$ pro
stránku $v$. Nezajímá jej ale pouhý počet hlasů --- analyzuje i~stránku, která hlas dává:
\emph{hlasy důležitých stránek váží více a~pomáhají činit ostatní stránky důležitějšími}.
Protože stránka může odkazovat na mnoho jiných, musí své skóre \emph{rozdělit}.

Web jako orientovaný graf $G=(V,E)$, $n=|V|$. PageRank stránky $i$:
\[ P(i)=\sum_{(j,i)\in E}\frac{P(j)}{O_j}, \]
kde $O_j$ je počet výstupních odkazů $j$. Maticově: nechť $A_{ij}=1/O_i$ pro $(i,j)\in E$
a~0 jinak; pak $P=A^{\!\top}P$.

\textbf{Intuice:} PageRank je jako \uv{kapalina}, která cirkuluje sítí, přechází od uzlu
k~uzlu po hranách a~shromažďuje se v~nejdůležitějších uzlech. \emph{Celkový PageRank v~síti
zůstává konstantní} --- každá stránka svůj PageRank rozdělí a~pošle po odkazech dál; PageRank
se nikdy nevytváří ani nezaniká, jen se přemísťuje. Proto \emph{není nutná normalizace} ---
na rozdíl od HITS.
\end{defbox}

\textbf{Problém základní definice.} V~mnoha sítích skončí veškerý PageRank ve \uv{špatných}
uzlech. Odkazují-li si $F$ a~$G$ navzájem místo na $A$, PageRank, který k~nim přiteče z~$C$, se
už nikdy nemůže vrátit --- pro velká $k$ dostaneme $\frac12$ pro každý z~$F$, $G$ a~0 pro
všechny ostatní. Jde o~\uv{pomalý únik} (\emph{slow leak}): malé množiny uzlů dosažitelné
z~ostatního grafu, ale bez cest zpět. \textbf{Řešení:} posílat malý zlomek PageRanku
\emph{všem} uzlům.

\begin{defbox}[Škálovaný PageRank a~tlumicí faktor]
Zvol \emph{škálovací faktor} $d$ mezi 0 a~1 a~nahraď základní pravidlo aktualizace:
\begin{enumerate}
\item nejprve aplikuj základní pravidlo;
\item všechny hodnoty PageRanku sniž faktorem $d$ (celkový PageRank v~síti se tím zmenší
      z~1 na $d$);
\item zbývajících $1-d$ jednotek rozděl \emph{rovnoměrně} mezi všechny uzly, tj.\ $(1-d)/n$
      každému.
\end{enumerate}
\textbf{Proč to funguje:} je to \uv{vodní cyklus}, který v~každém kroku $1-d$ jednotek
PageRanku \uv{odpaří} a~rovnoměrně je \uv{sprchne} na všechny uzly.

Opakovaná aplikace škálovaného pravidla \emph{konverguje} k~limitním hodnotám PageRanku
a~ekvilibrium je \emph{jednoznačné} --- hodnoty ale závisejí na volbě $d$. V~praxi se používá
$d$ mezi 0,8 a~0,9 ($d=0{,}85$ v~původním článku); $d$ se nazývá \emph{tlumicí faktor}
(\emph{damping factor}).
\end{defbox}

\begin{defbox}[PageRank jako markovský řetězec]
Web jako \emph{stochastický proces}: stránka $=$ \emph{stav}, hyperodkaz $=$ \emph{přechod}.
\emph{Náhodné surfování}: přechodová pravděpodobnost je $1/O_i$, klikáli surfař na odkazy
uniformně náhodně (bez tlačítka \uv{zpět} a~bez zadávání URL).

Nechť $A$ je matice přechodových pravděpodobností a~$p_0$ počáteční rozdělení; pak
$p_k=A^{\!\top}p_{k-1}$. \emph{Věta o~markovských řetězcích:} konečný řetězec se stochastickou
maticí $A$ má \textbf{jednoznačné stacionární rozdělení}, je-li $A$ \emph{ireducibilní}
a~\emph{aperiodická} --- $p_k$ pak konverguje k~$\pi$ nezávisle na volbě $p_0$ a~v~ustáleném
stavu $\pi=A^{\!\top}\pi$, tedy $\pi$ je hlavní vlastní vektor $A^{\!\top}$ s~vlastním číslem 1.
Právě $\pi$ je vektor PageRanku: odráží \emph{dlouhodobou pravděpodobnost}, že náhodný surfař
stránku navštíví.

\textbf{Pro webový graf neplatí ani jedna podmínka:}
\begin{itemize}
\item \emph{$A$ není stochastická} --- stránky bez výstupních odkazů dají v~matici nulové řádky
      (\emph{visící stránky}, \emph{dangling pages}). Řešení: odstranit je během výpočtu (pořadí
      jiných stránek přímo neovlivňují), nebo z~nich přidat odkazy na všechny stránky
      ($A_{ij}=1/n$).
\item \emph{$A$ není ireducibilní} --- to by znamenalo \emph{silně souvislý} graf (cesta z~$i$
      do $j$ pro každý pár), což web není.
\item \emph{$A$ není aperiodická} --- stav $i$ je \emph{periodický} s~periodou $k>1$, mají-li
      všechny cesty z~$i$ zpět do $i$ délku dělitelnou $k$; řetězec je aperiodický, jsou-li
      aperiodické všechny stavy. Příklad $k=3$: z~1 vede zpět jen cesta $1\to2\to3\to1$.
\end{itemize}
\textbf{Poslední dvě se řeší jedinou strategií:} přidat odkaz z~každé stránky na každou s~malou
přechodovou pravděpodobností řízenou parametrem $d$ --- rozšířená matice je pak ireducibilní
i~aperiodická.
\end{defbox}

\begin{defbox}[Finální PageRank]
Po augmentaci má náhodný surfař na stránce dvě možnosti: s~pravděpodobností $d$ zvolí náhodně
odkaz, který bude následovat, a~s~pravděpodobností $1-d$ skočí na náhodnou stránku:
\[ P=\Big(\frac{1-d}{n}\Big)\,\mathbf{E}+d\,A^{\!\top}P, \]
kde $\mathbf E=\mathbf{ee}^{\!\top}$ ($\mathbf e$ je sloupcový vektor jedniček), tj.\
$\mathbf E$ je matice $n\times n$ samých jedniček. Matice
$\big(\frac{1-d}{n}\big)\mathbf E+dA^{\!\top}$ je stochastická (transponovaná), ireducibilní
i~aperiodická. Po přeškálování tak, aby $\mathbf e^{\!\top}P=n$:
\[ P=(1-d)\,\mathbf e+d\,A^{\!\top}P,\qquad\text{tj.\ pro každou stránku } i: \quad
   P(i)=(1-d)+d\sum_{j=1}^{n}A_{ji}P(j), \]
což je ekvivalentní formuli z~původního článku o~PageRanku
\[ P(i)=(1-d)+d\sum_{(j,i)\in E}\frac{P(j)}{O_j}. \]

\textbf{Výpočet mocninnou iterací:}
\begin{algorithmic}[1]
\State $P_0\gets \mathbf e/n$; $k\gets 1$
\Repeat
  \State $P_k\gets (1-d)\mathbf e+dA^{\!\top}P_{k-1}$
  \State $k\gets k+1$
\Until{$\|P_k-P_{k-1}\|_1<\varepsilon$}
\State \Return $P_k$
\end{algorithmic}
\end{defbox}

\textbf{Výhody PageRanku.} \emph{Boj se spamem:} stránka je důležitá, jsou-li důležité stránky,
které na ni odkazují --- a~protože pro majitele stránky není snadné přidat \emph{vstupní}
odkazy z~jiných důležitých stránek, není snadné PageRank ovlivnit. \emph{Globální
a~dotazově nezávislá míra:} hodnoty PageRanku všech stránek se počítají a~ukládají
\emph{offline}, nikoli v~době dotazu. \textbf{Kritika:} právě dotazová nezávislost --- PageRank
neumí rozlišit stránky autoritativní \emph{obecně} od stránek autoritativních \emph{k~tématu
dotazu}.

\textbf{Důležité:} hyperodkazové řazení není jediný algoritmus ve vyhledávačích --- kombinuje se
s~mnoha faktory založenými na obsahu, aby vzniklo finální pořadí. (Yahoo! a~MSN mají vlastní
nepublikované odkazové algoritmy.)

\subsection{Tematické komunity a~jejich objevování}

Odkazy lze použít i~k~nalezení \emph{tematických komunit} --- skupin tvůrců obsahu či lidí se
společnými zájmy (webové, e-mailové, komunity pojmenovaných entit). Souvisí \emph{cílené
procházení} (\emph{focused crawling}): následují se odkazy a~klasifikátor rozhoduje, zda je
stránka k~tématu.

\begin{defbox}[Tematická komunita]
Pro danou konečnou množinu entit $E=\{e_1,\dots,e_n\}$ \emph{téhož typu} je \emph{tematická
komunita} pár $K=(T,M)$, kde $T$ je \emph{téma} komunity a~$M\subseteq E$ množina všech entit
z~$E$, které téma $T$ sdílejí. Je-li $e\in M$, říkáme, že $e$ je \emph{členem} komunity $K$.
\begin{itemize}
\item \emph{Téma definuje komunitu} --- pro dané $T$ je množina členů jednoznačná, takže dvě
      komunity se stejným tématem jsou si rovny.
\item \emph{Téma lze definovat libovolně} --- událost (skandál) nebo koncept (dolování z~webu).
\item \emph{Entita může být v~libovolném počtu komunit} $\Rightarrow$ komunity se mohou
      \emph{překrývat} a~sdílet členy.
\item \emph{Entity v~$E$ jsou téhož typu} --- lidé a~organizace nemohou být v~téže komunitě.
\end{itemize}
Komunity mohou tvořit \emph{hierarchii}: $(T,M)$ může obsahovat \emph{subkomunity}
$\{(T_1,M_1),\dots,(T_k,M_k)\}$, kde $T_i$ je podtéma $T$ a~$M_i\subseteq M$; $(T,M)$ je pak
\emph{superkomunita} a~každou subkomunitu lze dále dekomponovat.

\textbf{Cíl objevování:} v~datové sadě najít skryté komunity --- pro každou její \emph{téma}
(obvykle množina klíčových slov) a~\emph{členy}.
\end{defbox}

\textbf{Manifestace komunity v~datech:} členové jsou nějak propojeni a~přidružené texty obsahují
slova indikativní pro téma. \emph{Webové stránky} --- propojení \emph{hyperodkazy}, téma ze slov
na stránkách; \emph{e-maily} --- členové komunikují mezi sebou, téma ze slov v~e-mailech;
\emph{textové dokumenty} --- členové se objevují \emph{spolu} v~téže větě (či dokumentu), téma
ze slov ve větě.

\begin{defbox}[Bipartitní jádra komunit]
$(i,j)$-\emph{bipartitní jádro} je úplný bipartitní podgraf s~alespoň $i$ \emph{vějířovými}
uzly (\emph{fans}) a~alespoň $j$ \emph{středovými} uzly (\emph{centers}), tj.\ obsahuje všechny
možné hrany mezi $i$ vějířovými a~$j$ středovými vrcholy. Vějíře reprezentují
\emph{specializované huby}, středy odpovídají \emph{autoritám}.

\textbf{Postup hledání bipartitních jader:}
\begin{itemize}
\item \emph{Krok 1 --- prořezání.} Odstraň velmi silně odkazované stránky (vstupní stupeň $>k$,
      např.\ $k=50$). Pak iterativně: odstraň vějíře s~výstupním stupněm $<j$, středy se
      vstupním stupněm $<i$ a~hrany spojené s~odstraněnými uzly.
\item \emph{Krok 2 --- generování $(i,j)$-jader.} Najdi všechna $(1,j)$-jádra (vrcholy
      s~výstupním stupněm $\ge j$); $(2,j)$-jádra zkonstruuj kontrolou každého vějíře, který
      odkazuje také na některý střed z~$(1,j)$-jádra; pokračuj až do $i$.
\end{itemize}
\textbf{Omezení:} algoritmus najde jen \emph{jádrové} stránky komunit, ne všechny členské
stránky, a~nenajde ani \emph{témata} komunit, ani jejich hierarchickou strukturu.
\end{defbox}

\begin{defbox}[Překrývající se komunity pojmenovaných entit]
\textbf{Cíl:} najít komunity entit z~textového korpusu. \textbf{Hlavní idea:} dvě entity jsou
propojeny, \emph{vyskytnou-li se v~téže větě} (jsou-li dva lidé zmíněni v~téže větě, obvykle mezi
nimi existuje vztah); entita může být ve více komunitách a~téma reprezentují nejvýznamnější
klíčová slova. Slibné výsledky u~komunit politiků a~celebrit z~webových dokumentů.

\textbf{Algoritmus:}
\begin{enumerate}
\item \emph{Postav graf vazeb} --- v~každé větě identifikuj pojmenované entity a~má-li věta více
      než jednu, propoj je \emph{po párech}; klíčová slova z~věty připoj k~párům jako textový
      obsah, ostatní věty zahoď.
\item \emph{Najdi všechny trojúhelníky} --- trojúhelník tří svázaných entit je základní stavební
      blok; komunita se rozšiřuje trojúhelníky sdílejícími hranu.
\item \emph{Najdi jádra komunit} --- jádro je skupina těsně svázaných trojúhelníků, tj.\
      \emph{uvolněná} klika: množina těsně propojených členů.
\item \emph{Klastruj okolo jader} --- trojúhelníky a~páry mimo jádra přiřaď k~jádrům podle
      \emph{podobnosti textového obsahu}.
\end{enumerate}
\end{defbox}

\subsection{Detekce komunit}

\begin{defbox}[Komunita a~detekce komunit]
\emph{Komunity} jsou tvořeny jednotlivci tak, že ti \emph{uvnitř} skupiny interagují mezi
sebou \emph{častěji} než s~těmi vně skupiny. V~různých kontextech se jim říká také
\emph{skupina}, \emph{klastr}, \emph{kohezní podgrupa}, \emph{modul}.

\emph{Detekce komunit} je objevování skupin v~síti, kde příslušnost jednotlivců ke skupinám
\emph{není explicitně dána}.

\textbf{Definice komunity je subjektivní:} hustě propojený podgraf může být komunitou, ale
i~každá komponenta souvislosti může být komunitou.

\textbf{Dva typy skupin v~sociálních médiích:} \emph{explicitní} (vzniklé přihlášením
uživatelů) a~\emph{implicitní} (sociální interakce). Proč extrahovat skupiny z~topologie, i~když lze do
skupin vstoupit? Ne všechny weby platformu pro komunity mají; ne všichni jsou ochotni do skupiny
vstoupit; skupiny se \emph{dynamicky mění}. Interakce navíc doplňují ostatní druhy informací,
pomáhají vizualizaci i~navigaci a~jsou základem pro další úlohy.
\end{defbox}

\begin{defbox}[Metody minimálního řezu]
\textbf{Minimální řez} (Ahuja et al., 1993). \emph{Řez} je rozdělení vrcholů do dvou
disjunktních množin; \emph{úloha minimálního řezu} minimalizuje počet hran (u~vážených sítí
součet vah) mezi nimi --- \emph{velikost řezu}. Motivace: většina interakcí je \emph{uvnitř}
skupiny, mezi skupinami jsou vzácné $\Rightarrow$ detekce komunit $\approx$ minimální řez. Pro
$k>2$ se použije \emph{iterativní bisekce} (skupiny se postupně dělí na dvě).

\textbf{Nevýhoda: produkuje nevyvážené skupiny} (jedna množina často jednoprvková).

\textbf{Složitost:} $s$-$t$ řez přes maximální tok $O(n^5)$; Kargerův algoritmus $O(n^4)$, ale
bez záruky nalezení minima. Pro jednotkové váhy a~bez balančních omezení je 2-cestný řez
polynomiální, $O(n^2\log n)$ (Kargerův randomizovaný minimální řez). Obecné váhy \emph{a}
balanční omezení činí úlohu \textbf{NP-těžkou}.
\end{defbox}

\begin{defbox}[Poměrový a~normalizovaný řez]
Minimální řez často vrací nevyvážené rozdělení (jedna množina je singleton) $\Rightarrow$
změníme účelovou funkci tak, aby brala v~úvahu \emph{velikost komunity}. Pro rozdělení
$\pi=\{C_1,\dots,C_k\}$:
\[ \mathrm{RatioCut}(\pi)=\frac{1}{k}\sum_{i=1}^{k}\frac{\mathrm{cut}(C_i,\bar C_i)}{|C_i|},\qquad
   \mathrm{NCut}(\pi)=\frac{1}{k}\sum_{i=1}^{k}\frac{\mathrm{cut}(C_i,\bar C_i)}{\mathrm{vol}(C_i)}, \]
kde $|C_i|$ je počet uzlů v~$C_i$, $\mathrm{vol}(C_i)$ součet stupňů v~$C_i$
a~$\mathrm{cut}(A,B)=\sum_{i\in A,\,j\in B}a_{ij}$. \textbf{Obě míry preferují vyvážené
rozdělení.}
\end{defbox}

\begin{defbox}[Girvanové--Newmanův algoritmus (2002)]
\emph{Divizivní hierarchický} algoritmus (přístup shora dolů): dekomponuje celý výchozí graf
na postupně menší souvislé části, dokud nezůstanou žádné hrany k~odstranění a~každý uzel
nereprezentuje sám komunitu.

\textbf{Idea:} identifikovat hrany spojující vrcholy z~\emph{různých} komunit (tzv.\ mosty)
a~iterativně je z~grafu odstraňovat. \textbf{Kritérium:} \emph{mezilehlost hrany}
(\emph{edge betweenness}) --- podíl nejkratších cest procházejících danou hranou; hrany
s~vysokou mezilehlostí leží na velkém počtu nejkratších cest mezi vrcholy, a~jsou tedy mosty.

\begin{algorithmic}[1]
\State spočítej mezilehlost všech hran v~síti
\State odstraň hranu s~\emph{nejvyšší} mezilehlostí
\State opakuj předchozí kroky, dokud v~grafu zbývají hrany k~odstranění
\end{algorithmic}
\textbf{Vstup:} celý graf. \textbf{Výstup:} hierarchická struktura, reprezentovaná např.\
dendrogramem.

\textbf{Volba počtu komunit:} algoritmus vrátí \emph{množinu} řešení, ale neurčí nejlepší ---
pro každé dělení se spočítá \emph{modularita} a~vybere se rozdělení s~nejvyšší.

\textbf{Nevýhoda:} vhodný pouze pro sítě \emph{střední velikosti} kvůli vysokým výpočetním
nákladům --- $O(m^2n)$ pro graf o~$m$ hranách a~$n$ vrcholech, resp.\ $O(n^3)$ pro řídké sítě.
\end{defbox}

\emph{Příklad výpočtu mezilehlosti hrany:} mezilehlost hrany $e(1,2)$ je 4, protože všechny
nejkratší cesty z~2 do $\{4,5,6,7,8,9\}$ musí procházet buď $e(1,2)$, nebo $e(2,3)$, a~$e(1,2)$
je nejkratší cesta mezi 1 a~2. Po odstranění $e(4,5)$ se mezilehlost $e(4,6)$ zvýší na 20
(nejvyšší); po odstranění $e(4,6)$ má nejvyšší mezilehlost 4 hrana $e(7,9)$.

\begin{defbox}[Modularita $Q$]
Kvantifikuje kvalitu daného rozdělení grafu na komunity. \textbf{Základní idea:} síť má
smysluplnou komunitní strukturu, je-li počet hran \emph{mezi} komunitami menší, než by se
očekávalo při náhodné volbě. Očekávaný počet hran mezi uzly $i$ a~$j$ se stupni $k_i$, $k_j$
je $k_ik_j/2m$.
\[ Q=\frac{1}{2m}\sum_{i,j}\Big(A_{ij}-\frac{k_ik_j}{2m}\Big)\delta(c_i,c_j), \]
kde $A_{ij}$ je prvek matice sousednosti (počet hran mezi $i$ a~$j$), $m$ počet hran, $k_i$
stupeň vrcholu $i$, $c_i$ skupina, do níž $i$ patří, a~$\delta$ Kroneckerovo delta.

\textbf{Hodnoty:} $Q\in[-1,1]$. Pozitivní $Q$ znamená, že v~síti lze komunitní strukturu najít;
velké pozitivní hodnoty naznačují, že rozdělení odráží skutečnou strukturu. Pro smysluplné
komunity by mělo být $Q\ge 0{,}3$ (Clauset et al., 2004).
\end{defbox}

\begin{defbox}[Louvainova metoda (Blondel et al., 2008)]
\emph{Hladová} optimalizační metoda pro \emph{aglomerativní hierarchickou} maximalizaci
modularity, ve dvou fázích.

\textbf{Fáze 1} --- každý uzel je zprvu samostatná komunita; pak se opakovaně přesouvá izolovaný
uzel do sousední komunity, počítá se zisk/ztráta modularity $Q$ a~při zisku se uzel v~nové
komunitě ponechá. Opakuje se, dokud lze modularitu zlepšit. \emph{Výstup:} rozdělení $k$-té
úrovně ($k$ je číslo iterace).

\textbf{Fáze 2} --- z~původní sítě se vytvoří \emph{supergraf}: každý \emph{supervrchol} je
agregací uzlů jedné komunity z~fáze 1 a~dva supervrcholy spojuje hrana, existuje-li alespoň jedna
hrana mezi odpovídajícími komunitami. Obě fáze se opakují až do maximální modularity.

\textbf{Výhody:} dobrý výkon i~na velkých grafech při nízkých nákladech --- $O(n\log n)$.
\textbf{Nevýhoda:} výsledky jsou \emph{citlivé na pořadí} zpracování.

\textbf{Další nevýhody:} \emph{nesouvislé komunity} --- přesun uzlu může původní komunitu učinit
\emph{vnitřně nesouvislou}, a~přesto v~ní ostatní uzly zůstanou (jsou lokálně optimálně
přiřazeny); \emph{uváznutí v~lokálních optimech}; \emph{selhání v~upřesnění} komunit po
počátečním slučování.
\end{defbox}

\begin{defbox}[Leidenův algoritmus (Traag, Waltman, van Eck, 2019)]
Upřesnění Louvainovy metody pro aglomerativní hierarchickou maximalizaci modularity. Tři fáze:
\emph{(1)}~lokální přesuny uzlů (obdobně jako Louvain); \emph{(2)}~\textbf{upřesnění
rozdělení}; \emph{(3)}~agregace sítě (obdobně jako Louvain).

\textbf{Výhody:} efektivnější než Louvainova metoda, ale pro obrovské grafy může vést
k~delším dobám zpracování; paralelní vícejádrové implementace mohou rychlost zvýšit.
\textbf{Nevýhoda:} \emph{tvrdé rozdělení} --- uzly mohou patřit jen do jedné komunity.
\end{defbox}

\textbf{Výzvy detekce komunit.} Přes množství algoritmů zůstává úloha náročná ze dvou důvodů:
počet komunit obvykle není znám a~je třeba jej určit; komunity jsou \emph{heterogenní} co do
velikosti a~hustoty.

\subsubsection{Taxonomie metod detekce komunit}

Metody lze rozdělit do čtyř kategorií (kritéria se mohou podle úlohy lišit):
\emph{uzlově orientované} (každý uzel ve skupině splňuje určité vlastnosti),
\emph{skupinově orientované} (spojení uvnitř skupiny jako celku splňují určité vlastnosti,
bez zoomu na úroveň uzlů), \emph{síťově orientované} (rozdělení celé sítě do několika
disjunktních množin) a~\emph{hierarchicky orientované} (konstrukce hierarchické struktury
komunit).

\begin{defbox}[Uzlově orientovaná detekce komunit]
Uzly mohou být požadovány, aby splňovaly různé vlastnosti používané v~tradiční SNA:

\textbf{Úplná vzájemnost --- kliky.} \emph{Klika} je maximální úplný podgraf tří a~více
navzájem sousedních uzlů. Nalezení maximální kliky je \textbf{NP-těžké}. \emph{Rekurzivní
prořezávání:} v~klice velikosti $k$ má každý uzel stupeň $\ge k-1$, takže uzly se stupněm
$<k-1$ v~maximální klice nebudou --- opakovaně vzorkuj podsíť, najdi v~ní kliku (hladově),
a~je-li velikosti $k$, odstraň všechny uzly se stupněm $\le k-1$; opakuj, dokud není síť dost
malá. V~sociálních médiích se prořeže mnoho uzlů, protože stupně následují mocninné rozdělení.
Kliky se obvykle používají jako \emph{jádro} či \emph{semeno} pro větší komunity.

\textbf{Metoda perkolace klik} (\emph{clique percolation method}, CPM) --- klika je striktní
a~v~praxi nestabilní definice; CPM hledá \emph{překrývající se} komunity. \emph{Vstup:} $k$
a~síť. \emph{Postup:} najdi všechny kliky velikosti $k$; zkonstruuj \emph{graf klik} (dvě kliky
jsou sousední, sdílejí-li $k-1$ uzlů); komponenty souvislosti grafu klik jsou komunity.
\emph{Příklad:} kliky velikosti 3 jsou $\{1,2,3\}$, $\{1,3,4\}$, $\{4,5,6\}$, $\{5,6,7\}$,
$\{5,6,8\}$, $\{5,7,8\}$, $\{6,7,8\}$ $\Rightarrow$ komunity $\{1,2,3,4\}$ a~$\{4,5,6,7,8\}$
(uzel 4 je v~obou).

\textbf{Dosažitelnost --- $k$-klika, $k$-klub, $k$-klan.} Každý uzel ve skupině má být
dosažitelný v~$k$ skocích. \emph{$k$-klika} je maximální podgraf s~největší geodetickou
vzdáleností dvou uzlů $\le k$ (pozor: \emph{uvnitř podgrafu} může mít průměr $\ge k$ --- 2-klika
$\{12,4,10,1,6\}$ má $d(1,6)=3$). \emph{$k$-klub} má průměr $\le k$; \emph{$k$-klan} je
$k$-klika s~průměrem $\le k$. \emph{Příklad:} kliky $\{1,2,3\}$; 2-kliky $\{1,2,3,4,5\}$,
$\{2,3,4,5,6\}$; 2-kluby $\{1,2,3,4\}$, $\{1,2,3,5\}$, $\{2,3,4,5,6\}$.

\textbf{Stupně uzlů --- $k$-jádro, $k$-plex.} Každý uzel má mít daný počet spojení \emph{uvnitř}
skupiny. \emph{$k$-jádro} ($k$-core): každý uzel je spojen s~alespoň $k$ dalšími členy.
\emph{$k$-plex}: ve skupině o~$n_s$ uzlech je každý uzel sousední s~ne méně než $n_s-k$ uzly.
Definice jsou komplementární ($k$-jádro je $(n_s-k)$-plex), ale množina všech $k$-jader pro dané
$k$ \emph{není} identická s~množinou všech $(n_s-k)$-plexů --- velikost $n_s$ se mezi $k$-jádry
může lišit.

\textbf{Vazby uvnitř vs.\ vně --- LS množiny.} \emph{LS množina} (Luccio--Sami): každá vlastní
podmnožina má více vazeb \emph{ve} skupině než \emph{mimo} ni. Vznikla pro návrh obvodů na
čipech, pro síťovou analýzu je \emph{příliš striktní}. Uvolněním je \emph{$k$-komponenta} ---
maximální podgraf, který zůstává souvislý i~po odstranění méně než $k$ vrcholů či hran (vyžaduje
výpočet hranové souvislosti algoritmem minimálního řezu / maximálního toku).

\textbf{Omezení uzlově orientovaných definic:} příliš striktní (použitelné jako jádro komunity);
neškálovatelné, jen pro malé sítě; někdy nekonzistentní s~vlastnostmi rozsáhlých sítí (stupně
uzlů v~bezškálových sítích).
\end{defbox}

\begin{defbox}[Skupinově orientovaná detekce: kvazikliky]
Skupinově orientovaná kritéria vyžadují, aby \emph{celá skupina} splňovala určitou podmínku ---
např.\ aby hustota skupiny byla dostatečně velká ($\ge$ daný prah). Podgraf
$G_S=(V_S,E_S)$ je \emph{$\gamma$-hustá kvaziklika}, je-li
\[ |E_S|\ \ge\ \gamma\cdot\frac{|V_S|\,(|V_S|-1)}{2}. \]
Lze použít strategii podobnou klikám: \emph{lokální prohledávání} --- vzorkuj podgraf a~najdi
maximální $\gamma$-hustou kvazikliku (velikosti $k$); \emph{heuristické prořezávání} ---
odstraň uzly se stupněm menším než průměrný stupeň ($<\gamma(k-1)$).
\end{defbox}

\begin{defbox}[Síťově orientovaná detekce komunit]
Kritéria uvažují spojení v~síti \emph{globálně}; cílem je rozdělit uzly sítě do
\emph{disjunktních} množin. Přístupy: shlukování podle podobnosti vrcholů; modely latentního
prostoru; aproximace blokovým modelem; spektrální shlukování; maximalizace modularity.

\textbf{Shlukování podle podobnosti vrcholů.} Na uzly se aplikuje $k$-means, kde podobnost
vrcholů je dána podobností jejich \emph{sousedství} (spojení se chápou jako příznaky).
\emph{Strukturní ekvivalence} (spojení s~\emph{touž} množinou aktérů) je pro praxi příliš
restriktivní, používají se míry
\[ \mathrm{Jaccard}(v_i,v_j)=\frac{|N_i\cap N_j|}{|N_i\cup N_j|},\qquad
   \cos(v_i,v_j)=\frac{|N_i\cap N_j|}{\sqrt{|N_i|\,|N_j|}}. \]

\textbf{Modely latentního prostoru.} Zobrazí uzly do \emph{nízkorozměrného} prostoru se
zachováním jejich blízkosti a~pak aplikují $k$-means. \emph{Vícerozměrné škálování} (MDS):
z~párových (např.\ geodetických) vzdáleností se sestaví matice blízkosti; souřadnice
$S\in\mathbb R^{n\times l}$ vyjdou jako $S=V\Lambda^{1/2}$, kde $V$ je $l$ hlavních vlastních
vektorů a~$\Lambda$ diagonální matice jejich vlastních čísel (z~dvojitě centrované matice
blízkosti).

\textbf{Blokové modely.} $A\approx S\,\Sigma\,S^{\!\top}$, kde $\Sigma$ je \emph{míchací matice}
a~$S$ \emph{indikátorová matice komunit}; při uvolnění $S$ na reálné hodnoty jsou optimem hlavní
vlastní vektory $A$.

\textbf{Spektrální shlukování.} Poměrový i~normalizovaný řez lze přeformulovat jako minimalizaci
$\mathrm{tr}(S^{\!\top}LS)$ s~\emph{grafovým laplaciánem} $L=D-A$ (poměrový řez), resp.\
\emph{normalizovaným} $L=D^{-1/2}(D-A)D^{-1/2}$; $D$ je diagonální matice stupňů. Po
\emph{spektrální relaxaci} je optimem matice vlastních vektorů s~\emph{nejmenšími} vlastními
čísly (první je nepoužitelný --- říká, že všechny uzly patří do téhož klastru).

\textbf{Maximalizace modularity.} \emph{Síla komunity} je
$\sum_{i,j\in C}\big(A_{ij}-\frac{d_id_j}{2m}\big)$, modularita
$Q=\frac{1}{2m}\sum_{C}\sum_{i,j\in C}\big(A_{ij}-\frac{d_id_j}{2m}\big)$. Analogicky ke
spektrálnímu shlukování se použije \emph{matice modularity} $B=A-\frac{dd^{\!\top}}{2m}$;
optimem jsou její \emph{hlavní} vlastní vektory, na něž se postprocessingem pustí $k$-means.

\textbf{Jednotný pohled: maticová faktorizace.} Modely latentního prostoru, blokové modely,
spektrální shlukování a~maximalizace modularity lze všechny zapsat jako faktorizaci
$X\approx S\,\Sigma\,S^{\!\top}$, kde
\[ X=\begin{cases}
\text{(dvojitě centrovaná) matice blízkosti} & \text{modely latentního prostoru},\\
\text{sociomatice }A & \text{aproximace blokovým modelem},\\
\text{grafový laplacián} & \text{minimalizace řezu},\\
\text{matice modularity} & \text{maximalizace modularity}.
\end{cases} \]
\textbf{Omezení:} všechny vyžadují, aby uživatel zadal \emph{počet komunit} předem.
\end{defbox}

\begin{defbox}[Hierarchicky orientovaná detekce komunit]
\textbf{Cíl:} vybudovat hierarchickou strukturu komunit na základě topologie sítě; umožňuje
analýzu sítě v~různých \emph{rozlišeních}. Dva reprezentativní přístupy:

\textbf{Divizivní hierarchické shlukování} --- např.\ Girvanové--Newman ($O(m^2n)$, resp.\
$O(n^3)$ pro řídké sítě): rozděl uzly do několika množin a~každou dále dělej na menší (na
rozdělení lze aplikovat síťově orientované metody). Konkrétně: rekurzivně odstraňuj
\uv{nejslabší} vazbu --- najdi hranu s~nejmenší silou, odstraň ji, aktualizuj síly hran
a~opakuj, dokud se síť nerozpadne na požadovaný počet komponent souvislosti; každá komponenta
je komunita. Sílu vazby lze měřit \emph{mezilehlostí hrany} --- vyšší mezilehlost značí most.

\textbf{Aglomerativní hierarchické shlukování} --- např.\ Louvain, $O(n\log n)$: inicializuj
každý uzel jako komunitu; postupně slučuj komunity do větších podle konkrétního kritéria (např.\
na základě \emph{přírůstku modularity}). Dalším příkladem je \emph{Walktrap} --- spouští krátké
náhodné procházky a~podle jejich výsledků slučuje komunity zdola nahoru.
\end{defbox}

\subsection{Predikce vazeb}

\begin{defbox}[Úloha predikce vazeb]
\textbf{Cíl:} predikovat \emph{budoucí} vazby mezi páry uzlů v~síti. V~komerčních sociálních
sítích (Facebook, LinkedIn) se uživatelům doporučují potenciální přátelé. Použít lze podobnost
obsahu i~strukturu:
\begin{itemize}
\item \textbf{Míry založené na obsahu} využívají princip \emph{homofilie} --- uzly s~podobným
      obsahem se spíše propojí. Zlepšují predikci vazeb, ale jsou citlivé na konkrétní síť
      (např.\ krátké a~šumivé tweety s~nestandardními akronymy na Twitteru). \emph{Obsahová
      homofilie nutně neimplikuje strukturní propojenost!}
\item \textbf{Strukturní míry} využívají princip \emph{tranzitivního uzávěru} --- dva uzly, které
      sdílejí sousedství, se spíše propojí. Tranzitivní uzávěr je velmi častý napříč doménami,
      míry jsou účinné v~různých typech sítí a~přímo aplikovatelné. Patří k~nim
      \emph{sousedské} míry, \emph{Katzova} míra, míry založené na \emph{náhodných procházkách}.
      \emph{Strukturní propojenost obvykle implikuje i~obsahovou homofilii.}
\end{itemize}
\end{defbox}

\begin{defbox}[Sousedské míry predikce vazeb]
Využívají počet \emph{společných sousedů} mezi párem uzlů $i$ a~$j$ ke kvantifikaci
pravděpodobnosti budoucí vazby. Nechť $S_i$ je množina sousedů uzlu $i$.

\textbf{Míra společných sousedů:} $\mathrm{CN}(i,j)=|S_i\cap S_j|$.
\emph{Slabina:} nezohledňuje \emph{relativní} počet společných sousedů --- dvě velmi populární
osobnosti spojené s~mnoha aktéry mohou mít mnoho společných sousedů \emph{jen náhodou}.

\textbf{Jaccardova míra:} $\mathrm{Jaccard}(i,j)=\dfrac{|S_i\cap S_j|}{|S_i\cup S_j|}$ ---
podobnost sousedství \emph{normalizovaná} vůči stupňům uzlů, které se kvůli mocninnému
rozdělení silně liší; vyšší stupně obou uzlů koeficient sníží. \emph{Slabina:} nepřizpůsobuje se
\emph{prostředním} sousedům --- je-li společný soused velmi populární, vyskytuje se jako
společný soused i~u~mnoha jiných párů, a~je tedy pro predikci \emph{méně důležitý}.

\textbf{Adamicova--Adarova míra} --- vážená verze míry společných sousedů zohledňující jejich
různou důležitost: váha souseda je \emph{klesající} funkcí jeho stupně, např.\ $1/\log|S_k|$:
\[ \mathrm{AA}(i,j)=\sum_{k\in S_i\cap S_j}\frac{1}{\log|S_k|}. \]
\emph{Příklad:} pro čtyři společné sousedy se stupni 4, 2, 2, 4 je
$\frac{1}{\log 4}+\frac{1}{\log 2}+\frac{1}{\log 2}+\frac{1}{\log 4}=\frac{3}{\log 2}$.
\end{defbox}

\begin{defbox}[Katzova míra]
Sousedské míry \emph{nejsou účinné} při malém počtu společných sousedů: sdílejí-li Alice a~Bob
jediného společného souseda a~Alice a~Jim také, sousedské míry mezi těmito případy nerozliší ---
existuje ale \emph{nepřímá} propojenost přes delší cesty. Pak jsou vhodnější míry založené na
\emph{procházkách}, jako Katzova.

Nechť $W_{ij}^{(t)}$ je počet procházek délky $t$ mezi uzly $i$ a~$j$. Pro předem zvolený
parametr $\beta<1$:
\[ \mathrm{Katz}(i,j)=\sum_{t=1}^{\infty}\beta^{t}\,W_{ij}^{(t)}, \]
kde $\beta$ je \emph{diskontní faktor}, který potlačuje význam delších procházek; pro malá
$\beta$ nekonečná suma konverguje. Je-li $A$ symetrická matice sousednosti neorientované sítě,
lze matici Katzových koeficientů $K$ rozměru $n\times n$ spočítat jako
\[ K=\sum_{t=1}^{\infty}\beta^{t}A^{t}=(I-\beta A)^{-1}-I \]
(vlastní čísla $A^t$ jsou $t$-té mocniny vlastních čísel $A$). Silně centrální uzly mají
tendenci tvořit vazby s~mnoha uzly $\Rightarrow$ součet Katzových koeficientů uzlu $i$ vůči
ostatním uzlům se nazývá jeho \textbf{Katzova centralita}.
\end{defbox}

\begin{defbox}[Predikce vazeb jako klasifikační úloha]
Přítomnost či absence vazby mezi párem uzlů se chápe jako \emph{binární} indikátor třídy;
pro každý pár uzlů se zkonstruuje vícerozměrný datový záznam:
\begin{itemize}
\item \emph{příznaky} zahrnují všechny různé podobnosti mezi uzly --- sousedské, Katzovy,
      založené na procházkách --- plus řadu příznaků preferenčního připojování, např.\ stupně
      obou uzlů v~páru;
\item výsledkem je \emph{pozitivně-neoznačkovaná} klasifikační úloha: páry uzlů s~hranou jsou
      pozitivní příklady a~zbývající páry jsou neoznačkované;
\item neoznačkované příklady lze při trénování přibližně považovat za negativní;
\item protože je negativních párů v~rozsáhlých a~řídkých sítích příliš mnoho, použije se jen
      \emph{vzorek} negativních příkladů.
\end{itemize}

\textbf{Algoritmus.} \emph{Trénovací fáze:} vygeneruj data s~jedním záznamem pro každý pár uzlů
s~hranou a~vzorek záznamů z~párů bez hrany; příznaky odpovídají extrahovaným podobnostním
a~strukturním rysům; označení třídy je přítomnost či absence hrany; postav a~natrénuj model
klasifikátoru. \emph{Testovací fáze:} převeď každý testovaný pár uzlů na vícerozměrný záznam
a~použij natrénovaný klasifikátor k~predikci označení. Běžnou volbou bázového klasifikátoru je
\emph{logistická regrese}; kvůli nevyváženosti dat se často používají \emph{cenově citlivé}
verze klasifikátorů.

\textbf{Výhody supervizovaného přístupu:} obsahové příznaky lze snadno přidat; klasifikátor se
\emph{sám} naučí relevanci příznaků; pro orientované sítě lze příznaky extrahovat
\emph{asymetricky} (vstupní a~výstupní stupně, personalizovaný PageRank); vyšší flexibilita
díky schopnosti učit se vztahy mezi vazbami a~příznaky různých typů.
\end{defbox}

\textbf{Souhrn predikce vazeb.} Účinnost se liší podle datové sady: \emph{sousedské} míry se
efektivně spočítají i~pro rozsáhlá data a~fungují téměř tak dobře jako ostatní nesupervizované;
\emph{Katzova} míra a~míry na \emph{náhodných procházkách} jsou pro \emph{velmi řídké} sítě,
kde nelze počet společných sousedů robustně změřit; \emph{supervize} dá lepší přesnost
a~adaptabilitu napříč doménami, ale je drahá; obsah predikci vylepšuje, strukturní míry jsou
však \emph{výrazně silnější} (kvůli tranzitivitě reálných sítí). Obsahové míry stavějí na
\uv{obrácené homofilii} --- podobný či s~vazbami korelovaný obsah se využívá k~predikci vazeb.

\subsection{Shrnutí ke~zkoušce}

\begin{itemize}
\item \textbf{SNA} $=$ sociální vědy $+$ analýza sítí $+$ teorie grafů; soustředí se na
      \emph{vztahy}, ne na entity a~jejich atributy. Reprezentace: obrázek, seznam hran, matice
      sousednosti (u~orientovaných grafů nesymetrická). \emph{Ego síť} vs.\ \uv{celá} síť ---
      \emph{problém specifikace hranic}: žádná studovaná síť není celá.
\item \textbf{Síla vazeb}: silné vs.\ slabé; \emph{síla slabých vazeb} (Granovetter 1973).
      Homofilie, tranzitivita a~přemostění; homofilie $+$ tranzitivita $\to$ \emph{kliky}.
      Most $=$ hrana, jejíž odstranění rozpojí koncové uzly; \emph{zkratkový most} $=$
      vzdálenost po odstranění vzroste nad 2. \textbf{Překryv sousedství}
      $O(i,j)=\frac{|N_i\cap N_j|}{|N_i\cup N_j|-2}$ --- čím větší, tím silnější vazba.
\item \textbf{Čtyři míry centrality}: \emph{stupňová} (kolik lidí přímo dosáhne),
      \emph{mezilehlá} (na kolika nejkratších cestách leží; kde se síť rozpadne),
      \emph{blízkostní} (převrácená průměrná délka nejkratších cest; rychlost šíření),
      \emph{vlastní} ($Ax=\lambda x$, jen \emph{největší} vlastní číslo dává nezáporný vektor;
      jak dobře je spojen s~dobře propojenými). Umět interpretační tabulku a~vysvětlit, proč
      nestačí jedna míra (uzly 3 a~5 \emph{společně} $>$ uzel 10). \textbf{Strukturní
      ekvivalence}: pozice lze vyměnit bez změny struktury.
\item \textbf{Koheze}: reciprocita (jen orientované grafy), hustota
      ($|E|/\binom{|V|}{2}$; klika $=1$; orientované mají poloviční), \textbf{klastrovací
      koeficient} $C_i=\frac{2|\{e_{jk}\}|}{k_i(k_i-1)}$ a~sítě jako průměr, průměr sítě
      (\emph{diameter}) $=$ nejdelší nejkratší cesta, průměrná vzdálenost, excentricita.
      \textbf{Malý svět} $=$ vysoký $C$ $+$ krátká průměrná cesta.
\item \textbf{Šest vlastností reálných komplexních sítí}: efekt malého světa (Milgram, medián
      6), tranzitivita/klastrování, mocninné rozdělení stupňů, odolnost (robustní k~náhodnému,
      zranitelná cíleným odstraněním hubů), vzorce míchání (asortativní míchání $=$ homofilie),
      komunitní struktura. \textbf{Bezškálové sítě}: růst $+$ \emph{preferenční připojování}
      (rich-get-richer, Matoušův efekt, kumulativní výhoda); důvody --- popularita, kvalita,
      smíšený model (halo efekt). \textbf{Centralizace} a~struktura \emph{jádro--periferie}
      (bow-tie analýza).
\item \textbf{Modelování vlivu}: společné rysy (orientovaný graf, aktivní/neaktivní,
      \emph{aktivovaný uzel nelze deaktivovat}). \textbf{LTM} --- prah $\theta_v$ uniformně
      z~$[0,1]$, aktivace při $\sum_{\text{aktivní}}w_{u,v}\ge\theta_v\sum w_{u,v}$; soustředí
      se na \emph{příjemce}. \textbf{ICM} --- \emph{jediná} příležitost aktivovat každého
      souseda s~pravděpodobností $p_{v,u}$; soustředí se na \emph{vysílající}.
      \textbf{Maximalizace vlivu} je NP-těžká, hladový přístup dává \textbf{63~\%} optima.
\item \textbf{Vliv vs.\ korelace}: \emph{test korelace} (podíl hran mezi různými hodnotami
      atributu vs.\ $2p(1-p)$; příklad kuřáků $14\,\%<49\,\%$); tři procesy --- homofilie,
      \emph{zmatení}, vliv; logistický model $\ln\frac{p(a)}{1-p(a)}=\alpha\ln(a+1)+\beta$,
      $\alpha$ $=$ koeficient sociální korelace; \textbf{shuffle test} --- zpřeházej časové
      značky; liší-li se $\alpha'$ významně od $\alpha$, jde o~důkaz \emph{vlivu}.
\item \textbf{Test homofilie}: heterogenních hran je významně méně než $2pq$ $\Rightarrow$
      homofilie; významně \emph{více} $\Rightarrow$ \emph{inverzní} homofilie (partnerské
      vztahy). Umět příklad: 5 z~18 vs.\ $2pq=8$ z~18.
\item \textbf{Výběr vs.\ sociální vliv}: výběr $=$ charakteristiky řídí tvorbu vazeb; vliv $=$
      vazby formují \emph{mutabilní} charakteristiky. \emph{Imutabilní} (pohlaví, rasa, věk)
      vs.\ \emph{mutabilní} (chování, názory). Rozlišit lze jen \emph{longitudinálními}
      studiemi (obezita: Christakis \& Fowler, 12\,000 lidí / 32 let $\to$ důkaz pro
      \emph{sociální vliv}).
\item \textbf{Afiliační sítě}: bipartitní (osoby $\times$ fokusy), sociálně-afiliační má dva
      typy hran. \textbf{Tři uzávěry}: tranzitivní (společný přítel), \emph{fokální} (společný
      fokus; \emph{výběr}), \emph{členský} (přítel už ve fokusu; \emph{sociální vliv}).
      Metodika měření $T(k)$ ze dvou snímků; bázový model $T_{\text{base}}(k)=1-(1-p)^k$;
      empiricky velký nárůst z~1 na 2 a~poté \emph{klesající výnosy}.
\item \textbf{Schellingův model}: agenti dvou typů v~mřížce, spokojeni při $\ge t$ sousedech
      téhož typu, nespokojení se přesouvají $\Rightarrow$ \emph{samovolná segregace}. Ukazuje,
      jak se \emph{imutabilní} charakteristika stane silně korelovanou s~\emph{mutabilní}
      (volbou místa) a~jak lokální homofilie vytvoří globální vzorec.
\item \textbf{Strukturní balance}: balancované trojúhelníky jsou $+{+}{+}$ a~$+{-}{-}$
      (tj.\ všechny tři pozitivní, nebo právě jedna pozitivní). \textbf{Věta o~balanci}: úplný
      balancovaný graf $\Rightarrow$ všichni přátelé, nebo rozdělení na dvě skupiny vnitřních
      přátel a~vzájemných nepřátel --- umět důkaz ($X$ $=$ přátelé $A$, $Y$ $=$ nepřátelé
      $A$, tři kroky). \textbf{Slabá balance}: zakázána jen konfigurace \emph{dvě pozitivní
      a~jedna negativní} $\Rightarrow$ rozdělení na \emph{libovolný počet} skupin. Aplikace:
      evropské aliance 1872--1907, Bangladéš 1972.
\item \textbf{HITS}: dotazově \emph{závislý}; kořenová množina $S$ ($t=200$) $\to$ bázová
      množina $B$; $a=L^{\!\top}h$, $h=La$, mocninná iterace s~\textbf{nutnou normalizací};
      $a_k=(L^{\!\top}L)^{k-1}a_0$ konverguje k~hlavnímu vlastnímu vektoru $L^{\!\top}L$
      $=$ \emph{ko-citační} matice ($LL^{\!\top}$ $=$ bibliografické párování). Slabiny: snadné
      spamování, drift tématu, neefektivnost v~době dotazu.
\item \textbf{PageRank}: dotazově \emph{nezávislý}, $P(i)=\sum_{(j,i)\in E}P(j)/O_j$,
      \emph{nepotřebuje normalizaci} (celkový PageRank je konstantní --- \uv{kapalina}).
      Problém \uv{pomalého úniku} $\Rightarrow$ \emph{škálované} pravidlo s~tlumicím faktorem
      $d$ (0,8--0,9; $0{,}85$ v~článku), \uv{vodní cyklus}. Markovský řetězec potřebuje
      \emph{stochastickou, ireducibilní a~aperiodickou} matici --- webový graf nesplňuje ani
      jednu (visící stránky, není silně souvislý, periodicita); vše se řeší \emph{jedinou}
      strategií (odkaz z~každé stránky na každou s~malou pravděpodobností). Finální tvar
      $P(i)=(1-d)+d\sum_{(j,i)\in E}P(j)/O_j$, výpočet mocninnou iterací. Výhody: odolnost
      proti spamu (vstupní odkazy nelze snadno získat), globální offline míra; kritika:
      dotazová nezávislost.
\item \textbf{Tematické komunity}: $K=(T,M)$, téma určuje komunitu, komunity se mohou
      \emph{překrývat}, hierarchie subkomunit. \emph{Bipartitní jádra} --- $(i,j)$-jádro,
      vějíře $=$ huby, středy $=$ autority; prořezání $+$ generování jader; najde jen jádra,
      ne témata ani hierarchii. \emph{Komunity pojmenovaných entit}: graf ze společných výskytů
      ve větě $\to$ trojúhelníky $\to$ jádra $\to$ klastrování podle textové podobnosti.
\item \textbf{Detekce komunit --- definice je subjektivní}; explicitní vs.\ implicitní skupiny.
      \emph{Minimální řez}: nevyvážené skupiny, $O(n^5)$ / Karger $O(n^4)$, s~obecnými vahami
      a~balančními omezeními \textbf{NP-těžké} $\Rightarrow$ \emph{poměrový} a~\emph{normalizovaný
      řez} (dělení $|C_i|$, resp.\ $\mathrm{vol}(C_i)$).
\item \textbf{Girvanové--Newman}: divizivní, odstraňuje hrany s~nejvyšší \emph{mezilehlostí},
      $O(m^2n)$ / $O(n^3)$; počet komunit se vybere podle nejvyšší \textbf{modularity}
      $Q=\frac{1}{2m}\sum_{i,j}(A_{ij}-\frac{k_ik_j}{2m})\delta(c_i,c_j)$, $Q\in[-1,1]$,
      smysluplné komunity $Q\ge 0{,}3$.
\item \textbf{Louvain}: hladová aglomerativní maximalizace modularity ve dvou fázích (lokální
      přesuny $\to$ supergraf), $O(n\log n)$; nevýhody: citlivost na pořadí, \emph{nesouvislé
      komunity}, lokální optima. \textbf{Leiden} přidává fázi \emph{upřesnění rozdělení}; jen
      tvrdé rozdělení. Výzvy: neznámý počet komunit, heterogenní velikosti a~hustoty.
\item \textbf{Čtyři kategorie metod}: \emph{uzlově orientované} (kliky --- NP-těžké, rekurzivní
      prořezávání uzlů se stupněm $<k-1$, \textbf{CPM} pro \emph{překrývající se} komunity
      přes graf klik sdílejících $k-1$ uzlů; $k$-klika / $k$-klub / $k$-klan; $k$-jádro
      a~$k$-plex jsou komplementární; LS množiny a~$k$-komponenty), \emph{skupinově
      orientované} ($\gamma$-husté kvazikliky), \emph{síťově orientované} (podobnost vrcholů ---
      Jaccard, kosinus $+$ $k$-means; MDS; blokové modely; spektrální shlukování přes grafový
      laplacián s~\emph{nejmenšími} vlastními čísly; maximalizace modularity přes matici
      modularity s~\emph{největšími} --- vše sjednoceno jako maticová faktorizace
      $X\approx S\Sigma S^{\!\top}$, vyžaduje zadat počet komunit), \emph{hierarchicky
      orientované} (divizivní $=$ Girvanové--Newman, aglomerativní $=$ Louvain, Walktrap).
\item \textbf{Predikce vazeb}: obsahové míry (homofilie) vs.\ strukturní (tranzitivní uzávěr;
      silnější). Sousedské míry: \emph{společní sousedé} $|S_i\cap S_j|$ (nezohledňuje stupně),
      \emph{Jaccard} $\frac{|S_i\cap S_j|}{|S_i\cup S_j|}$ (normalizuje stupně obou uzlů, ne
      prostředních sousedů), \emph{Adamic--Adar} $\sum_{k}\frac{1}{\log|S_k|}$ (váha klesá se
      stupněm společného souseda). \textbf{Katzova míra}
      $\sum_t\beta^tW_{ij}^{(t)}=(I-\beta A)^{-1}-I$ pro \emph{řídké} sítě s~málo společnými
      sousedy; součet přes ostatní uzly $=$ \emph{Katzova centralita}. \textbf{Klasifikační
      pojetí}: pozitivně-neoznačkovaná úloha, vzorek negativ, logistická regrese jako bázový
      klasifikátor, cenově citlivé varianty; výhoda --- klasifikátor se sám naučí relevanci
      příznaků a~lze použít \emph{asymetrické} příznaky u~orientovaných sítí.
\end{itemize}

%=====================================================================
\newpage
\section{Praktické využití technik}
\label{sec:praxe}
%=====================================================================

Poslední věta okruhu žádá spojení obou polovin dohromady: k~čemu se probrané techniky
\emph{skutečně} používají. U~zkoušky se hodí umět ke každé aplikaci říct \emph{jakou úlohu}
řeší (kap.~\ref{sec:paradigmata}), \emph{jakou metodou} (kap.~\ref{sec:metody}) a~\emph{jak se
vyhodnocuje} (kap.~\ref{sec:reprez}).

\subsection{Dobývání znalostí}

\begin{center}
\small
\begin{tabular}{@{}p{3.5cm}p{4.4cm}p{6.9cm}@{}}
\toprule
\textbf{Aplikace} & \textbf{Typ úlohy / metoda} & \textbf{Poznámka} \\
\midrule
Analýza nákupního košíku & asociační pravidla, Apriori / FP-Growth &
  plánování a~rozvržení prodejny, kupony a~časově omezené slevy, \uv{balíčkování} produktů;
  srovnání nových a~existujících prodejen pomocí \emph{virtuálních položek} \\
Detekce podvodů & klasifikace, detekce odlehlých hodnot &
  20 mil.\ výkonů zdravotní pojišťovny $\to$ 214 tis.\ podivných $\to$ 14 tis.\ předaných
  k~ruční kontrole $\to$ 8 tis.\ potvrzených podvodů; stroj \emph{zaostřuje} pozornost
  experta \\
Kreditní skórování & rozhodovací stromy, naivní Bayes, logistická regrese &
  průběžný příklad v~celém textu; požaduje se \emph{srozumitelnost} (regulace) a~zacházení
  s~chybějícími hodnotami \\
Segmentace zákazníků & klastrová analýza ($k$-means, hierarchické) &
  reprezentace segmentu \emph{centroidem} (\uv{typický zákazník}) je pro byznys velmi
  intuitivní \\
Predikce odchodu (churn) & klasifikace na nevyvážených datech &
  accuracy nestačí $\Rightarrow$ precision/recall, $F_1$, ROC/AUC; \emph{lift analýza} pro
  rozhodnutí, kolik zákazníků kontaktovat \\
Cílené marketingové kampaně & skórování a~řazení, lift &
  ekonomická úvaha nad lift křivkou (příklad s~hodinkami: 1.~přihrádka $+\$550$,
  10.~přihrádka $-\$475$) \\
Klasifikace textu & SVM s~kernely &
  patří k~nejlepším klasifikátorům pro text (vysokorozměrná data); srozumitelnost se
  neočekává \\
Predikce časových řad & regrese, neuronové sítě &
  řada se transformuje posuvným oknem na tabulku vstupů a~výstupů \\
Klasifikace chemických látek & klasifikace strukturních dat &
  reprezentace řetězci SMILES nebo fakty v~Prologu \\
Web usage mining & sekvenční vzory (GSP) &
  navigační vzory uživatelů podle posloupností návštěv stránek \\
Doporučovací systémy \mimo{} & klastrová analýza, klasifikace, predikce vazeb &
  kolaborativní filtrování (podobnost uživatelů či položek) vs.\ obsahové filtrování;
  studený start řeší obsah \\
\bottomrule
\end{tabular}
\end{center}

\textbf{Poznámka k~motivaci rozhodovacích stromů.} Banka Barclays zavedla telefonní systém
rozdělující volající podle úvěrového hodnocení: klienti s~vysokými úsporami nebo dostatečným
limitem šli na britského operátora, ostatní do indického call centra --- cílem bylo
\uv{odfiltrovat} zákazníky bez peněz na produkty banky. Kritika: systém preferuje bohatší
klienty a~nemá za cíl dát nejlepší radu. Ilustruje, že \emph{technicky správný} model může mít
problematické důsledky a~že fáze \emph{vyhodnocení} v~CRISP-DM musí zahrnovat i~pohled mimo
přesnost.

\subsection{Analýza sociálních sítí}

\begin{center}
\small
\begin{tabular}{@{}p{3.7cm}p{11.2cm}@{}}
\toprule
\textbf{Oblast} & \textbf{Využití} \\
\midrule
Podniky & analýza a~zlepšení komunikačních toků v~organizaci, s~partnery či zákazníky;
  marketingové kampaně založené na sociálních sítích (vizualizace komunikačních toků před
  a~po zavedení systému pro správu obsahu) \\
Bezpečnostní složky a~armáda & identifikace kriminálních a~teroristických sítí ze zachycené
  komunikace; identifikace \emph{klíčových hráčů} v~nich (míry centrality,
  odd.~\ref{sec:sna}) \\
Sociální sítě (Facebook) & identifikace a~doporučení potenciálních přátel na základě
  \uv{přátel přátel} --- přímo úloha \emph{predikce vazeb} \\
Občanská společnost & odhalování konfliktů zájmů ve skrytých vazbách mezi státními orgány,
  lobbisty a~podniky \\
Síťoví operátoři & optimalizace struktury a~kapacity komunikačních sítí \\
Medicína & nakažlivé nemoci a~jejich prevence; sexuálně přenosné nemoci --- epidemie
  syfilis v~Rockdale County (Atlanta, 1996): skupina bílých dívek okolo 16 let, skupina
  bílých chlapců 17--21 let a~skupina afroamerických chlapců 17--21 let; syfilis
  diagnostikována u~6 bílých žen, 2 bílých mužů a~2 afroamerických mužů \\
Bezškálové sítě & \emph{výpočetní technika:} SF sítě (WWW) jsou odolné proti náhodným poruchám,
  ale zranitelné koordinovaným útokem --- vymýcení internetových virů je v~podstatě nemožné;
  \emph{medicína:} vakcinace cílená na huby (lidé s~mnoha kontakty, obtížně identifikovatelní),
  léky mířící na klíčové molekuly nemoci podle vazebných sítí v~buňkách; \emph{business:} vazby
  mezi firmami a~odvětvími pro eliminaci kaskádových krachů, šíření \uv{nákazy} v~marketingu \\
Detekce komunit & \emph{Belgický mobilní operátor} (Louvainova metoda): cca 2 miliony
  zákazníků; velikost uzlu odpovídá počtu jednotlivců v~komunitě, barva hlavnímu jazyku
  komunity (červená francouzština, zelená nizozemština); zobrazeny jen komunity nad
  100 zákazníků. Zoom ve vyšším rozlišení odhalí \emph{přechodovou} komunitu tvořenou
  subkomunitami s~méně zřetelným jazykovým oddělením \\
\bottomrule
\end{tabular}
\end{center}

\textbf{Kdy a~proč SNA použít.} Studovat sociální síť nebo zvýšit její efektivitu; vizualizovat
vzorce ve vztazích a~interakcích; sledovat cesty, kterými proudí informace; identifikovat typy
aktérů (klíčové hráče); najít příčiny dysfunkčních komunit a~podpořit kohezi a~růst online
komunity. Klíčová myšlenka: rozsah akcí a~příležitostí dostupných jednotlivci je často funkcí
jeho \emph{pozice} v~síti, takže odhalení pozic (otec, učitel, pracovník) přináší i~překvapivé
výsledky.

\begin{poznbox}[Úvahy o~návrhu (\emph{thoughts on design})]
Teorie SNA inspirovala některé z~prvních sociálních sítí (např.\ SixDegrees), stále se však
při návrhových rozhodnutích nepoužívá tak často. Otevřené otázky:
\begin{itemize}
\item Jak může online platforma (nebo její správci) využít metody a~poznatky SNA?
\item Jak může platforma povzbudit uživatele, aby jednali (\emph{lokálně}) tak, že tím
      rozšíří komunitu?
\item Jaká opatření může online komunita přijmout k~optimalizaci své struktury? Příklad:
      \emph{kliky} mohou být nežádoucí, protože odhánějí nově příchozí.
\item Jaké struktury jsou žádoucí pro různé typy online platforem?
\item Jak mohou online komunity identifikovat a~využít klíčové hráče ke svému prospěchu?
\end{itemize}
\end{poznbox}

\subsection{Provozní a~etické aspekty nasazení}

\mimo{} Nasazení modelu (fáze \emph{Deployment} v~CRISP-DM, odd.~\ref{ssec:metodologie}) není
konec projektu:
\begin{itemize}
\item \textbf{Monitoring a~drift.} Rozdělení dat se v~čase mění (\emph{data drift}), stejně
      jako vztah mezi vstupy a~cílem (\emph{concept drift}) --- model je nutné pravidelně
      přeučovat a~sledovat pokles kvality na produkčních datech.
\item \textbf{Zpětná vazba modelu do dat.} Nasazený model ovlivňuje realitu, kterou měří:
      skórovací model pro schvalování úvěrů vidí jen ty žadatele, kterým úvěr schválil, takže
      další trénovací data jsou jím samým zkreslená.
\item \textbf{Soukromí.} Analýza sociálních sítí pracuje s~daty \emph{o~vztazích}, které
      uživatel často nemohl ovlivnit; i~anonymizovaný graf lze deanonymizovat podle struktury.
      Predikce vazeb a~modelování vlivu jsou v~podstatě nástroje pro předpovídání a~ovlivňování
      chování konkrétních lidí.
\item \textbf{Spravedlnost a~diskriminace.} Model naučený na historických datech reprodukuje
      historickou nerovnost (viz případ Barclays výše); chráněný atribut lze často
      \emph{rekonstruovat} z~ostatních atributů, takže jeho vypuštění samo nestačí.
\item \textbf{Právní rámec.} GDPR čl.~22 (právo na vysvětlení automatizovaného rozhodnutí)
      a~AI Act pro vysoce rizikové systémy --- v~regulovaných doménách je interpretovatelnost
      právním požadavkem, nikoli pohodlím.
\end{itemize}

\subsection{Shrnutí ke~zkoušce}

\begin{itemize}
\item Ke každé aplikaci umět říct \textbf{trojici}: typ úlohy (deskriptivní / prediktivní,
      s~učitelem / bez učitele), použitá metoda, způsob vyhodnocení.
\item \textbf{Vlajkové aplikace dobývání znalostí}: analýza nákupního košíku (asociační
      pravidla), detekce podvodů (příklad zdravotní pojišťovny --- stroj \emph{zaostřuje}
      pozornost experta, nenahrazuje jej), kreditní skórování (stromy a~Bayes; požadavek
      srozumitelnosti), segmentace zákazníků (klastrování, centroid $=$ typický zákazník),
      churn a~cílený marketing (nevyvážená data $\to$ ROC/AUC a~\emph{lift}), klasifikace textu
      (SVM), predikce časových řad (posuvné okno), web usage mining (sekvenční vzory).
\item \textbf{Vlajkové aplikace SNA}: identifikace klíčových hráčů v~kriminálních
      a~teroristických sítích (centralita), doporučování přátel (predikce vazeb), optimalizace
      komunikačních sítí, epidemiologie (Rockdale County 1996), vakcinace cílená na \emph{huby},
      marketing a~kaskádové finanční krachy v~bezškálových sítích, segmentace zákazníků
      operátora podle jazykových komunit (Louvain, 2~mil.\ zákazníků).
\item \textbf{Klíčová vlastnost SF sítí pro praxi}: odolnost proti náhodným poruchám (až 80~\%
      uzlů) vs.\ zranitelnost cíleným útokem (5--15~\% hubů) --- z~toho plyne jak strategie
      vakcinace, tak nemožnost vymýtit internetové viry.
\item \textbf{Nasazení není konec}: monitoring driftu, zpětná vazba modelu do dat, soukromí
      (i~anonymizovaný graf lze deanonymizovat), spravedlnost (chráněný atribut lze
      rekonstruovat z~ostatních), GDPR čl.~22 a~AI Act.
\end{itemize}

\end{document}
